% SOURCES. Binding final4/GOAL_alt.md: geometry is hydrodynamic and obeys
% continuum Gauss laws, with no lattice-scale geometric law. REQUIREMENTS
% 061ccbe55ee items 11,19,22,25,26,27. [X] Continuum exterior: this author's
% radial principal-curvature and chart calculation agrees with the independent
% Gauss-law route in strand_E_declared_equations.tex E11,E15,E16, actual source
% and SOURCES read in <internal>/.
% Its E18 agrees with the local sine ray Hamiltonian. It has no counterpart
% periodic-orbit coefficient. The normalized-averaging theorem and periodic-ray
% coefficients below remain [R]. The massless axial branch, negative-frequency
% group-zero locus and cubic near-surface relation are [X]: this author's
% inverse-flow monotonicity proof agrees with the independent endpoint-sign
% and upward-root proof in strand_F_shell_modes.tex FM1,FM6--FM9, with its
% review_F_shell_modes.tex, both actual complete sources and SOURCES read in
% <internal>/. The transferred argument
% uses only that source's N=psi=1 free sine plus two-step flow operator, which
% is exactly the present operator. No retired field, source, ruler, interaction,
% stiffness or temperature argument transfers. Its massive extension does not
% establish a massive light-orbit coefficient here.
% [X] Directional principal-ray, periastron and radial-delay coefficients
% have the named complete T7v67/T7e/T7g argument comparisons near the end;
% their point-coupling and numerical-production limits are explicit there.
% Predecessor lane_z/B11-007: original g30b source read; its cos(2ap)
% velocity dressing is retained. No induced exterior, ruler or stiffness is used.
% The same-pair independent RED derived the orbit from its Hamiltonian
% linearization, agreeing with all dispersion and endpoint-ordering harmonics.
% Actual strand_E_hydrodynamic_sources.tex HA3,HA10, its review and scope/SOURCES
% in <internal>/solve were read. Their adjoint-identity route
% agrees with integral preservation for a FIXED bare source array. The reference
% can change with ell even at fixed physical C, so that result does not fix
% the charge along a coupled branch. The metric/ray comparisons here fix E.
% CURRENT BINDING SCOPE: the fully reread goal uses the coherent -J row,
% spherical source projection, adjoint back-averaging and a continued
% occupied SPECTRAL reference commuting with h. The fixed prescribed MP8
% packet/core/protocol below is unchanged. Its point-coupled propagation
% is a control for the binding finite-ell matter coupling; replacing the
% geometric flow by its l=1 back-average is mandatory for that comparison.
% Geometry at fixed E and matter propagation at finite ell are distinct
% statements. The earlier reference/sign alternatives retain only their
% explicitly proved control scopes.
%
% CONTINUUM GEOMETRY. For radial v, the principal curvatures of Kplus are
% v'/N,v/(Nr),v/(Nr). The coherent momentum row is
% v N'/(4pi G r N^2)=-J_r and the energy row is
% (r v^2)'=8pi G N^2 r^2 rho. On the J=0 connected branch N=1 by the
% asymptotic clock. For a fixed bare source array, normalized compact W_ell preserves
% E=int rho_ell d^3x=int rho d^3x exactly. Outside source+averaging support,
% rv^2=2GE=r_s. If N varies in the source, int rho is not automatically the
% boundary charge: the energy row contains N^2. The fixed-E exterior remains
% the appropriate comparison, but the source matching must then be supplied.
% The inward branch v=-sqrt(r_s/r)n gives ds^2=-dt^2+|dx-vdt|^2 and the
% Schwarzschild chart. Isotropic coordinates give beta=gamma=1 exactly.
% At fixed boundary E there is no a-dependent vacuum metric coefficient.
% For this metric the zero outward normal speed and zero outgoing expansion
% both occur at r_s. No lattice-scale geometric derivative has been inserted.
%
% EXACT WEYL SYMBOL AT POINT COUPLING. At N=1, h0(p) is constant
% in x. The Weyl symbol of {v_i,Q_i}/2, whose hopping length is 2a, is
% Q_i(p)[v_i(x+a e_i)+v_i(x-a e_i)]/2 exactly in a smooth interior.
% Thus the upper massless Weyl band is
% omega_W=E(p)+sum_i Q_i(p) bar v_i(x), bar v_i=[v_i(x+a e_i)+v_i(x-a e_i)]/2.
% This formula defines the symbol-ray observable used below; it is not a claim
% that a polarized finite-wavelength quantum packet obeys the scalar equations
% without Berry, interband or finite-width corrections.
% Put C=sqrt(r_s), v0_i=-C x_i r^(-3/2). The endpoint average is
% bar v_i=v0_i+a^2 delta v_i+O(a^4 C r^(-9/2)),
% delta v_i=(1/2)partial_i^2 v0_i
% =C r^(-5/2)[(9/4)n_i-(21/8)n_i^3].
% The scalar Hamiltonian through a^2 is H0+a^2(h_ord+h_disp), where
% H0=p-u p_r, u=sqrt(r_s/r), pvec=p_r n+(L/r)t, p=|pvec|,
% h_disp=-sum_i pvec_i^4/(6p)+(2u/3)sum_i n_i pvec_i^3,
% h_ord=delta v.pvec.
%
% CONSTRUCTED PERIODIC LIGHT ORBIT. Reflection symmetry keeps z=p_z=0 invariant.
% For positive circulation in the xy plane and fixed positive asymptotic omega,
% the continuum orbit has r0=3r_s/2 and L0=sqrt(3)r0 omega.
% The radial Poincare multipliers are exp(+/-2pi), so its small-perturbation
% periodic continuation is unique locally. Cubic matter dispersion generally
% gives a fourfold deformation, not an exact circle or conserved angular L.
% Write n=(cos phi,sin phi), t=(-sin phi,cos phi). Then
% delta v_r=C r^(-5/2)(9-21cos4phi)/32,
% delta v_t=C r^(-5/2)(21sin4phi)/32.
% Choose units r0=L0=1, r_s=2/3, p_r0=sqrt2, p0=sqrt3,
% omega0=Omega0=1/sqrt3. Derivatives of H0 divided by Omega0 are
% H_pp=1/3, H_pr=5sqrt2/6, H_pL=-sqrt2/3, H_rr=7/6, H_rL=-5/3.
% With r=1+a^2 R, p_r=sqrt2+a^2 P and L=1+a^2 L1, fixed energy gives
% L1=-sqrt3 h. Eliminating P,L1 from the phi-parameter Hamilton equations,
% R''-R=F, F=sqrt3[-h_r/3+(5sqrt2/6)h_p+h_{p phi}+(sqrt2/3)h_phi],
% where h=h_ord or h_disp and partial derivatives precede base evaluation.
% The separate forcing and periodic solutions are
% F_disp=19/4-(77/12)cos4phi-(11sqrt2/3)sin4phi,
% R_disp=-19/4+(77/204)cos4phi+(11sqrt2/51)sin4phi;
% F_ord=15/16-(7/16)cos4phi+(329sqrt2/64)sin4phi,
% R_ord=-15/16+(7/272)cos4phi-(329sqrt2/1088)sin4phi.
% Restoring units gives the explicit orbit
% r_W(phi)/r0=1+(a/r0)^2 R_ord(phi)+3(a omega)^2 R_disp(phi)+...
% and its azimuthal mean 1-(5/12)(a/r_s)^2-(57/4)(a omega)^2+...
% Remainders of this defined analytic symbol-ray expansion are quartic in
% a/r_s and a omega, including mixed terms; require a/r_s<<1 and a omega<<1.
% Using a frozen local Bloch symbol instead drops R_ord; it is a different
% approximation. Geometrical optics additionally needs omega r_s>>1.
% Berry/interband terms of orders (omega r_s)^-1 and (omega r_s)^-2 can combine
% with a^2 omega^2 at order a^2/r_s^2. Therefore the displayed -5/12 is an
% endpoint-ordering coefficient of the specified scalar symbol, not a complete
% polarization-independent microscopic packet-orbit coefficient. A quantum
% orbit observable must specify state, polarization and its measured center.
% Frequency dependence already disproves a universal scalar light-ring shift.
%
% EXACT LOCAL AXIAL BRANCH THEOREM. Put z=ak in (0,pi/2), w=a omega in (0,1),
% u>0 and use omega=sin z(1-u cos z)/a. The exact Weyl version on an axis has
% the same formula with u replaced by its endpoint average, also positive.
% There is a branch with s=sin z in [w,1) satisfying
% u=(1-w/s)/sqrt(1-s^2). Its derivative is positive:
% du/ds=[s^3+w(1-2s^2)]/[s^2(1-s^2)^(3/2)]. For s<=1/sqrt2 the numerator
% is positive; for larger s it exceeds s^3+1-2s^2=(s-1)(s^2-s-1)>0.
% Thus u runs monotonically from 0 to infinity on this branch. Its group speed
% cos z-u cos2z=[s^3+w(1-2s^2)]/[s sqrt(1-s^2)] is positive throughout.
% This proves absence of a local group-velocity turning surface for that
% positive-frequency branch, not preparation of a quantum state in the core.
% A putative zero group speed at 0<z<pi/4 instead requires u=cos z/cos2z,
% giving omega=-sin^3 z/(a cos2z)<0. The formal fixed-k radius
% R/r_s=1-3(ak)^2+O((ak)^4) belongs to this negative-frequency test and cannot
% be reported as the positive-frequency trapping radius.
% At u=1 and small w, w=sin z(1-cos z)=z^3/2+O(z^5), hence
% z=(2w)^(1/3)[1+O(w^(2/3))],
% dot r=(3/2)(2w)^(2/3)+O(w^(4/3)). The near-surface dispersive region has
% fractional radial width O((a omega)^(2/3)); it is not a universal analytic
% a^2 shift of the continuum characteristic. These are point-coupling
% formulas at fixed E; binding finite-ell propagation instead uses the
% actual l=1 back-average, as specified in the binding comparison below.
%
% SHORT CHECK EVIDENCE. Actual declared_exterior_short_check.py and .json were
% read in full by author and RED. Sol's authorized remote SymPy 1.14.0 job
% used CPU1/nice10 and returned twelve exact zero residuals in 1.2784 seconds.
% Its dispersion forcing and R_disp checks support the formulas above. The
% other checks concern the distinct centered-derivative constraint control
% below; none establishes a finite-a geometric law or a quantum packet orbit.
%
% INDEPENDENT MICROSCOPIC PACKET PROGRESS [R]. This calculation was saved
% before reading or exchanging the RED's new microscopic-packet route.
% It keeps the full massless four-component Hamiltonian, with N=1 and the
% declared endpoint coupling. Choose a bounded smooth core extension of the
% fixed-E exterior, so h is self-adjoint on the native one-particle space.
% Define s_i=sin(ap_i)/a, C_i=cos(ap_i), D_i=cos(2ap_i),
% F=sum_j {v_j,Q_j}/2 and W_i={v_i,D_i}/2. Exact commutators give
% V_i=i[h,X_i]=alpha_i C_i+W_i,
% A_i=i[h,V_i]=2 sum_jk epsilon_ijk Sigma_k s_j C_i
%       +i alpha_i[F,C_i]+i sum_j alpha_j[s_j,W_i]+i[F,W_i].
% Every flow-ordering term is present. Only the first term contains Sigma;
% the others contain alpha or the spin identity. Thus no projected-band,
% Berry or interband contribution is being dropped from these exact initial
% velocity and acceleration observables.
%
% Prepare two normalized one-particle spinors psi_chi(x)=phi(x)u_chi, chi=+/-1,
% with the same scalar wavepacket phi. Let d_i=sin(ap0_i)/a and choose u_chi
% to have positive central Bloch energy, chirality chi, and polarization
% <alpha_i>=d_i/|d|, <Sigma_i>=chi d_i/|d|. Both preparations have exactly
% the same position/momentum probability laws, energy expectation, and initial
% centroid velocity under the full nonuniform Hamiltonian. Their acceleration
% difference is nevertheless exactly
% Delta A_i=4 sum_jk epsilon_ijk (d_k/|d|) <s_j C_i>_phi.
% This formula is independent of the core extension because the other three
% acceleration terms cancel between the two states. It is an exact statement
% about these full spinor preparations, not a pure positive-spectrum packet
% theorem: finite-width constant spinors generally contain both local bands.
% Positive central energy or positive mean Killing energy is not a substitute
% for a stated projection onto the actual unoccupied spectral subspace.
%
% For a narrow Gaussian momentum law with mean p0 and covariance V, the
% Brillouin-periodized preparation differs from the following Gaussian
% characteristic-function formula only by exponentially small band-face
% tails when a|p0|,a sqrt(||V||)<<1. For i!=j,
% <s_j C_i>=exp[-a^2(V_ii+V_jj)/2]/a times
% [sin(ap0_j)cos(ap0_i)cosh(a^2 V_ij)
%  -cos(ap0_j)sin(ap0_i)sinh(a^2 V_ij)].
% Consequently the chirality-dependent part of the full acceleration is
% A_i^spin=-chi a^2/|p0| sum_jk epsilon_ijk
%       [p0_j p0_k V_jj+2 p0_i V_ij p0_k]+O(a^4),
% with fixed p0,V; the central-dispersion terms cancel against the expansion
% of d/|d|. This is a closed full-matrix state-qualified lattice correction.
% In particular put the centroid on the x axis, p0=(sqrt2 k,k,0), and choose
% V_xx=V_yy=V_zz=b, V_xz=V_zx=t, all other off-diagonal entries zero,
% with |t|<b. The two radial centroid accelerations differ by
% Delta A_x=4 sqrt(2/3) a^2 k t+O(a^4).
% The initial centroid positions and velocities agree, so the same expression
% is their difference in d^2|<X>|/dt^2, including its centripetal term.
% At the continuum light-ring phase-space point r0=3r_s/2, k=sqrt3 omega,
% set t=eta/r0^2 and omega=kappa/r0. The difference is
% 4sqrt2 kappa eta a^2/r0^3+O(a^4), precisely a relative a^2/r0^2 effect.
% Choose 1<<sqrt(b)r0<<kappa and |eta|<b r0^2 to obtain a localized packet
% with narrow relative momentum spread; its mean energy is positive for this
% subcritical exterior preparation. Changing chi at fixed t changes the
% coefficient without changing the common mean energy or initial speed.
%
% Scope: this disproves a preparation-independent initial centroid-acceleration
% coefficient even for packets sharing the light-ring central phase-space
% data. It also demonstrates why finite-width/interband preparation terms can
% contribute at the same a^2/r_s^2 order as the scalar endpoint coefficient.
% It does not yet compute the radius of a periodic positive-spectrum quantum
% packet, nor assert that these constant-spinor preparations are an adiabatic
% circular orbit at a=0. Their exact dynamics can contain rapid interband
% transients. A claimed full light-ring coefficient must define its spectral
% preparation, polarization, packet covariance, center observable and averaging
% time; applying the scalar periodic-ray coefficient to all these preparations
% fails this exact commutator test. The positive-spectrum periodic-packet
% continuation and a full quantum trapping/scattering readout remain active
% work, separate from this completed initial-observable counterexample.
%
% COMPLETE MATRIX PROPAGATION AND ADMISSIBLE PREPARATION [R]. The actual
% redteam_declared_exterior.tex MP1--MP7 and all SOURCES were read after both
% independent packet saves. The author checked its stencil expansion,
% Duhamel sign, free-band projector, centering translation and root proof.
% These are named same-pair derivations, not other-lane [X] evidence.
% On smooth packets with fixed physical scales, away from zone faces,
% h_a=h_c+a^2 h_2+O(a^4),
% h_c=alpha.p+sum_i{v_i,p_i}/2,
% h_2=-sum_i alpha_i p_i^3/6-sum_i{v_i,p_i^3}/3.
% Its Weyl form is -alpha.p^3/6-(2/3)sum_i v_i p_i^3
% +(1/2)sum_i(partial_i^2 v_i)p_i. This unprojected matrix contains the
% endpoint term and all band/polarization response. A separate Berry term
% cannot be added without also transforming the Hamiltonian and detector.
% For a specified smooth detector F and preparation rho_a=rho_0+a^2rho_2,
% the complete finite-time coefficient is
% B_F(T)=Tr[rho_2 U_c(T)^dag F U_c(T)]
% +i int_0^T Tr[rho_0[h_2I(s),F_I(T)]]ds.
% A varying detector adds Tr(rho_0 F_2I(T)). Uniform Sobolev control (H^8
% suffices here), bounded smooth flow derivatives, and an actual controlled
% finite-time collar justify the integrated O(a^4) residual. This is an
% exact coefficient FUNCTIONAL, not an evaluated scattering or orbit result.
%
% A Pauli-admissible incoming-sea probe can be constructed explicitly.
% Take P_a=(1+alpha.s/|s|)/2 and a normalized smooth positive-continuum-band
% spinor phi_R with compact momentum support in one cone away from zero.
% With n=p/|p|, S4=sum_i p_i^4,
% P_a=P_0+a^2P_2+O(a^4), P_2=alpha.delta n/2,
% delta n_i=-p_i^3/(6|p|)+p_i S4/(6|p|^3), n.delta n=0.
% Set delta x_j=2Re<P_2 phi_R,X_j phi_R> and
% b_R=P_2 phi_R+i sum_j delta x_j p_j phi_R.
% The normalized and translated P_a phi_R then has native centroid R e_1
% and expansion phi_R+a^2 b_R. Its covariance
% C_a=P_-^a+|psi_a><psi_a| is admissible exactly, and under the common
% transport it remains P_T+|psi_a(t)><psi_a(t)|. This is a fixed-background
% probe, not a self-consistent object or an instantaneous spectral refilling.
% For its initial outward centroid drift,
% D_0(R)=<phi_R,(alpha_1+v_1)phi_R>,
% D_2(R)=<phi_R,[-alpha_1p_1^2/2-{v_1,p_1^2}]phi_R>
%        +2Re<b_R,(alpha_1+v_1)phi_R>.
% For a translated forward packet with c=<alpha_1> in (0,1], width b,
% 1/k<<b<<r_s, controlled tails, and all these scales fixed as a->0,
% D_0=c-sqrt(r_s/R)+O(b^2/r_s^2),
% D_0'=sqrt(r_s)/(2R^(3/2))+O(b^2/r_s^3).
% The actual continuum root exists uniquely in a fixed annulus around
% R_0=r_s/c^2[1+O(b^2/r_s^2)]. Its continued initial-drift root is
% R_a=R_0-a^2D_2(R_0)/D_0'(R_0)+O(a^4).
% This coefficient includes projection, centering, width and ordering; it
% fixes the carrier/profile, not mean frequency. It is still an INITIAL
% drift diagnostic, not a trapping barrier or a periodic/scattering orbit.
%
% EVALUATED NATIVE TRANSFER ACROSS A SURFACE [R]. This new calculation was
% saved before sending its argument to the RED. It concerns the unrestricted
% microscopic theory, not a low-energy expansion of a positive-frequency ray.
% Fix a>0, the massless native Hamiltonian above and a bounded smooth core
% extension. Let B be a nontrivial spatial interior and P_B its site projector.
% An exactly impermeable quantum surface would require
% (I-P_B)exp(-i h_a t)P_B=0 for every t>=0. Differentiating at t=0 gives
% (I-P_B)h_a P_B=0. Every nearest-neighbor link crossing the surface instead
% has the nonzero matrix +/-i alpha_i/(2a), independent of the two-step
% flow. The latter connects different endpoints and cannot cancel this
% nearest-link matrix. Thus no spatial cut is an exactly trapping surface
% for all microscopic preparations. This is an exact support statement at
% fixed a; it does not deny long-wave trapping or bound a low-energy signal.
%
% The transfer can be evaluated with admissible number-conserving Gaussian
% states and the declared common reference. Choose a normalized local mode
% b=|x,w>, ||w||=1, inside B, and B_x=|b><b|. For any incoming projector R,
% set C_-= (I-B_x)R(I-B_x), C_+=C_-+B_x. Both satisfy 0<=C_+/-<=I since
% C_- is supported on I-B_x and bounded there by that identity. They differ
% by occupation of this initially reset local mode; both use the same R
% and subsequently the same U. No forbidden addition to an already filled
% sea is assumed. Their excess exterior count at time t is exactly
% Delta N_out(t)=Tr[(I-P_B)U(C_+-C_-)U^dag]
%              =||(I-P_B)U b||^2
%              =t^2||(I-P_B)h_a b||^2+O(t^3||h_a||^3).
% This is an actual finite-time detector probability, not an initial
% acceleration or an unevaluated Duhamel functional. It is a prescribed-
% geometry probe comparison; source backreaction is not inferred absent.
%
% For an explicit spherical detector choose x=q e_1, q an integer multiple
% of a, B={|y|<R}, R=q+a/2, R>5a. Only x+a e_1 and x+2a e_1 are exterior
% endpoints reachable in one application of h_a. Orthogonality of these
% endpoint states gives, for every polarization w,
% Delta N_out(t)=t^2/a^2[1/4+(v_1(x)+v_1(x+2a e_1))^2/64]+O(t^3||h_a||^3).
% Both endpoints can be in the fixed-E Schwarzschild collar, where
% v_1(q e_1)=-sqrt(r_s/q). At a detector near R=r_s this coefficient tends
% to 5/(16a^2) as a/r_s->0. Equivalently t=tau a gives a leading transfer
% 5 tau^2/16 for tau<<1, with the exact displayed finite-a flow coefficient.
% The norm bound is controlled at fixed a; for a family with bounded flow,
% ||h_a||=O(1/a), so the scaled-time remainder is O(tau^3).
% This deliberately local preparation has ultraviolet momenta and no
% fixed positive frequency. It rules out an exact microscopic support
% horizon or a universal shifted one, while leaving the physically separate
% low-energy, band-prepared finite-time ring/scattering calculation active.
%
% The same reset construction also makes the earlier chirality-acceleration
% comparison admissible relative to a filled incoming sea, without falsely
% projecting its constant-spinor packets onto an upper band. The two chirality
% modes psi_+,psi_- are orthogonal. Put B=B_++B_-, B_chi=|psi_chi><psi_chi|,
% C_base=(I-B)R(I-B), C_chi=C_base+B_chi. These are positive contractions
% with common reference R. Their physical excess difference is B_+-B_-,
% so the common reset background cancels in the equal-energy, initial-velocity
% and acceleration comparisons. This precise admissibility extension was
% proposed independently by the RED after reading the native-transfer proof
% and checked here by compression and orthogonality. It remains a two-mode
% reset preparation with particle/hole changes, not a pure added upper-band
% packet or a periodic state. No claim about its self-consistent source is made.
%
% EVALUATED FULL-SPINOR WEAK-FIELD TRANSITION KERNEL [R]. This derivation
% was saved before exchanging it with the RED. It is a distinct scattering
% control, not a calculation at the strong-field light ring. Let
% h(lambda)=h_0,a+lambda V_a, v(lambda)=lambda grad Phi, so outside the
% smooth radial core Phi=-2sqrt(r_*)sqrt(r) and r_s=lambda^2 r_*. The object computed
% here is the on-shell first derivative in lambda (the Born transition
% kernel), with two positive free-band states at the SAME exact lattice
% energy. No scalar-band replacement is made in computing this derivative.
% Use <x|p>=(2pi)^(-3/2)exp(ip.x), and report the Fourier kernel before its
% overall factor (2pi)^(-3). All momentum transfers below avoid q=0 and
% all external states lie in a fixed cone away from zero and zone faces.
%
% A useful exact operator subtraction exposes the long-range pure-gradient
% part. Put B={Phi,h_0,a}/2. Since h_0,a^2 is scalar in spin space,
% i[B,h_0,a]=-i[h_0,a^2,Phi]/2. Consequently conjugation by exp(i lambda B)
% removes V_g=i[h_0,a^2,Phi]/2 to first order in lambda. Define the residual
% R_a=V_a-V_g. This is an exact finite-parent identity; on the infinite
% lattice it is first an identity on finitely supported spinors. An
% asymptotic scattering construction for the long-range channels is a
% separate domain/limit statement, not inferred merely from this identity.
% On an exact equal-energy pair the matrix element of V_g vanishes, so
% V_a and R_a have precisely the same elastic first-order kernel.
%
% The exact sampled Fourier vertex is
% R_a(p',p)=sum_i[Q_i(p')+Q_i(p)]v_tilde_i^a(q)/2
%            -i[E_a(p')^2-E_a(p)^2]Phi_tilde^a(q)/2.
% Writing q=p'-p and P=p'+p, its low-cone unaliased part is
% i Phi_tilde(q)/(2a^2) sum_i sin(aP_i)
%                         [a q_i cos(aq_i)-sin(aq_i)].
% The identities Q_i(p')+Q_i(p)=sin(aP_i)cos(aq_i)/a and
% E_a(p')^2-E_a(p)^2=sum_i sin(aP_i)sin(aq_i)/a^2 are exact. For a FIXED
% smooth radial core, Poisson aliases in v_tilde_i^a-iq_i Phi_tilde^a are
% smaller than any prescribed power of a on the retained compact momentum
% set. Repeated derivatives of the smooth core/exterior are integrable from
% sufficiently high order onward, giving that Fourier decay. No exact
% continuum-derivative Fourier identity is assumed for a sampled field.
% Thus, with these alias remainders retained,
% R_a=a^2 R_2+O(a^4), R_2=(1/6)sum_i{p_i,partial_i^3 Phi},
% R_2^W=(1/3)sum_i p_i partial_i^3 Phi,
% R_2(p',p)=-(i/6)Phi_tilde(q)sum_i P_i q_i^3.
% The remainder is controlled on fixed smooth momentum-support test packets.
% This residual contains the complete first-lambda, order-a^2 effect:
% corrections to the external free-band spinors multiply R_a and begin only
% at a^4. The actual equal-energy radial momentum adjustment is retained by
% the exact V_g cancellation. Substituting continuum equal radii into the
% un-subtracted V_a before expanding would miss that adjustment.
% In real space the oriented block <x+2ae_i|R_a|x> is
% i[partial_i Phi(x)+partial_i Phi(x+2ae_i)
%       -(Phi(x+2ae_i)-Phi(x))/a]/(8a),
% which is O(a partial_i^3 Phi) and decays as r^(-5/2) in the exterior.
% The long-range pure-gradient vertex has therefore been removed from this
% first-order kernel; no long-range wave-operator theorem is asserted here.
%
% The complete positive-band spin factor is explicit. For fixed chirality
% chi, choose normalized massless free eigenspinors u_chi(n). Then
% T_1,a(p',p)=-(i a^2/6)Phi_tilde(q)
%             [sum_i(p_i'+p_i)(p_i'-p_i)^3]
%             u_chi(n')^dag u_chi(n)+O(a^4),
% while opposite chiralities do not mix. Its modulus squared spin factor
% is (1+n'.n)/2. No Berry, interband or endpoint term is omitted: this is
% a direct matrix element of the exact microscopic perturbation between
% its native free eigenstates. Actual transitions at finite lambda also
% have order-lambda^2 gravity terms, which are not computed by this derivative.
%
% For the pure exterior distribution, the Abel Fourier transform is
% Phi_tilde(q)=3sqrt2 pi^(3/2)sqrt(r_*) |q|^(-7/2), q!=0.
% Indeed 4pi/q int r Phi(r)sin(qr)dr uses
% int r^(3/2)sin(qr)dr=Gamma(5/2)sin(5pi/4)q^(-5/2)
% after an Abel regulator. A smooth core of radius r_c adds an analytic
% compact-support transform; for a fixed bounded scaled core its relative
% correction is O((|q|r_c)^(7/2)) as |q|r_c->0. That core term is retained
% at finite |q|r_c; the unmodified singular exterior is not silently used
% as a complete microscopic object.
%
% An evaluated orientation contrast is useful. Set the common continuum
% energy k, n=e_1, n'=(1/2,sqrt3/2,0), and take the exact lattice radial
% momenta at that energy before expanding. Then |q|=k and
% sum_i P_i q_i^3=3k^4/8. In the coplanar spinor gauge with overlap sqrt3/2,
% T_1,a= -i[3sqrt6 pi^(3/2)/32]a^2 sqrt(r_* k)
%         [1+O((k r_c)^(7/2))]+O(a^4).
% Rotating BOTH directions to n=(sqrt3/2,-1/2,0),
% n'=(sqrt3/2,1/2,0) keeps the energy, scattering angle and chirality but
% gives P_i q_i=0 for each i. For the sampled radial core/exterior, reflection
% symmetry makes v_tilde_x^a=v_tilde_z^a=0 at this purely y transfer, while
% Q_y(p')+Q_y(p)=0. The exact native vertex therefore vanishes at every a
% for the correspondingly equal exact radial momenta, including aliases. The same
% native on-shell Born kernel is therefore not rotation independent even
% in this spherically symmetric exterior. This is an evaluated low-energy
% full-spinor order-a^2 transition coefficient and a precise cubic-orientation
% counterexample, not an evaluated finite-time light-ring packet, a full
% scattering probability, or a theorem about the solved object's core.
%
% EVALUATED LOW-ENERGY FINITE-TIME RING NONUNIQUENESS [R]. This new
% argument was saved before sending it to the RED. It goes beyond the
% initial-drift and UV-transfer controls: it gives an explicit difference
% between the order-a^2 radii of two admissible low-energy one-circuit
% packet experiments, including their full native matrix propagation.
% It excludes a coefficient determined only by the continuum packet and
% mean frequency. It does not claim to evaluate an absolute coefficient
% for one completely fixed preparation through order a^2.
%
% First fix the actual return experiment from the independently saved RED
% MP8, whose full source/SOURCES and primary Egorov hypotheses were read.
% Its free-positive, reflection-symmetrized mixed probe has variables
% z=(R,p_r,L,T), fixed mean frequency omega, and conditions
% F_a=(r_a(T/2)-R,r_a(T)-R,y_a(T),Tr(sigma_a h_a)-omega)=0,
% with x_a(T)>0 and one winding. The author independently verified the
% full classical radial Hessian and the nonzero determinant
% -2H_pp sinh(pi)[cosh(pi)-1]/lambda, lambda=1/(sqrt3 r_*).
% On a compact gapped phase-space tube, the finite-time native scalar-
% detector Egorov estimate and coherent-state derivative bounds give C^1
% convergence to that classical return map. The actual continuum packet
% root z_0 exists for sufficiently small fixed epsilon=(omega r_*)^(-1).
% The exp(2pi) instability constant must be controlled in the localization
% estimates; no moderate-frequency accuracy is inferred. With epsilon then
% held fixed and a omega<<1, the full matrix expansion and implicit-function
% theorem give an actual root z_a and mean centroid radius
% Rcal_a=Rcal_0+a^2 Rcal_2+O(a^4),
% Rcal_a=T_a^(-1)int_0^{T_a}|Tr(sigma_a X(t))|dt.
% This is the time mean of the NORM OF THE CENTROID in the fixed lattice
% Cartesian frame, not Tr(sigma_a |X|) and not a radial density peak. The
% native X detector, fixed y=0 return section and observation-time convention
% are unchanged between experiments; only the preparation is rotated.
% This names a finite-time position return, not recurrence of the complete
% quantum state. MP8's unevaluated Rcal_2 functional is retained as such.
%
% Now let J_z be the continuum total rotation generator, orbital plus
% Sigma_z/2, and choose any fixed real eta. In the same low-momentum
% embedding define theta_a=eta a/r_s and rotate the initial packet by
% U_+=exp(-i theta_a J_z) and U_-=exp(+i theta_a J_z), with equal weights.
% Project each component back into the SAME exact native free P_+^a and
% normalize it. This is a new specified preparation, not a change to h_a
% or to the native detector X. Rotations are defined on the smooth retained
% momentum support; fixed-time sampling/high-momentum remainders are bounded
% as in the complete matrix expansion. Both mixtures have trace one, are
% supported in P_+^a and are positive contractions. Therefore both
% P_-^a+sigma_a and P_-^a+sigma_a^{(eta)} are Pauli-admissible incoming
% covariances, with the same sea and its common subsequent transport.
%
% The a=0 operator h_c is exactly rotation invariant for the chosen radial
% core/exterior: [h_c,J_z]=0. Also [P_0,J_z]=0. With ad_J A=[J,A],
% sigma_a^{(eta)}=sigma_a-(eta^2 a^2/(2r_s^2))ad_J^2 sigma_0+O(a^4).
% The double-commutator term stays inside P_0. Exact P_+^a reprojection
% therefore changes this expansion only at O(a^4); the individual O(a^3)
% off-band changes cancel in the equal mixture. This verifies the native
% Pauli constraint rather than silently rotating an occupied sea subspace.
%
% Propagating this DIFFERENCE needs only U_c at order a^2, since h_a-h_c
% already carries a^2. Thus the full lattice matrix correction h_2 and the
% original preparation correction are identical in the two experiments
% and cancel in their difference. For the native position,
% [J_z,[J_z,X_x]]=X_x, [J_z,[J_z,X_y]]=X_y,
% [J_z,[J_z,X_z]]=0. Their centroid has z=0 by the retained reflection
% preparation. Hence, uniformly over the specified finite observation time,
% x_a^{(eta)}(t)=[1-eta^2 a^2/(2r_s^2)]x_a(t)+O(a^4),
% r_a^{(eta)}(t)=[1-eta^2 a^2/(2r_s^2)]r_a(t)+O(a^4).
% The mean energy is unchanged through order a^2 because
% Tr(h_c ad_J^2 sigma_0)=Tr([J_z,[J_z,h_c]]sigma_0)=0.
% All quantities here are full matrix expectations, so no omitted spin,
% band-gradient or endpoint term can restore equality of these two radii.
%
% Label the new family by its ACTUAL initial centroid
% R^{(eta)}=[1-eta^2 a^2/(2r_s^2)]R+O(a^4), not the unrotated carrier's
% label. The two radial return conditions and final y condition are then
% multiplied by that same factor, while the fixed-energy condition changes
% only at O(a^4). The nondegenerate actual packet root moves by O(a^4)
% beyond this explicit relabeling. Its evaluated mean-radius coefficient is
% Rcal_2^{(eta)}-Rcal_2= -eta^2 Rcal_0/(2r_s^2),
% equivalently
% Rcal_a^{(eta)}/Rcal_0-Rcal_a/Rcal_0
%       =-(eta^2/2)(a/r_s)^2+O(a^4).
% This is an explicit evaluated low-energy finite-time packet coefficient
% difference, not a formal Duhamel answer. The two families have the same
% continuum limiting packet, mean frequency and one-circuit return protocol;
% they differ by a specified subleading angular preparation variance.
% The remaining p_r,L variables label the UNROTATED preparation before the
% symmetric mixture; they are not silently reinterpreted as the mixture's
% measured carrier. Holding the same actual launch centroid instead of this
% explicit relabeling would require a further state adjustment and would
% define a different coefficient.
% Their individual mode momenta remain in the same low-energy cone, with
% rotation-induced changes O(a omega/r_s), not ultraviolet localization.
% Consequently continuum packet data alone do not fix the requested
% microscopic light-ring coefficient. A completely fixed preparation to
% order a^2 still has its own definite Rcal_2, requiring the actual matrix
% propagation to evaluate; this theorem neither supplies that absolute
% value nor promotes a finite-time return into an eternally periodic state.
% Indeed every smooth radial density detector f(|X|), as well as
% X_x^2+X_y^2, commutes with J_z, so its order-a^2 difference under this
% preparation change is ZERO. The evaluated result is angular-mixture
% broadening of this centroid protocol; it is not a shift of the invariant
% radial distribution or an intrinsic fixed-preparation lattice coefficient.
%
% PRIMARY SCOPE CHECK for the preceding packet-root step: Panati, Spohn and
% Teufel, Space-Adiabatic Perturbation Theory, math-ph/0201055v3,
% https://arxiv.org/html/math-ph/0201055v3 , actual (CG)/(IG) assumptions,
% Theorem 2.1, Theorem 5.1, Corollary 5.2 and the native-position warning
% were read. They supply the gapped symbol/finite-time estimate, not this
% model's return Jacobian, core construction or angular-mixture coefficient.
% The local extension and packet parameter derivatives are model checks;
% no electromagnetic Dirac coefficient is imported from that article.
%
% FIXED-PREPARATION RADIAL-DETECTOR ROUTE [R; specified readouts N below]. This
% independent continuation was saved before exchanging its new reduction.
% It supplies a tractable full-matrix coefficient for a precisely defined
% finite-time trapping detector; it is not the centroid ring experiment above.
% The native h_2 is a cubic invariant whose angular decomposition contains
% a rotational scalar and rank four. Averaging an incoming single kappa
% multiplet over all m makes the rank-four contribution vanish in every
% first-order expectation of a radial detector: h_c and that detector
% preserve j,m,kappa, and the trace of a non-scalar irreducible tensor over
% the full j multiplet is zero. This is angular trace selection, not an
% assertion that each microscopic angular channel evolves independently.
%
% Explicitly, the rotational average of sum_i e_i^a e_i^b e_i^c e_i^d is
% (delta_ab delta_cd+delta_ac delta_bd+delta_ad delta_bc)/5. Hence
% (h_2,kin)_scalar=-(alpha.p)p^2/10,
% (h_2,flow)_scalar=-[(v.p)p^2+p^2(p.v)]/5.
% For psi=(f(r)Omega_{kappa m},i g(r)Omega_{-kappa m})/r, inner product dr,
% p_r=-i partial_r as a differential expression, the full continuum radial
% matrix on the regular full-Dirac domain is
% D_kappa=sigma_y p_r+(kappa/r)sigma_x,
% h_c,kappa=D_kappa+{v,p_r}I/2,
% L_kappa=D_kappa^2=p_r^2+(kappa^2+kappa sigma_z)/r^2.
% The scalar operator v.p becomes A=v p_r+i v/r, with
% A^dag=p_r v-i v/r. Thus the COMPLETE order-a^2 radial-detector response
% for the m-averaged preparation uses
% h_2,kappa=-D_kappa^3/10-(A L_kappa+L_kappa A^dag)/5.
% The exact endpoint allocation remains in the latter anticommutators.
% No adiabatic band substitution or discarded polarization term is involved.
%
% NATIVE STATE AND READOUT. This packet occupies one low-momentum valley of
% the one lattice fermion; no degeneracy multiplier is introduced. Embed
% the compact-momentum continuum spinors by lattice sampling, project each
% with the actual free upper-band projector P_+^a, and normalize it. Call
% the resulting modes psi_m^a and set sigma_a=(1/40)sum_m|psi_m^a><psi_m^a|.
% It is a positive trace-one operator with range in P_+^a, hence sigma_a<=P_+^a
% even if the projected magnetic modes cease to be exactly orthogonal.
% The incoming covariance C_in=P_-^a+sigma_a is therefore Pauli admissible.
% Use the reference R_in=P_-^a and evolve C and R with the SAME static
% microscopic h_a in the stated prescribed background. Their difference is
% U_a sigma_a U_a^dagger; no instantaneous refilling or reference variation
% is hidden in the probe measurement. The detector is the site multiplication
% F(|X|), or |X| for the mean, and the measured excess is
% Tr[F(|X|)(C(T)-R(T))]. The continuum radial integral computes its a^2
% coefficient. Smoothness, compact incoming momentum support, regular-center
% powers and finite T control the sampling remainder at this order.
% The probe's negligible-backreaction limit is part of this prescribed-field
% experiment; no self-consistent microscopic source for the core is asserted.
%
% For a positive free radial preparation, the scalar part of the native
% free-projector correction P_2 vanishes: -D_kappa^3/10 commutes with
% D_kappa and changes no spectral projector away from zero. Its rank-four
% part has zero m-averaged contribution to a radial readout, also after
% propagation by h_c. Thus the native free-band preparation correction
% contributes zero to these first-order radial expectations. The actual
% projected native states still exist; their angularly off-diagonal
% corrections have not been declared absent as states.
%
% A completely specified probe for an evaluated check is as follows, with
% r_s=1, kappa=-20, j=39/2 and equal weight 1/40 over m. Choose a regular
% core with C^11 cumulative mass and flow, with mass fraction
% f(r)=I_{(r/r_c)^2}(3/2,12) for r<r_c=1/2 and f=1 outside,
% v(r)=-sqrt(f(r)/r), N=1. I is the regularized incomplete beta function.
% It gives nonnegative compact density proportional to
% (1-r^2/r_c^2)^11 (C^10 at r_c), a regular center, and the exact declared vacuum
% exterior. C^11 regularity suffices for the finite-derivative response
% estimates here; no C-infinity core is asserted. This fixes a prescribed
% H2-compatible background, not a
% self-consistent microscopic core state.
%
% Let K=20, R_*=3/2, p_*=2K/R_* and Delta=8. For |z|<1 set
% g(p)=C exp[-1/(1-z^2)] exp[-i S_K(p,R_*)], z=(p-p_*)/Delta,
% S_K(p,R)=sqrt(p^2 R^2-K^2)-K arccos(K/(pR)), and set g=0 otherwise.
% C>0 normalizes int_0^infinity |g|^2 dp=1. Its support is strictly above
% K/R_*, so this real phase is smooth. Define the EXACT free positive
% radial packet by
% u_0(r)=int_0^infinity dp g(p)/sqrt(pi)
%             (p r j_{K-1}(p r),-p r j_K(p r))^T.
% Each displayed radial basis spinor is a +p eigenvector of D_{-K};
% spherical-Bessel orthogonality gives unit norm. The phase selects the
% outgoing localized branch near the classical light-ring radius,
% but its actual radial distribution and continuum detector value must
% be measured from u_0, not set equal to a classical circle by definition.
% Its radial carrier exceeds the circular one. This deliberately avoids
% the zero classical energy derivative at the exact circular phase point
% when p_* alone is used for frequency matching; no periodicity is assumed.
%
% Fix the mean coordinate frequency to its actual continuum value
% omega_0=<u_0,h_c u_0>. The native free-band-projected packet adjusts only
% the envelope parameter p_* at order a^2 to preserve that frequency.
% Let u_p=partial_{p_*}u_0 at fixed R_*,Delta,K. Since the support remains
% positive and C is translation independent, its integrand is the one above
% multiplied by 2z/[Delta(1-z^2)^2]. Set
% E_0'=2Re<u_p,h_c u_0>, E_2=<u_0,h_2,kappa u_0>,
% delta p_*=-E_2/E_0', provided the ACTUAL derivative is nonzero.
% Then u_2(0)=delta p_* u_p gives the compulsory frequency-matching term.
% E_0' must be computed and checked, not assumed from the classical carrier.
% In addition to a times the occupied momentum being small, the expansion
% requires a^2 |delta p_*| much smaller than Delta. The small but nonzero
% energy derivative can make this preparation condition more restrictive.
% Evolve i dot u_0=h_c u_0,
% i dot u_2=h_c u_2+h_2,kappa u_0 for 0<=t<=T=1/2.
% Define F(r)=exp[-(r-R_*)^2/(2b_D^2)], b_D=1/4, in the fixed radial frame.
% The actual finite-time annular-retention coefficient is
% P_F(a,T)=<u_0(T),F u_0(T)>+a^2 c_F+O(a^4),
% c_F=2Re<u_0(T),F u_2(T)>.
% Also record the radial mean coefficient c_r=2Re<u_0(T),r u_2(T)>.
% These are absolute coefficients for a specified low-energy preparation,
% frequency matching, detector and time. Their numerical values require
% actual matrix propagation with controlled discretization; the accepted
% and rejected readouts are distinguished in the evidence block below.
% They measure radial probability/mean, not the norm of a centroid, an
% eternal light ring, a universal horizon shift or a solved object's state.
%
% BOUNDED SHORT-CHECK SPECIFICATION (implementation/execution belongs to Sol).
% Target: evaluate the preceding c_F and c_r, including fixed-frequency
% preparation, by continuum full-matrix response. This is NOT a fitted
% native-spacing scan or a claim of a Cartesian radial decomposition.
% Use one CPU, no GPU, nice 10, a hard total wall budget of 180 seconds and
% a memory budget of 512 MiB; return unaccepted partial rows if exceeded.
% The sparse matrices are O(n) banded; the largest preparation quadrature
% uses at most 1024 by 2047 values. No many-body/eigensolver job is needed.
%
% Numerical continuum regulator: r_j=j dr, j=1,...,Rmax/dr-1, Rmax=8,
% dr=1/64,1/128,1/256. The derivative is the skew fourth-order centered
% matrix (u_{j-2}-8u_{j-1}+8u_{j+1}-u_{j+2})/(12dr), with zero extension
% at the two ends, and p=-i derivative. This is a regulator of the regular
% full radial Dirac operator, not a self-adjoint half-line momentum
% observable. The regular-center packet has f=O(r^20), g=O(r^21);
% retain that origin domain and test the finite-box regulator by refinement
% and the outer-domain comparison. Use the actual v values specified above.
% Build D=sigma_y tensor p+sigma_x tensor diag(kappa/r),
% hc=D+I tensor (v p+p v)/2, L=D^2,
% A=I tensor (v p+i diag(v/r)),
% h2=-D^3/10-(A L+L A^dag)/5.
% Using L=D^2 is a consistent Hermitian discretization of the continuum
% identity; its finite-dr errors are included in the refinement test.
% Check hc and h2 Hermiticity to relative 1e-12 before propagation.
%
% Prepare u and u_p by 512-point Gauss--Legendre quadrature in z in (-1,1),
% p=p_*+Delta z, using the exact spherical Bessel functions above. Compute
% C from the same quadrature with dp=Delta dz. Include sqrt(dr) in the
% vector so all matrix inner products are ordinary Euclidean ones. Record
% its pre-normalization finite-box vector norm n=||u_raw||, then set
% u=u_raw/n and differentiate that normalization in u_p: replace the raw
% derivative u_p by u_p/n-u Re<u,u_p/n>.
% Compute omega0, E0prime, E2, delta p_*, and initial u2=delta p_* u_p.
% Reject frequency matching if |E0prime|<1e-3 or its sign/magnitude is not
% stable under refinement; no classical derivative may be substituted.
% For magnitude stability require the two finest main-grid values to differ
% by less than 1e-3 times |E0prime_finest|, with the same sign on all main
% grids. The raw values remain subject to paired scientific review.
%
% To keep the driven exponential well scaled, put s=dr^2/(1+kappa^2) and
% propagate (u,s u2) with the sparse block generator
% -i [[hc,0],[s h2,hc]] from t=0 to T=1/2. scipy.sparse.linalg.expm_multiply
% with traceA=0 is suitable; it must handle this non-Hermitian triangular
% block, not call a Hermitian eigensolver on it. Read times 0,1/8,1/4,3/8,1/2.
% Undo s only when forming the response. At each time report norm, the
% response norm constraint 2Re<u,u2>, and the matched energy response
% 2Re<u,hc u2>+<u,h2 u>, which should remain zero.
% Require |norm-1|<1e-9 and both response residuals below 1e-8 times their
% natural sum-of-term magnitudes (with a floor 1). Record absolute residuals
% too, so cancellation is visible rather than hidden by a relative number.
%
% Record at least omega0,E0prime,E2,delta p_*, initial <r>, final <r>,
% P0F(T), c_F,c_r, max residuals, outer/inner probability tails and timings.
% Repeat the middle grid with 1024 quadrature points; require preparation
% L2 difference of the normalized preparation vectors <1e-8 and coefficient
% differences <1e-4 max(1,|c|).
% Repeat dr=1/128 on Rmax=12 (1535 sites) to bound outer-boundary and initial
% tail effects; require coefficient differences <1e-3 max(1,|c|).
% For each coefficient require the two finest main-grid values to differ
% by <3e-3 max(1,|c|), with a decreasing refinement error and an estimated
% residual dr^4 error <1e-3 max(1,|c|). The finest difference divided by 15
% is an asymptotic fourth-order estimate, not a bound implied by decreasing
% differences alone; report and review the observed ratio of successive
% differences. Report the raw rows, not only a fitted number. If these
% checks fail, keep the result unaccepted and state
% which bound failed; additional/full production belongs to solve.
% This check evaluates only the specified radial-detector protocol. It
% does not evaluate the distinct MP8 centroid-return coefficient.
%
% ACTUAL FIRST RADIAL CHECK AND BOUNDED CONTINUATION. The complete actual
% producer declared_exterior_radial_short_check.py and all five rows in
% results/declared_exterior_radial_short_check.json were read. Its 7.832 s
% wall time and 121792 KiB peak RSS satisfy the stated resource bounds.
% The three final c_F values are 271.0414647672,272.5283822540,272.6251971167;
% successive differences have ratio 15.35836. Domain and quadrature changes
% are 8.93e-5 and 8.53e-13. This supports c_F at the requested precision.
% The final c_r values are -291.2097908863,-290.5026030469,-290.0011440291;
% their difference ratio 1.41026 does NOT support the fine/15 estimate.
% Its continuum mean-radius differences also change sign. No converged c_r
% value is accepted from those rows, despite the program's threshold flags.
% All propagation invariants pass, and the positive E0prime values
% .03951253518,.03983508454,.03985564222 meet the explicit finest-grid test.
% Neither good conservation nor small quadrature/domain error proves radial
% spatial convergence. The two observables have separate acceptance decisions.
%
% Exact continuation preset, checked with the same RED before execution:
% preserve the original five-row output; use the same producer and one_grid
% formulas, state, detector, core, time and frequency adjustment. Run ONLY
% (dr,Rmax,quadrature)=(1/512,8,512),(1/1024,8,512), under one total hard
% 180 s / 512 MiB cap, one CPU, nice 10, CUDA disabled. The largest Bessel
% matrix is 512 by 8191; sparse propagation remains banded. Save a distinct
% small continuation result, retaining partial rows as unaccepted if capped.
% Keep all original per-row diagnostics. At each saved time additionally
% partition c_r=2Re<u,r u2> into r<.25, .25<=r<.5, and r>=.5, and verify
% their sum equals the total to 1e-10 times max(1,sum absolute parts).
% These partitions diagnose, rather than assume, inner-region sensitivity.
%
% Compare the continuation with the original finest rows. Require unchanged
% Hermiticity/conservation bounds and positive stable E0prime, with its new
% finest difference <1e-3 |E0prime_finest|. For EACH observable separately,
% retain the <3e-3 relative finest-difference and <1e-3 estimated-residual
% requirements (scale max(1,|c|)). A fourth-order estimate additionally needs
% the ratio of the 1/256-to-1/512 and 1/512-to-1/1024 differences in [8,24],
% unless both differences are already below 1e-6 times that scale, in which
% case report their absolute smallness without inferring an order. Report
% actual rows and ratios; no asymptotic estimate is a rigorous error bound.
% If c_r remains unresolved it stays unaccepted; this specification does not
% authorize further refinement or change any preparation to force acceptance.
% The absolute MP8 centroid-return coefficient remains a different task.
%
% EVALUATED NATIVE RADIAL RESPONSES [N]. The author and the original RED read
% the full actual producer, baseline output and continuation output
% results/declared_exterior_radial_fine_check.json, including every saved
% time, partition, normalization, Hermiticity and response invariant. The
% continuation took 73.310 s and 304092 KiB peak RSS; no resource limit was
% relaxed. These are full-matrix a^2 coefficients of the native observable
% under the angular-trace theorem above, not scalar-symbol trajectory fits.
% Restoring r_s and writing epsilon_a=a/r_s, the specified annular detector
% at T=r_s/2 gives
% P_F(a,T) ~= .74246 + 272.6 epsilon_a^2 + O(epsilon_a^4).
% The continuum coefficient and a^2 coefficient are numerical estimates at
% fixed preparation/core/detector, not exact decimal identities. The actual
% finest c_F values are 272.6314231854 and 272.6318227845; the last two
% difference ratio is 15.58078 and the asymptotic fine/15 estimate is 2.66e-5.
% The independent middle-grid domain and quadrature changes were 8.93e-5
% and 8.53e-13. We quote 272.6, within the specified empirical precision;
% no certified interval follows solely from these convergence tests.
% The fixed mean coordinate frequency has omega_0 r_s ~=7.996734, obtained
% from the preparation rather than imposed as the round number 8.
% At the finest grid E0prime=.03985701435, E2=3565.619538 and
% delta p_*=-89460.27684 in r_s=1 units. The preparation's large frequency-
% adjustment derivative makes epsilon_a^2*89460 much smaller than 8 a
% separate requirement of this expansion. All native band and endpoint
% ordering contributions that survive the rotational trace are included.
%
% The predeclared earlier sample T=r_s/4 also resolves a radial-mean response:
% <|X|>(a,T)/r_s ~=1.5130 -132.0 epsilon_a^2+O(epsilon_a^4).
% Its five c_r values are -132.4768252592,-132.0507098999,-132.0241534647,
% -132.0224889900,-132.0223880229. The last difference is 1.00967e-4,
% with ratio about16.5; the original domain difference is 2.18e-4 and the
% quadrature difference is 5.3e-12. This is an explicitly earlier-time
% readout of the SAME preparation, not a substituted success at the target
% time. Its continuum mean is about1.512999 r_s. No new execution was used.
%
% In contrast c_r(T=r_s/2) remains UNACCEPTED. Its newest values
% -286.4030429761 and -280.6751793899 differ by 5.72786, with ratio .62818.
% The r<.25 contribution changes 1.469389->7.322506, the middle part
% -9.973099->-10.241312, and r>=.5 changes -277.899334->-277.756374.
% These observations locate most of the measured refinement sensitivity
% in the inner region; they do not establish its physical or numerical cause.
% Passing norm/energy identities does not repair spatial convergence.
% No limiting target-time mean-radius coefficient is inferred, and no
% further refinement is specified here. The accepted radial measurements
% are not a self-consistent object, universal trapping-surface displacement,
% fixed-preparation centroid-return coefficient or eternal periodic packet.
%
% SEQUENTIAL LIGHT-RING COEFFICIENT [R, independent continuation]. This
% calculation was saved before reading the RED's new sequential-limit
% argument. It evaluates a leading limit of the MP8 coefficient, not its
% absolute value at fixed finite frequency and packet width.
% Use the actual MP8 native-free-band projected/recentered, reflection-paired
% coherent family, with a fixed scaled smooth envelope and polarization,
% no extra spacing-dependent preparation deformation, and the same native
% centroid and position-return section. Let epsilon=(omega r_s)^-1 and
% zeta=(a omega)^2. First expand at a=0 with epsilon fixed; only afterwards
% take epsilon->0 in the sufficiently localized one-circuit regime.
% Write Rcal_a=Rcal_0+a^2 Rcal_2+O(a^4). Then
% lim_{epsilon->0} Rcal_2/[omega^2 Rcal_0]=-57/4.
% This is an evaluated sequential limit for that class of preparations;
% it neither assigns -57/4 to a finite-width packet nor determines its
% order-a^2/r_s^2 coefficient.
%
% Derivation. Set xbar=x/r_s, k=p/omega, tbar=t/r_s. Quantization has
% parameter epsilon. Use the actual one-valley Fourier embedding and native
% sampling of MP4/MP8: compact smooth momentum preparation is sampled
% isometrically for small a, and smooth-field alias/cutoff remainders are
% controlled to higher order before taking the spacing derivative. No
% unrelated finite parents at different spacings are compared.
% The common compact gapped tube must contain both bands' finite-time
% trajectories from this preparation. At zeta=0 both obey
% dot k=-(grad v)^T k. Bounded derivatives of the fixed smooth core keep
% initially nonzero k bounded away from zero for the fixed time; sufficiently
% small zeta preserves that bound. This also controls the small lower-band
% component if it enters the core. No global massless spectral gap is used.
% On this common tube, the exact native
% Weyl matrix divided by omega has the smooth two-parameter expansion
% hbar=alpha.k+v.k+zeta[-sum_i alpha_i k_i^3/6
%                         -(2/3)sum_i v_i k_i^3]
%       +(zeta epsilon^2/2)sum_i(partial_xbar_i^2 v_i)k_i
%       +O(zeta^2),
% with the endpoint and matrix terms retained before taking either limit.
% Its upper principal eigenvalue is
% H_+=|k|+v.k+zeta[-sum_i k_i^4/(6|k|)
%                         -(2/3)sum_i v_i k_i^3]+O(zeta^2).
% The endpoint term is zeta epsilon^2, not an additional leading zeta
% term. The gap 2|k| stays uniformly positive on the retained tube for
% small zeta, independently of epsilon. All symbols, cutoffs and the
% specified preparation have bounded finite zeta derivatives there.
%
% The finite-time band construction can be differentiated in zeta for
% native scalar position measurements. To justify this step, construct the
% superadiabatic projector/intertwiner to more orders than the derivatives
% used: the homological inverse is bounded by the same gap, so its symbol
% coefficients and their zeta derivatives remain uniform. A zeta derivative
% of the propagator can lose one epsilon power; an arbitrarily higher
% finite-order residual absorbs it on the fixed one-orbit interval. The
% incoming free-band preparation differs from the transported band by
% O(epsilon). For a native scalar detector its interband block is also
% O(epsilon), so the cross-band expectation is O(epsilon^2); differentiating
% its rapid relative phase loses one power and still gives O(epsilon).
% Within the degenerate upper band, the leading scalar detector is the
% classical pullback times the spin identity. The spin transport therefore
% contributes only subleading terms, with its full native value retained
% in the finite-epsilon coefficient rather than discarded.
% More explicitly, the block-Egorov scalar observable can be constructed
% through its order-epsilon correction with an O(epsilon^2) evolution
% remainder and uniformly smooth zeta-dependent symbols. Differentiating
% the remainder costs at most epsilon^-1 and leaves O(epsilon). The
% leading symbol is scalar in spin, so the order-epsilon spin transport
% cannot produce an unaccounted leading zeta derivative.
%
% Pairing the resulting symbols with the specified coherent family gives
% convergence of the position-return map and its zeta derivative to their
% classical counterparts. Center/carrier derivatives needed for J_0 can
% cost epsilon^-1/2; the higher-order construction and coherent moments
% still give a conservative O(sqrt(epsilon)) remainder. No zeta derivative
% of a velocity-return condition is required: MP8 uses positions only.
% The time derivative of the unperturbed map is its actual velocity and
% has the usual leading band symbol. The energy row is an initial matrix
% expectation and has the same smooth parameter limit. The native coordinate
% tails use the weighted finite-time estimates already specified in MP8.
% Instability constants are finite for one circuit but contain its exp(2pi)
% growth; epsilon must be small enough to control them.
%
% Thus F_zeta(epsilon,z_0)->F_zeta,classical and
% J_0(epsilon)->J_classical, with the latter invertible by the previously
% evaluated radial determinant and energy/section rows. The actual matched
% root derivative -J_0^-1 F_zeta and the derivative of its mean centroid
% radius converge to the corresponding classical quantities. This proves
% that the sequential coefficient is obtained from the FULL matrix response,
% not by assuming the scalar ray to be an exact microscopic orbit.
%
% The remaining classical derivative is already evaluated by the dispersion
% forcing above: the angle-mean fractional radius derivative is -57/4.
% Its identification with MP8's four classical conditions needs a check.
% The unique nearby hyperbolic periodic ray is invariant, as an orbit,
% under the cubic pi/2 rotation in the xy plane. Each successive rotated
% quarter therefore takes the same time. At half its period it has the
% same radius as at launch; at its full period it returns to the initial
% x axis with positive x. It satisfies all four MP8 conditions at the fixed
% frequency. Their nonzero Jacobian identifies it with the local classical
% return root, rather than assuming an arbitrary position return is periodic.
% MP8 uses a time mean, but at zeta=0 the radius is the constant r_0 and
% angular speed is constant. First-order changes of the time measure and
% period multiplying r_0 cancel against the normalization by that period;
% the surviving radius derivative is its uniform angle mean. Hence the
% time-mean derivative is also -57/4, proving the displayed limit. Under
% the uniform estimates just stated one may write
% Rcal_2/(omega^2 Rcal_0)=-57/4+O(sqrt(epsilon)).
% The native Pauli projection and centering belong to the preparation and
% enter these subleading terms. The angular-mixture counterexample above
% changes this ratio only by -eta^2 epsilon^2/2, consistently with the same
% leading limit. The scalar endpoint number -5/12 likewise occurs at that
% subleading epsilon^2 scale, where polarization and preparation effects
% cannot be inferred from this leading theorem. The absolute fixed-epsilon
% MP8 coefficient remains unevaluated; the radial [N] probe does not supply it.
% ACTUAL SAME-PAIR COMPARISON: after both independent saves, the complete
% MP15 argument in redteam_declared_exterior.tex was read. Its exact scaled
% native symbol, differentiated block-Egorov route, free-band leakage check,
% C4 return-root identification and normalized time mean agree with the
% derivation above. Its epsilon=(omega r_*)^-1 differs by a fixed factor,
% and its denominator r_* omega^2 has the same limit and error order since
% Rcal_0/r_*=1+O(sqrt(epsilon)). The common two-band tube and native
% embedding precisions from that comparison are explicit above. This is
% same-pair [R], not an other-lane [X] or a finite-frequency evaluation.
%
% FIXED-FREQUENCY ANGULAR REDUCTION [R, independent continuation]. This
% reduction was derived before exchanging the RED's next finite-frequency
% argument and saved before reading its actual new source. It retains the
% MP8 preparation, four return conditions, native centroid detector and
% frequency adjustment. It is a constructive reduction of the remaining
% absolute coefficient, not its numerical evaluation.
%
% The full continuum h_c is rotationally invariant for the fixed smooth
% radial core/exterior. The order-a^2 perturbation is exactly
% h_2=-(1/6) C_ijkl alpha_i p_j p_k p_l
%     -(1/3) C_ijkl {v_i,p_j p_k p_l},
% where C_ijkl=sum_{b=1}^3(e_b)_i(e_b)_j(e_b)_k(e_b)_l.
% Its unique rotation-scalar part is
% C^(0)_ijkl=(delta_ij delta_kl+delta_ik delta_jl
%                                      +delta_il delta_jk)/5.
% The remainder C^(4)=C-C^(0) is traceless and fully symmetric, hence
% a rank-four irreducible tensor. There is no rank-two part. Contracting
% C^(0) gives
% h_2^(0)=-(alpha.p)p^2/10-(1/5)sum_i{v_i,p_i p^2}.
% This agrees with the independently checked radial trace above, but here
% the complete rank-four part is RETAINED. In fixed lattice axes its
% spherical components have q=0,+4,-4. No angular average is performed
% on the MP8 state, Hamiltonian or detector.
%
% The native position is a rank-one tensor. Since conjugation by U_c
% preserves tensor rank, the dynamical centroid response
% i int Tr sigma_0 [h_2I(s),X_I(t)]ds sees only the rank-one part of
% sigma_0 from h_2^(0), and ranks 3,4,5 from h_2^(4). This follows by
% decomposing the tensor product 4 tensor 1, not by treating the packet
% as a rotational mixture. The preparation response needs separate care.
% P_2=(alpha.delta n)/2 is pure rank four: the scalar contractions of
% -p_i^3/(6|p|)+p_i sum_j p_j^4/(6|p|^3) cancel exactly.
% Before recentering, sigma_2={P_2,sigma_0}; its position response again
% uses ranks 3,4,5. With
% delta x_j=Tr[sigma_0{P_2,X_j}], exact native recentering adds
% i delta x_j[p_j,sigma_0]. Its contribution is
% delta x_j K_ij(t), K_ij(t)=i Tr[sigma_0[X_Ii(t),p_j]],
% and K_ij(0)=-delta_ij. The K tensor has ranks 0,1,2. Thus the COMPLETE
% centered coefficient uses initial density multipoles of ranks at most
% five, including the rank-zero/four energy row and the centering product;
% it is not restricted to ranks 1,3,4,5 after that state-dependent centering.
% For the reflection pair this formula is applied to each prepared member
% and the two expectations averaged. Its z centroid still vanishes exactly.
%
% Low tensor rank is not a cutoff on particle angular momentum. Write the
% native continuum spinor in the standard reduced radial basis
% psi(r,n)=r^-1 sum_{kappa,m}(F_{kappa m}(r) Omega_{kappa m}(n),
%                          i G_{kappa m}(r) Omega_{-kappa m}(n))^T,
% with j=|kappa|-1/2 and l_kappa=j+sgn(kappa)/2. On the regular full-Dirac
% domain its independent radial propagators have generators
% h_{c,kappa}=D_kappa+{v,p}/2,
% D_kappa=sigma_y p+kappa sigma_x/r, p=-i partial_r.
% The symbol p here is only a differential expression; no self-adjoint
% momentum on the half-line is assumed. The regular paired spinor boundary
% condition, the fixed core, and the actual outer domain define h_c.
%
% All angular blocks needed for the response can be constructed without
% numerically differentiating angular functions. Use spherical position
% X_q=r sqrt(4pi/3)Y_1q, q=0,+/-1. Let N_q be its angular matrix, diagonal
% in the two reduced radial components. Its upper block is the finite sum
% (N_q)_{kappa'm',kappa m}=sum_{s=+/-1/2}
%  CG(l',m'-s;1/2,s|j',m') CG(l,m-s;1/2,s|j,m)
%  *(-1)^(m'-s)*sqrt((2l'+1)(2l+1))
%  *(l' 1 l;0 0 0)*(l' 1 l;-(m'-s) q m-s).
% The parentheses are Wigner 3j symbols. In the lower block replace
% kappa',kappa by -kappa',-kappa. This is obtained directly by expanding
% both spinor harmonics in Y_lm times spin-1/2 and integrating their three
% scalar harmonics. It fixes normalization and phase with the stated
% Condon--Shortley CG convention; no shell degeneracy weight is inserted.
%
% Let L_kappa=diag(-partial_r^2+l_kappa(l_kappa+1)/r^2,
%                 -partial_r^2+l_{-kappa}(l_{-kappa}+1)/r^2).
% The full scalar p^2 has precisely these radial blocks. The exact momentum
% and Dirac-alpha blocks are therefore
% (P_q)_{ab}=(i/2)[L_a(rN_q)_{ab}-(rN_q)_{ab}L_b],
% (A_q)_{ab}=i[D_a(rN_q)_{ab}-(rN_q)_{ab}D_b].
% These follow from p_q=(i/2)[p^2,X_q] and alpha_q=i[h_0,X_q]. The
% apparent second derivatives cancel in P_q, and all radial derivatives
% cancel in A_q. Transform the q components to Cartesian components using
% X_{+1}=-(X_x+iX_y)/sqrt2, X_{-1}=(X_x-iX_y)/sqrt2, X_0=X_z.
% Set V_i=v(r)N_i. Then ordinary block multiplication gives every term of
% h_2=-sum_i A_i P_i^3/6-sum_i{V_i,P_i^3}/3, with all intermediate angular
% channels retained. This provides explicit finite differential blocks of
% order at most three, including endpoint ordering and all spin matrices.
% The spectral selection rules are |j'-j|<=4 and m'-m=0,+4,-4 for h_2,
% with the parity required by this even operator. The X blocks have
% |j'-j|<=1 and m'-m=q. These rules reduce work; they do not discard the
% rank-four response or replace the initial coherent angular profile.
%
% A finite angular approximation can be CONTROLLED independently of the
% radial propagation. Let P_J retain all j<=J, and first truncate the
% specified continuum preparation by P_J, with its normalization and
% parameter derivatives treated explicitly. Since [h_c,J^2]=0, unperturbed
% evolution never leaves that finite angular set. A single h_2 insertion
% or P_2 preparation insertion reaches only j<=J+4; the recentering momentum
% insertion reaches at most J+1. Consequently all order-a^2 propagation
% for that truncated state is EXACT within j<=J+4. There is no repeated
% anisotropy cascade at this derivative order. One must nevertheless retain
% the original state's discarded tail, rather than redefine its envelope.
%
% A useful tail norm is
% ||psi||_{3,1}=||(1+r^2)^(1/2)(1+p^2)^(3/2)psi||,
% with sufficient additional derivatives for parameter-differentiated
% estimates. Both factors are rotational scalars and commute with P_J.
% For any fixed integer s and smooth prepared state,
% ||(1-P_J)psi||_{3,1}
% <=[(J+1)(J+2)]^(-s/2)||[J^2]^(s/2)psi||_{3,1},
% when J is a retained half-integer (the first excluded angular momentum
% is J+1). The same estimate holds for each preparation parameter derivative.
% Fixed-time Sobolev propagation under the smooth bounded h_c, together
% with h_2:H^3->L^2 and the weighted position estimate from
% dot X=alpha+v, turns these tails into errors on F_0,F_2,J_0 and the
% mean-radius numerator. Constants depend on the actual core, time and
% derivative bounds. They cannot be replaced by the geometric-optics
% assertion that the packet starts narrow. In particular the actual
% continuum return root and its inverse-Jacobian bound remain necessary.
%
% The time insertion can also be integrated exactly in a radial spectral
% representation. For any two continuum energies E,E', the propagator
% derivative uses the divided difference
% D_t(E,E')=[exp(-iEt)-exp(-iE't)]/(E-E'),
% D_t(E,E)=-it exp(-iEt).
% Thus U_2(t)_{E a,E' b}=D_t(E,E') (h_2)_{E a,E' b}, with the actual
% radial spectral measure and the preparation terms added separately.
% This is equivalent to the complete matrix response, not a scalar-band
% approximation, and avoids subtracting nearby lattice runs. The radial
% spectral kernels, actual MP8 root, parameter derivatives and final
% centroid mean still require evaluation at the chosen finite frequency.
% No numerical value or execution request follows from the tensor reduction
% alone. SOURCES: direct isotropic tensor contraction, spinor-harmonic Gaunt
% integration, commutator construction of momentum/alpha blocks and spectral
% divided-difference integration above; the already checked MP3/MP4/MP8
% definitions supply the state and observable, without a new model input.
%
% EXPLICIT MP8 MEMBER, fixed before any finite-frequency result [R, WIP].
% This candidate was derived before reading the RED's actual new candidate
% source. It instantiates the previously symbolic core and envelope; it
% changes none of the native projection, centering, frequency or return
% rules. Its coefficient and actual return root are not yet evaluated.
%
% Fix r_s>0, r_*=3r_s/2, epsilon=2^(-24), and omega=1/(epsilon r_*).
% For b(t)=exp(-1/t) when t>0 and b(t)=0 otherwise, put
% eta(t)=b(t)/[b(t)+b(1-t)]. Thus eta=0 for t<=0 and eta=1 for t>=1.
% Let r_a=6r_s/5, r_b=13r_s/10, and
% f(r)=r_b^(-3/2)+eta((r-r_a)/(r_b-r_a))
%                              [r^(-3/2)-r_b^(-3/2)]
% for r>r_a, with f=r_b^(-3/2) for r<=r_a. Set N=1 and
% v(x)=-sqrt(r_s) f(|x|) x. This is a fixed C-infinity regular core and
% exactly the declared vacuum exterior for r>=r_b, including the photon
% circle r_*. It obeys the GLOBAL bound |v|<=sqrt(5/6)<1: in the join
% its magnitude is a convex combination of r sqrt(r_s)/r_b^(3/2)
% and sqrt(r_s/r), each bounded by sqrt(5/6), and the inner/outer pieces
% satisfy the same bound. This removes a core band-crossing surface from
% this particular prescribed probe experiment. It is not a calculation of
% a trapped object's interior or a microscopic source realization.
% In fact its continuum constraint mass is nonnegative and monotone:
% M(r)/E=[eta+(1-eta)(r/r_b)^(3/2)]^2 in the join, (r/r_b)^3 inside,
% and 1 outside. The derivative is nonnegative since eta'>=0 and r<=r_b.
% This source-profile check does not assert the missing fermion preparation.
%
% The envelope is fixed by a concrete real symplectic matrix, including a
% constant squeeze to reduce the one-circuit hyperbolic growth. Use scaled
% q=x/r_*, k=p/omega, time tau=t/r_*, and the rotating classical symbol
% e_rot(q,k)=|k|-sqrt(2/3)(q.k)/|q|^(3/2)
%                                  -(q_x k_y-q_y k_x)/sqrt3.
% Its fixed point is q_*=(1,0,0), k_*=(sqrt6,sqrt3,0), with
% tau_*=2pi sqrt3. Put K=J Hess(e_rot)|_*, J=[[0,I],[-I,0]].
% To make this matrix explicit, write K=[[B^T,C],[-A,-B]], where
% A=[[-3/2,3sqrt2/2,0],[3sqrt2/2,3,0],[0,0,3]],
% B=[[sqrt(2/3)/2,-1/sqrt3,0],
%    [1/sqrt3,-sqrt(2/3),0],[0,0,-sqrt(2/3)]],
% C=(1/9)[[1,-sqrt2,0],[-sqrt2,2,0],[0,0,3]].
% Define S=exp[-pi sqrt3 K] and G_*=S S^T. These are fixed numerical
% matrices, independent of a and of the variable launch labels R,p_r,L,T.
% Let V=(G_*)_kk and C_xk=(G_*)_qk. Symplecticity makes C_xk V^-1
% symmetric. For u in R^3 set
% g(u)=c theta(|u|/2048)
%       exp[-u^T V^-1 u/2-i u^T C_xk V^-1 u/2],
% theta(s)=b(2-s)/[b(2-s)+b(s-1)], and choose c>0 so int|g|^2du=1.
% This is an explicit smooth compact momentum envelope. Before its cutoff,
% its scaled phase-space covariance is (epsilon/2)G_*. The Fourier phase
% sign follows from q=i epsilon partial_k, so the symmetric q,k covariance
% is +(epsilon/2)C_xk. The cutoff is retained in the actual state and every
% evaluation; it is not silently replaced by an uncut Gaussian.
%
% For each launch z=(R,p_r,L,T), let p_c=(p_r,L/R,0). Use the standard
% Dirac representation, chi(n)=(sqrt((1+n_z)/2),
%                         (n_x+i n_y)/sqrt(2(1+n_z)))^T,
% w_+(n)=2^(-1/2)(chi(n),chi(n))^T, n=p_c/|p_c|.
% On the stated equatorial carrier patch this is a smooth unit positive-
% helicity, positive-free-energy spinor. Start from the momentum spinor
% g((p-p_c)/(omega sqrt(epsilon))) P_0(p)w_+(n), normalize it, and
% translate it to actual continuum centroid R e_1. Its native version is
% exactly the normalized P_a projection followed by the MP7 native-centroid
% translation, not a new spacing-dependent envelope deformation. Form the
% equal mixture with its S_z=(beta Sigma_z) spatial-z-reflected image.
% The incoming covariance is P_-^a plus this trace-one probe density;
% both it and P_-^a use the same subsequent full native h_a transport.
% The momentum support at the reference carrier lies in 2omega<=|p|<=4omega,
% since 4096 sqrt(epsilon)=1. A fixed neighborhood of the launch labels
% keeps it away from p=0; sufficiently small a omega keeps its one-valley
% embedding away from every Brillouin face. There is no extra valley weight.
%
% A check of the squeeze is possible by finite matrix algebra. In the
% horizontal (q_x,q_y,k_x,k_y) order, K_h^3=K_h/3 and its zero eigenspace
% is semisimple. Thus P_h=3K_h^2, P_0=I-P_h and
% S_h=P_0+cosh(pi)P_h-sqrt3 sinh(pi)K_h.
% The vertical block has K_z^2=-I/3 and S_z=-I (this matrix is unrelated
% to the spinor reflection symbol above). Explicitly ||P_h||_F=||P_0||_F
% =sqrt60 and ||K_h||_F=29/6, giving the simple bound ||S||<225.
% Its q_x row is
% (cosh(pi)-sinh(pi)/sqrt2,-sinh(pi),
%                       -sqrt3 sinh(pi)/9,+sqrt6 sinh(pi)/9).
% Under the QUADRATIC classical propagation the covariance is
% (epsilon/2)exp[K(tau-tau_*/2)]exp[K^T(tau-tau_*/2)].
% Hence the largest endpoint radial variance is
% (epsilon/2)[cosh(pi)^2+sqrt2 cosh(pi)sinh(pi)
%                                      +(29/18)sinh(pi)^2],
% and its rms is less than 17 sqrt(epsilon) r_* in physical units.
% This is an explicit localization estimate for the linearized Gaussian,
% not a proof of the nonlinear quantum return root at epsilon=2^(-24).
% Projection/cutoff, nonlinear field derivatives, native interband terms,
% parameter derivatives and the actual inverse return Jacobian still need
% quantitative control. The carrier label alone does not determine the
% sparse angular cost: the angular fluctuation and radial spectral accuracy
% must instead be bounded, as specified below. No implementation or execution is authorized by this
% candidate definition. Its purpose is to fix every datum needed for a
% controlled evaluation of the SAME MP8 observable before quoting a value.
% SOURCES: direct radial-core monotonicity and bound, rotating-symbol
% Hessian, symplectic Gaussian covariance and exact native preparation
% formulas above. Existing MP8 supplies the four fixed measurement rules.
%
% PAIRED MEMBER SELECTION AND ANGULAR CONTROL [R, WIP]. The explicit author
% member immediately above is selected unchanged, before any finite-frequency
% output: epsilon=2^-24, its globally subcritical C-infinity core, full
% half-period symplectic squeeze, exact 2048--4096 smooth cutoff and stated
% projected/centered/reflection-paired spinors. The independently saved RED
% MP17 alternative remains a record, not another input to this calculation.
% Actual MP17 was read completely after the author candidate was saved and
% self-read. Its canonical squeeze and quadratic angular-momentum warning
% are useful checks; no linearized-width estimate is promoted to a nonlinear
% root certificate. The exact target remains the MP8 coefficient at this
% finite epsilon, with normalized error target 0.05 for Rcal_2/(r_* omega^2).
% An actual root/inverse-Jacobian certificate is part of that target. These
% are accuracy objectives, not achieved bounds or an execution specification.
%
% An honest angular count can start with an exact Gaussian calculation.
% For the uncut scalar control at the classical launch, the linear part of
% physical L_z has row l=(sqrt3,-sqrt6,0,1) on
% (delta q_x,delta q_y,delta k_x,delta k_y), divided by epsilon.
% Since l K_h=0 and l S_h=l, its variance is exactly 5/epsilon. Thus its
% rms is sqrt5 times 4096, approximately 9159, not the huge carrier label
% sqrt3/epsilon. The exact L_z also contains the quadratic Cartesian
% fluctuation; native centering, the projector and the cutoff must all be
% included before assigning an actual retained j,m set. In particular a
% narrow LINEAR radial width does not establish a narrow angular spectrum.
%
% The complete uncut angular tail has a closed rotation-moment formula.
% Write M=V^-1, B=C_xk V^-1, A=M+iB, and use dimensionless k. The scalar
% state is proportional to
% exp[-(k-k_c)^T A(k-k_c)/(2epsilon)-i q_c.(k-k_c)/epsilon].
% Here q_c is its translation parameter; for the actual preparation it is
% the translation found by the stated native/continuum centroid rule, not
% automatically the requested measured centroid. Let O_t=R_z(-it) and
% D_t=A^*+O_t^T A O_t,
% b_t=A^*k_c+O_t^T A k_c+i(q_c-O_t^T q_c).
% Wherever Re D_t>0, with the square-root branch continuous from t=0,
% Z(t)=<exp(t L_z)>
%     =sqrt[det(2M)/det D_t]
%       exp{[b_t^T D_t^-1 b_t-2 k_c^T M k_c]/(2epsilon)}.
% This follows by the exact Gaussian integral for rotation of the momentum
% wavefunction and analytic continuation; it includes ALL quadratic angular
% fluctuations. For the fixed equatorial carrier spinor the corresponding
% unprojected J_z moment is Z(t) cosh(t/2). The Chernoff bounds are therefore
% Prob(J_z>=m_+)<=inf_{t>0} exp(-t m_+) Z(t)cosh(t/2),
% Prob(J_z<=m_-)<=inf_{t<0} exp(-t m_-) Z(t)cosh(t/2),
% with each infimum restricted to its actual Gaussian convergence domain.
% They may be evaluated from finite matrices; replacing Z by its linear
% variance would discard the non-Gaussian angular tail.
%
% Projection does not invalidate this angular comparison: P_0 commutes
% with J_z and J^2. Thus any angular spectral tail of P_0 psi is bounded
% by the corresponding unprojected tail, divided by the squared projected
% norm. At the reference carrier the actual compact momentum support has
% |k-k_c|<=1, |k_c|=3, so that norm is at least
% [1+sqrt(1-1/9)]/2=[1+2sqrt2/3]/2. The statement concerns squared norm.
% A launch neighborhood |k_c-k_*|<=1/4 similarly gives the lower bound
% [1+sqrt(1-16/121)]/2. The exact reflection mixture has the same angular
% tail bounds. No artificial angular averaging or loss of coherences occurs.
% This commutation statement concerns the continuum P_0 and its centered
% state. The finite-a P_a does not commute with J_z or J^2; its rank-four
% P_2 and the native-centering response remain in the spacing derivative.
%
% The actual smooth cutoff need not have an exponential angular moment;
% the uncut Gaussian is used only as a norm comparison. Since ||S||<225,
% the omitted scalar Gaussian probability is at most
% tau=erfc(2048/225)+(4096/(225 sqrt(pi)))exp[-(2048/225)^2].
% Normalizing the cut scalar state changes its L2 norm by at most 2sqrt(tau).
% Angular-tail norms can therefore be compared by the triangle inequality,
% before the projected-norm lower bound. For weighted and parameter-derivative
% tails the corresponding finite derivatives of the ACTUAL theta and
% Gaussian must be bounded separately; L2 agreement alone is insufficient
% for the h_2 insertion, position detector or return Jacobian. Likewise a
% J_z tail is not by itself a J^2 tail. These outstanding controls prevent
% assigning an unsupported retained dimension or a short-run resource cap.
% SOURCES: direct metaplectic Gaussian integral and angular commutation of
% the exact continuum free projector; actual RED MP17/full SOURCES and its
% independent check of l S_h=l. This new finite-member calculation is [R],
% not evidence supplied by the other lane or a new numerical result.
%
% FULL-MATRIX RESIDUAL CERTIFICATE FOR THE SELECTED MEMBER [R, WIP].
% This is an independent constructive control route, not an evaluated
% Gaussian approximation or a numerical execution specification. It keeps
% the selected finite epsilon, actual preparation and MP8 detector fixed.
% Write q=x/r_*, tau=t/r_*, k=-i epsilon partial_q and delta=(a omega)^2.
% The continuum and spacing-derivative operators in units of omega are
% H=alpha.k+{v,k}/2,
% B=partial_delta H_delta|_0
%   =-sum_i alpha_i k_i^3/6-sum_i {v_i,k_i^3}/3.
% Here {v,k}/2 denotes the SUM over components. In particular B includes
% the endpoint ordering. Its Weyl symbol contains the scalar term
% +(epsilon^2/2)sum_i (partial_i^2 v_i)k_i; it is not discarded.
% If u is the actual continuum spinor and w=partial_delta u|_0, then
% i epsilon dot u=H u,  i epsilon dot w=H w+B u,
% with the exact free-projector, normalization and centering derivatives
% in w(0). Delta is only a first-order response variable here; no global
% self-adjointness claim for a cubic truncated H+delta B is required.
%
% The chosen core gives a useful GLOBAL energy-graph control. Put
% v_max=sqrt(5/6), d=1-v_max>0 and b=div(v)/2. On Schwartz spinors,
% d ||k psi|| <= ||H psi||+epsilon ||b||_infinity ||psi||.
% This follows from ||alpha.k psi||=||k psi|| and the pointwise bound
% ||v.k psi||<=v_max||k psi||. The symmetric drift is relatively bounded
% with relative bound below one, so H is self-adjoint on H^1. Smooth bounded
% core derivatives give the higher graph domains. Explicit constants can
% be generated without a time-exponential Sobolev estimate: set
% P_s(psi)=sum_{|alpha|<=s}||k^alpha psi||,
% E_s(psi)=sum_{j=0}^s||H^j psi||, C_0=1, and
% D_s=sum_{|alpha|=s}sum_{0<beta<=alpha} binom(alpha,beta)
%       [epsilon^|beta| sum_i||partial^beta v_i||_infinity
%        +epsilon^(|beta|+1)||partial^beta b||_infinity],
% C_{s+1}=C_s[1+(sqrt3/d)(1+epsilon||b||_infinity+D_s)].
% Then P_s<=C_s E_s, by applying the first estimate to each k^alpha psi
% and commuting H through k^alpha. Only s<=3 is needed for B. These are
% conservative explicit finite constants; their numerical values and all
% required core-derivative suprema remain to be bounded before execution.
%
% For example, expanding the ordered B exactly gives
% B=-sum_i(alpha_i/6+2v_i/3)k_i^3
%       +i epsilon sum_i v_i' k_i^2
%       +epsilon^2 sum_i v_i'' k_i
%       -i epsilon^3 sum_i v_i'''/3,
% where primes mean derivatives in coordinate i. Thus
% ||B psi||<=c_B P_3(psi)<=c_B C_3 E_3(psi),
% c_B=max(1/6+2v_max/3, epsilon V_1, epsilon^2 V_2,
%                                      epsilon^3 V_3),
% V_j=max_i||partial_i^j v_i||_infinity.
% The zero-derivative term in P_3 occurs once; the three final B terms
% sum to at most epsilon^3 V_3||psi||, as used in c_B.
% Also [q_i,B]=i epsilon W_i with
% W_i=-alpha_i k_i^2/2-{v_i,k_i^2}. Consequently a weighted B error is
% controlled by E_3(q_i psi) plus a lower-order E_2(psi) term. No radial
% reduction, spin suppression or angular incoherent sum enters this bound.
%
% Let tilde u be ANY differentiable full four-spinor approximation in the
% stated graph domains, r=i epsilon dot(tilde u)-H tilde u, and e=u-tilde u.
% Unitarity and [H,H^j]=0 give the exact a posteriori majorants
% e_j(tau):=||H^j e(tau)||
% <=||H^j e(0)||+epsilon^-1 int_0^tau||H^j r(s)||ds  (0<=j<=3).
% For weighted errors it is convenient to put q BEFORE the graph powers:
% e_{ij}^q(tau):=||H^j q_i e(tau)||
% <=||H^j q_i e(0)||
%   +int_0^tau||H^j V_i e(s)||ds
%   +epsilon^-1 int_0^tau||H^j q_i r(s)||ds,
% V_i=alpha_i+v_i. The middle term has an explicit finite differential-
% coefficient bound times E_j(e); H^j V_i is only order j. Its coefficients
% are obtained recursively from H=sum_l(alpha_l+v_l)k_l-i epsilon b.
% This avoids an unquantified propagation constant for the exact native
% coordinate observable. The inequalities apply to each reflection member.
%
% For a candidate tilde w define
% r_1=i epsilon dot(tilde w)-H tilde w-B tilde u.
% Then e_w=w-tilde w obeys i epsilon dot e_w=H e_w+B e-r_1. Hence
% ||e_w(tau)|| <= ||e_w(0)||
%  +epsilon^-1 int_0^tau[||B e(s)||+||r_1(s)||]ds,
% ||q_i e_w(tau)|| <= ||q_i e_w(0)||
%  +int_0^tau||V_i e_w(s)||ds
%  +epsilon^-1 int_0^tau[||q_i B e(s)||+||q_i r_1(s)||]ds.
% All terms on the right are controlled by the preceding base-error graph
% majorants and ACTUAL residual norms. In particular, an L2 residual alone
% is insufficient because B has order three. The scaled centroid derivative
% x_delta=2 Re<u,q w> has the rigorous component error
% |x_delta,i-2 Re<tilde u,q_i tilde w>|
% <=2[||e|| ||q_i tilde w||+||q_i e_w||], since ||u||=1.
% The baseline position error is at most
% ||e|| ||q_i tilde u||+||q_i e||.
% Carrier/center derivatives have the same certificates with the exact
% differentiated initial data; H itself is fixed when those labels vary.
% Endpoint-time derivatives use the bounded V_i, not a falsely slow full
% spinor phase. The fixed-frequency row is the graph-controlled expectation
% of H and its B/preparation derivative, not a fixed-carrier substitute.
%
% A moving Gaussian construction can supply tilde u and tilde w while the
% preceding inequalities, not a presumed Gaussian shape, decide acceptance.
% One explicit algebraic route uses TWO complex phases S_sigma, sigma=+,-,
% and polynomial spinor amplitudes. Their real central characteristics use
% e_sigma(q,k)=v.k+sigma sqrt(k.k). The selected core keeps both band
% characteristics away from k=0 at their nonzero energies: along either,
% |e_sigma|/(1+v_max)<=|k|<=|e_sigma|/(1-v_max).
% The lower component is retained even though its initial amplitude starts
% at higher semiclassical order. Introduce formal delta in the PRINCIPAL
% matrix alpha.(k-delta k^3/6)+v.k-(2delta/3)sum_i v_i k_i^3 and in its
% two analytic eigenvalues/projectors on the controlled complex momentum
% tubes. No scalar band is substituted for the actual propagated spinor.
%
% The phase/amplitude recursion is explicit from conjugation. Define
% K_i[S]=partial_i S-i epsilon partial_i. EXACTLY,
% exp(-iS/epsilon)(H+delta B)exp(iS/epsilon)
% =sum_i alpha_i K_i+sum_i{v_i,K_i}/2
%    +delta[-sum_i alpha_i K_i^3/6-sum_i{v_i,K_i^3}/3].
% Expand this finite differential expression as h_pr(q,partial S)
% +sum_{j=1}^3 epsilon^j L_j[S,delta]. With
% partial_tau S_sigma=-e_sigma,delta(q,partial S_sigma),
% tilde u=sum_sigma exp(iS_sigma/epsilon)sum_{n=0}^K epsilon^n a_sigma,n,
% its coefficient equations are
% (e_sigma,delta-h_pr)a_sigma,n
%   =sum_{j=1}^{min(3,n)} L_j a_sigma,n-j
%                            -i partial_tau a_sigma,n-1,
% with a_sigma,-1=0. Projecting onto the band gives its transport equation;
% the complementary component is divided by the actual nonzero two-band
% gap. Taylor jets about each real central trajectory turn these equations
% into finite ODEs once phase/amplitude degrees are fixed. Coordinate
% cutoffs, positivity of the retained complex Gaussian phase and the
% omitted Taylor remainder belong to the actual residual norm, not to an
% unexamined asymptotic-order claim. Differentiating this construction in
% formal delta produces tilde w with the FULL ordered B above.
%
% The exact selected initial data are also retained. For the uncut scalar
% Gaussian with A=M+i B_phase, B_phase=C_xk V^-1, its Fourier transform has
% S_0(q)=k_c.(q-q_c)+(i/2)(q-q_c)^T A^-1(q-q_c).
% Acting with the free projector has the Gaussian stationary-phase series
% [exp{(epsilon/2)A^-1:partial_k partial_k}P_0(k)w_+(n)]
%     evaluated at k=k_c+i A^-1(q-q_c).
% This is only a method of constructing the initial polynomial amplitudes:
% their error is tested against the actual Fourier integral, including the
% fixed cutoff and normalization, rather than declaring the series exact.
% The global Fourier/translation phase is matched as well when certifying
% state or preparation-derivative errors, even though it cancels in the
% final density observables.
% At first order in delta the momentum multiplier is P_delta(k)P_0(k),
% followed by the stated centering translation; substituting P_delta alone
% on the unprojected scalar envelope would change the fixed preparation.
% The two initial phase branches and their parallel amplitudes must sum to
% these actual initial jets, including the higher-order lower-band part.
%
% The actual return root can be certified independently of the classical
% existence statement. Use scaled z=(R/r_*,p_r/omega,L/(omega r_*),T/r_*).
% Suppose the computation supplies a differentiable approximate return map
% tilde F on a ball B_rho(z_c), certified errors eta_F and eta_J on F and
% D F, and an explicit matrix A_c. If
% sup_B||I-A_c D tilde F||+||A_c||eta_J<=q<1,
% ||A_c tilde F(z_c)||+||A_c||eta_F<=(1-q)rho,
% then z->z-A_c F(z) is a contraction of that ball. There is a unique
% ACTUAL MP8 root there and ||(D F)^-1||<=||A_c||/(1-q). Winding, positive
% final x and exclusion of T=0 must also be certified on the centroid tube.
% The coefficient is evaluated with this root, the full F_delta, mean-
% radius derivative and time/energy rows. Validated quadrature and all
% state, response, preparation, root and interpolation errors must sum to
% at most 0.05 in the stated normalized coefficient. No phase/amplitude
% degree, retained dimension, runtime or successful certificate is asserted
% here; those require the quantitative bounds for the actual finite member.
%
% SOURCES: direct relative-bound/commutator, unitary Duhamel, conjugated
% differential-operator and contraction arguments above. The actual primary
% Fermanian-Kammerer--Lasser--Robert article, DOI10.1007/s00220-021-04147-2,
% https://link.springer.com/article/10.1007/s00220-021-04147-2,
% Sections2.1--2.2 and3.1 (Proposition2.3, Theorem3.2, Corollary3.4) was
% read for comparison. Its estimates require global subquadratic/gap
% hypotheses and unspecified constants. Our bare v(q).k symbol is not
% globally subquadratic in that sense, and the native massless spectrum
% has no global gap. Those theorems do not supply this finite-member error
% certificate, its dimensions or its root. The displayed direct residual
% bounds use the original fixed-core self-adjoint H instead.
%
% EXPLICIT CORE AND GAUSSIAN CONSTANTS FOR THAT CERTIFICATE [R, WIP].
% These bounds use the selected member, without changing its envelope or
% epsilon. In dimensionless q, its join radii are 4/5 and 13/15 and its
% join width is 1/15. The cutoff function eta obeys the global bounds
% |eta'|<=8, |eta''|<=128, |eta'''|<=4096. One direct proof uses symmetry
% about 1/2 and, for 0<t<=1/2, x=1/t>=2 and
% eta=1/(1+exp(h)), h=1/t-1/(1-t). Successive chain rules give bounds by
% e^2 exp(-x) times, respectively,
% x^2+4,
% x^4+2x^3+8x^2+32,
% x^6+6x^5+18x^4+24x^3+96x^2+352.
% Bounding each x^n exp(-x) by its maximum on x>=2 proves the stated
% convenient constants. Flat endpoints contribute no larger derivatives.
%
% Write v_i=-sqrt(2/3) f(r)q_i. On the join,
% h(r)=r^(-3/2)-(13/15)^(-3/2) has
% |h|<1/4, |h'|<3, |h''|<9, |h'''|<40, and |f|<3/2.
% Therefore |f'|<=33, |f''|<=7929 and |f'''|<=3718480.
% Cartesian differentiation of f(r)q_i yields, uniformly over every
% component/multi-index and over the entire core/exterior,
% M_1:=max_{i,|alpha|=1}||partial^alpha v_i||_infinity<=35,
% M_2<=8100, M_3<=3800000.
% For example the join bounds are |f|+r|f'|, r|f''|+3|f'|,
% and r|f'''|+9|f''|+9|f'|/r, before multiplication by sqrt(2/3).
% The constant inner f and the decaying exterior satisfy the same bounds.
% Thus ||b||<=52.5, ||partial b||<=12150 and
% ||partial^2 b||<=5700000 for b=div(v)/2.
%
% At epsilon=2^-24<6*10^-8 these give
% D_0=0,
% D_1<=315epsilon+36450epsilon^2<1.9*10^-5,
% D_2<=1260epsilon+291600epsilon^2+34200000epsilon^3<7.6*10^-5.
% Since sqrt3/(1-sqrt(5/6))<20, the preceding recurrence gives
% C_1<21.001, C_2<442, C_3<9300.
% In particular the explicit source bounds may be taken as
% ||B psi||<=7220 E_3(psi),
% ||W_i psi||<=1030 E_2(psi),
% ||q_i B psi||<=7220 E_3(q_i psi)+1030epsilon E_2(psi).
% These deliberately conservative constants certify a residual; they are
% not estimates of the physical coefficient or a claim that any candidate
% finite basis has achieved the required residual.
%
% The Gaussian tail can also be made considerably sharper than the earlier
% valid 225 bound. Let c=cosh(pi), s=sinh(pi). For the stated horizontal
% matrices, direct traces give
% ||P_h||_F^2=||I-P_h||_F^2=60,
% tr[(I-P_h)^T P_h]=-58,
% tr(P_h^T K_h)=tr[(I-P_h)^T K_h]=0,
% ||K_h||_F^2=841/36.
% Consequently, including the two-dimensional vertical -I block,
% ||S||_F^2=122-116c+(1561/12)s^2<128^2.
% Thus ||V^(1/2)||<128, and the probability omitted by the exact scalar
% cutoff is bounded more sharply by
% tau_128=erfc(16)+(32/sqrt(pi))exp(-256).
% This improves a bound, not the chosen cutoff. In particular its L2
% normalization comparison remains 2sqrt(tau_128).
%
% Polynomially weighted tail integrals are explicit too. For the normalized
% uncut scalar g and integers 0<=m<=10,
% int_{|u|>=2048}|u|^(2m)|g(u)|^2du
% <=128^(2m)(2/sqrt(pi))Gamma(m+3/2,256)
% <=20*2048^(2m) exp(-256).
% The last inequality follows by writing the incomplete-gamma integral
% from 256 to infinity and using
% (1+t/256)^(m+1/2)<=exp[(m+1/2)t/256].
% Purity also gives ||M||<=128^2, ||B_phase||<=128^2 and hence
% ||A||<=32768; indeed G_qq=V^-1+B_phase V B_phase and
% ||G||<128^2. Gaussian derivatives are polynomials generated by
% partial_u g=-A u g. The radial theta cutoff has derivative bounds
% |partial_u theta|<=8/2048,
% |partial_u^2 theta|<=136/2048^2,
% |partial_u^3 theta|<=5000/2048^3
% on its transition shell, componentwise. These formulas turn all required
% finite weighted/preparation-derivative cutoff errors into finite Gaussian
% moments. They do not license using the bare uncut P_0 at k=0: higher
% derivatives of that projector are singular there. For such a comparison
% one must use a specified smooth extension equal to P_0 on the actual
% compact support, or bound the cut state directly. The actual preparation
% itself has no such singularity, since its support stays away from zero.
% SOURCES: direct one-variable chain-rule bounds, Cartesian radial
% differentiation, the graph recurrence, exact finite-matrix traces and
% incomplete-gamma integration above. No code or numerical propagation was
% used to obtain these constants. They still do not determine a sufficient
% phase/amplitude degree, actual return root or achieved coefficient error.
%
% FINITE INPUT AND DUAL-RESIDUAL CONTROLS [R, WIP].
% Actual RED MP19--MP21, including their complete SOURCES, were read after
% the preceding author residual/constants derivation was saved. The
% J_+-ladder bound is an additional exact J^2 control: on joint sectors
% m>=m_min and j-m>=n_c, for integers 1<=s<=n_c,
% ||1_tail psi||^2 <= ||J_+^s psi||^2 /
%       [(n_c)_falling_s (2m_min+n_c+1)_rising_s].
% This follows by multiplying the exact raising factors
% (j-m-l)(j+m+l+1), l=0,...,s-1. It does not replace the separate m-tail
% bound, and P_0 rather than the finite-spacing P_a commutes with J_+.
% The same actual RED source derives the first adapted-band correction
% Pi_1=Sigma.(n cross F)/(4|k|^2), F_j=(partial_j v).k.
% At the selected carrier it is -Sigma_z/(12sqrt3). Thus the declared free
% preparation has a genuine higher-order opposite-band part; differentiating
% its relative fast phase can lower a remainder order by one power of
% epsilon. Both phases in the preceding construction are necessary.
%
% The RED's finite initial-projector bounds can be checked directly without
% an angular discretization. Let kappa=|k_c|, z=k-k_c, r=|z| and
% t=(2k_c.z+r^2)/kappa^2. In the fixed launch neighborhood,
% 11/4<=kappa<=13/4 and on the actual cutoff support r<=1, so
% |t|<=(104/121)r. With c_j=binom(-1/2,j), put
% A_N=1/2[I+(alpha.k/kappa)sum_{j=0}^N c_j t^j].
% For every N>=0 the actual cut-Gaussian input obeys
% ||(P_0-A_N)g_cut w||
% <=5[(13/484)sqrt((N+3)/2)]^(N+1)/sqrt(1-tau_128).
% Indeed |c_j|<=1 bounds the pointwise remainder by
% 5[(104/121)r]^(N+1), while |z| moments are bounded by those of
% |xi|/32, with density pi^(-3/2)exp(-|xi|^2); the product
% E|xi|^(2m)=product_{j=1}^m(j+1/2) is at most [(m+2)/2]^m.
% N=39 gives degree79 and error <3*10^-36 before positive-band
% normalization. This is an approximation to the chosen state, not a
% replacement of its exact projector or cutoff.
%
% Direct Clifford multiplication also confirms the required ordered input
% response, in delta=(a omega)^2 units,
% P_2 P_0=-alpha.k^(3)/(24p)+alpha.k S_4/(24p^3)
%                         -i Sigma.(k^(3) cross k)/(24p^2),
% p=|k|, S_4=sum_i k_i^4.
% Expanding its three inverse powers in t through degree39 gives degree83.
% For t_max=104/121, the p^-3 remainder factor is
% |t|^40[81/(1-t_max)+2t_max/(1-t_max)^2]. Its prefactor is bounded using
% the CORRELATED maximum (kappa+1)^5/(24kappa^3) at kappa=13/4.
% Adding the other two terms gives <1200|t|^40 and actual cut-Gaussian
% L2 error <10^-33. Separate extrema for numerator and denominator would
% lose this useful constant. Native-centering and launch derivatives still
% have to be added and controlled; these L2 bounds do not certify them.
% The full scalar monomial counts are binom(82,3)=88560 and
% binom(86,3)=102340, respectively. They are input upper bounds, not a
% chosen propagated dimension; N can be reduced if the complete error
% budget permits. The common scalar Gaussian and factor t also preserve
% algebraic sparsity that a full monomial count does not exploit.
%
% There is a useful exact goal-oriented improvement of a crude forward
% response norm estimate. Set Y=(u,w), Mcal=[[H,0],[B,H]], let
% Y_tilde be an approximation, R=(r,r_1)=(i epsilon partial_t-Mcal)Y_tilde,
% and E=Y-Y_tilde=(e,f). For a coordinate Q=q_i the response functional is
% 2 Re<u,Qw>. Its linearized error has terminal dual data
% Z_T=(Q w_tilde(T),Q u_tilde(T)); the remaining exact term is
% 2 Re<e(T),Q f(T)>. For ANY approximate backward dual Z_tilde, define
% S=(i epsilon partial_t-Mcal^dagger)Z_tilde. Differentiating the pairings
% gives the exact identity
% response_error = 2 Re<Z_tilde(0),E(0)>
%   +2 Re<Z_T-Z_tilde(T),E(T)>
%   -(2/epsilon) Im int_0^T[<Z_tilde,R>+<S,E>]dt
%   +2 Re<e(T),Q f(T)>.
% This also verifies the sign in RED MP20's exact-dual formula. The last
% term is not omitted; its modulus is at most 2||Qe(T)|| ||f(T)||.
% Integral observables and half-time return rows insert the corresponding
% dual source and jump, with the same identity on each time interval.
%
% In particular let W(t):C^n->L2(C^4) be a differentiable isometry whose
% columns lie in the required graph/coordinate domains, P=WW^dagger,
% K=W^dagger H W-i epsilon W^dagger dot W, and B_n=W^dagger B W.
% K and B_n are Hermitian. Galerkin primal coefficients obey
% i epsilon dot c_u=K c_u,
% i epsilon dot c_w=K c_w+B_n c_u;
% backward dual coefficients obey
% i epsilon dot c_zu=K c_zu+B_n c_zw,
% i epsilon dot c_zw=K c_zw.
% Both residual components are orthogonal to Ran W. If the primal initial
% data are their exact P(0) projections, without renormalizing that finite
% approximation, then E(0) is also orthogonal to this range. Choosing dual
% terminal data P(T)Z_T therefore cancels exactly
% the first and the <Z_tilde,R> terms above. The residual certificate then
% involves terminal projection defects and products of primal errors with
% dual residuals, as well as the retained quadratic term. This cancellation
% can substantially improve on bounding every forward response component
% in norm; it is an algebraic identity, not a claim of an achieved error.
% No finite L2 space is invariant under the unbounded coordinate operator,
% so the terminal Q-projection defect still requires its actual bound.
% Renormalizing the finite projection, approximate input quadrature or
% coefficient integration spoils exact
% orthogonality by a measurable defect, which belongs in the displayed
% identity. The energy row, launch derivatives and actual root certificate
% remain part of MP8 and cannot be dropped by this reformulation.
%
% A still simpler refinement, independently given in RED MP20, avoids a
% weighted RESPONSE error for the coordinate readout. If E_0>=||e||,
% Z_0>=||Qe|| and D_0>=||f||, then
% |x-x_tilde|<=2||Q u_tilde||E_0+E_0 Z_0,
% |x_delta-x_delta_tilde|<=2||Q w_tilde||E_0
%                         +2(||Q u_tilde||+Z_0)D_0.
% These follow by expanding the two exact quadratic forms. The graph-three
% base residual needed by B is still required. Neither refinement supplies
% a phase degree, finite residual value, return root or numerical result.
% Actual RED MP23/SOURCES additionally gives a concrete terminal-defect
% formula, checked by the creation-operator identity. In a full-spinor
% Gaussian/Hermite frame of total degree <=N, if
% q_i=q_ci+sqrt(epsilon/2)sum_j(Qcal_ij a_j^dagger+
%                                      conjugate(Qcal_ij)a_j),
% then ||(1-P_N)q_i phi|| is at most
% sqrt[epsilon(N+1)/2] ||Qcal_i|| ||P_=N phi||.
% Only creation from the actual top shell leaves the retained space.
% Its amplitudes must be evaluated; they are not presumed zero. This form
% uses a fixed four-component spin frame. A momentum-dependent spin frame
% contributes its coordinate connection and cannot omit that term.
% SOURCES: actual RED MP19--MP23/SOURCES, checked by the independent
% Gaussian moments, Clifford product and dual-pairing derivations above.
% The moving-isometry Galerkin cancellation is derived here explicitly.
% All comparisons are within the same author/RED pair and remain [R].
%
% ISOTROPIC GRAPH IMPROVEMENT FOR THE SAME CERTIFICATE [R, WIP].
% A tensor l2 derivative norm avoids the earlier conservative sum over
% multi-indices. Define T_s(psi)=|| |k|^s psi||. Equivalently T_s is the
% l2 sum of all 3^s ORDERED derivative tensors k_{i_1}...k_{i_s}psi.
% Applying d||k f||<=||Hf||+epsilon b_0||f|| to every such component and
% using Minkowski's inequality gives
% d T_{s+1}(psi)<=T_s(Hpsi)+epsilon b_0 T_s(psi)
% +sum_{j=1}^s binom(s,j)[epsilon^j V_j^F T_{s-j+1}(psi)
%                         +epsilon^(j+1) b_j^F T_{s-j}(psi)],
% where V_j^F=sup_q||D^j v(q)||_F<=3^((j+1)/2)M_j and
% b_j^F<=3^(j/2) max_{|alpha|=j}||partial^alpha b||_infinity.
% All coefficients are on the left of the remaining derivatives in this
% commutator expansion. The binomial counts refer to the ordered tensor
% derivative positions, so no hidden multi-index multiplicity is omitted.
%
% Put K_0=1 and, using the monotonicity of the following constants, set
% K_{s+1}=K_s/d [1+epsilon b_0+
%       sum_{j=1}^s binom(s,j)(epsilon^j V_j^F+epsilon^(j+1)b_j^F)].
% Induction gives T_s<=K_s E_s. With the previously proved core constants,
% V_1^F<=105, V_2^F<=24300sqrt3,
% b_1^F<=12150sqrt3, b_2^F<=17100000 and d^-1=6+sqrt30,
% the selected epsilon yields K_1<11.478, K_2<131.75, K_3<1513.
% The ordered B can now be bounded without an extra sum over components:
% ||Bpsi|| <=(1/6+2v_max/3)T_3
%             +sqrt3 epsilon M_1 T_2
%             +sqrt3 epsilon^2 M_2 T_1+epsilon^3 M_3 T_0
%          <1180 E_3(psi).
% Here Clifford anticommutation gives
% ||sum_i alpha_i k_i^3 psi||^2=sum_i||k_i^3 psi||^2<=T_3^2;
% the scalar cubic drift uses the pointwise vector bound |v|<=v_max.
% Similarly, expanding {v_i,k_i^2} exactly gives
% ||W_i psi|| <=(1/2+2v_max)T_2+2epsilon M_1 T_1
%                                      +epsilon^2 M_2 T_0
%             <307 E_2(psi).
% Thus the same global response certificate may use 1180 and 307 in place
% of 7220 and 1030. No locality, scalar-band replacement or omitted spin
% term is used in this improvement. It still bounds an actual residual,
% rather than asserting one has been achieved by a finite approximation.
% SOURCES: independent ordered-tensor commutator, Minkowski and Clifford
% calculations above, using only the already checked selected-core bounds.
%
% NATIVE INPUT CENTERING AND A PHASE-STABLE RESPONSE REPRESENTATION [R, WIP].
% Actual RED MP22/SOURCES were read completely. Its differentiated binomial
% and Gaussian-moment bounds agree: with n=N+1 and
% T_n=(13/484)^n sqrt[Gamma(n+3/2)/Gamma(3/2)], the pretranslation
% one-coordinate input error is bounded by
% [epsilon(5n+36)T_{n-1}
%       +5||A||sqrt(n+3/2)T_n/32]/sqrt(1-tau_128),
% plus the explicit cutoff-derivative tail. At N=39 it is <3*10^-32.
% The differentiated normalization and actual centering are still included;
% an arbitrary phase convention cannot be inserted into a derivative test.
%
% In particular, for n_c=k_c/|k_c|, z=k-k_c, p=|k|,
% s=(1+n_c.k/p)/2 and Z=E_g s, the ACTUAL cut projected continuum mean is
% mu=Z^-1[B_phase E_g(zs)-(epsilon/4)E_g((n_c cross z)/p^2)].
% Real-amplitude derivatives cancel from this real expectation. The stated
% moment/normalization bounds give |mu|<0.05 and ||D_{k_c}mu||<0.34;
% hence q_c=R e_1-mu is a specified centering parameter, not simply R e_1.
% The sign of the Berry term follows from
% (i epsilon/2)[P_0,partial_i P_0]
%     =-(epsilon/4)Sigma.(n_k cross partial_i n_k).
%
% For normalized phi with P_0 phi=phi, the exact spacing centering uses
% M_i=-i[partial_i P_0,P_2]
%    =(1/2)Sigma.(partial_i n_k cross delta n_k),
% delta n_k=-k^(3)/(6p)+k S_4/(6p^3),
% partial_delta<q_i>|_0=epsilon<phi,M_i phi>.
% Therefore the prescribed centered initial response is
% w(0)=P_2 phi+i sum_i<M_i> k_i phi.
% This follows directly from (1-P_0)q_i phi=i epsilon(partial_i P_0)phi.
% For the equatorial positive-chirality member, only <M_z> can be nonzero
% and <M>=E_phi[(k cross k^(3))/(12p^3)]. Its central value
% -sqrt2/36 is not substituted for the finite-width average. The exact
% vertical variance bounds the centering contribution by 7*10^-5 in norm.
% The inequality |k cross k^(3)|<=p^4/2 follows because its squared ratio
% to p^8 is the variance of n_i^2 under weights n_i^2, each in [0,1].
% It also verifies the RED's stated 700 binomial remainder for M_i:
% the two p^-4 prefactors together are at most 1.621 and the complete
% geometric-tail coefficient is <549<700. The corresponding degree83
% cut-Gaussian approximation error is <5*10^-34.
%
% The same variance bound gives ||P_2(k)||<=p^2/24. Since p<=17/4,
% ||P_2 phi||<=289/384. The centering term lies in Ran P_0 and P_2 phi
% in its orthogonal complement, pointwise. Their squared norms therefore
% add, giving the explicit uniform initial-response bound ||w(0)||<0.753.
% At this equatorial preparation <k_z>=0, so <phi,w(0)>=0 as well.
% This is a norm bound for the actual initial derivative, not for its
% propagated derivative or for an omitted preparation contribution.
%
% A large propagated response can contain an exactly removable global
% phase. For ANY prescribed real b_*(tau), let
% beta_*(tau)=int_0^tau b_*(s)ds,
% w_sharp=w+i beta_* u/epsilon.
% Then, exactly and in the full four-spinor space,
% i epsilon dot w_sharp=H w_sharp+(B-b_* I)u,
% w_sharp(0)=w(0),
% 2 Re<u,Qw_sharp>=2 Re<u,Qw>
% for every self-adjoint detector Q in the stated domain. The energy and
% frequency rows are unchanged too: their physical operator derivative
% remains B, not B-b_* I. Thus b_* is a representation choice,
% not a change of Hamiltonian, preparation or observable. It need not equal
% the unknown exact expectation <u,Bu>. If it does equal that expectation,
% the exact response stays orthogonal to u for the chosen initial phase.
%
% An explicit useful representation for the selected member is supplied by
% the fixed reference circular trajectory
% q=(cos theta,sin theta,0),
% k=sqrt6 q+sqrt3(-sin theta,cos theta,0), theta=tau/sqrt3.
% Its upper principal B symbol evaluates by direct polynomial contraction to
% b_*(tau)=45/8-(1/8)cos(4tau/sqrt3)-2sqrt2 sin(4tau/sqrt3),
% beta_*(tau)=45tau/8-(sqrt3/32)sin(4tau/sqrt3)
%                         +(sqrt6/2)[cos(4tau/sqrt3)-1].
% This reference function is usable for every nearby launch; equality with
% a finite-width mean is not assumed. For example
% sum_i k_i^4=243/4-(63/4)cos4theta-9sqrt2 sin4theta,
% sum_i q_i k_i^3=27sqrt6/4-(3sqrt6/4)cos4theta
%                                      -(15sqrt3/4)sin4theta,
% which give the displayed b_* in the full ordered-B principal symbol.
% Endpoint ordering, spin and lower-band terms remain in B-b_* I.
% Choosing this phase representation can prevent subtraction of a large
% unobservable central-action derivative. Its residual, dual and root
% certificates are the preceding ones with B replaced by B-b_* I and with
% the exact corresponding approximate-response definition. A simple valid
% graph bound is ||(B-b_* I)f||<=1180 E_3(f)+|b_*| ||f||; an actual
% improvement in the residual still has to be demonstrated. If a basis is
% differentiated with respect to launch or spacing, its W_z c or W_delta c
% contribution is retained. No finite-frequency coefficient is inferred.
% SOURCES: actual RED MP22/SOURCES, independent projected-position and
% Clifford checks above, and the direct full-matrix phase transformation
% and circular polynomial contractions derived here. No numerical work.
%
% BOUNDED NECESSARY INPUT-COMPONENT SPECIFICATION [R; accepted input N below].
% Purpose: supply the normalization and measured-centering data of the
% SELECTED native preparation, with their first carrier derivatives, for
% use by the full MP8 propagation/root calculation. This does not compute
% a different detector, a return root or a propagated spacing coefficient.
% No execution is requested until the actual implementation is read by the
% original author/RED pair. The scalar cutoff and native preparation remain
% exact definitions; polynomial/uncut evaluations have the bounds below.
%
% Fixed inputs are epsilon=2^-24 and the selected S, V and B_phase.
% For an unambiguous implementation, form the horizontal 4-by-4 matrix
% K_h=[[B_h^T,C_h],[-A_h,-B_h]], in order (q_x,q_y,k_x,k_y), with
% A_h=[[-3/2,3sqrt2/2],[3sqrt2/2,3]],
% B_h=[[sqrt(2/3)/2,-1/sqrt3],[1/sqrt3,-sqrt(2/3)]],
% C_h=[[1,-sqrt2],[-sqrt2,2]]/9.
% Put P_h=3K_h^2 and
% S_h=I+(cosh(pi)-1)P_h-sqrt3 sinh(pi)K_h.
% Embed S_h in indices (0,1,3,4) of the 6-by-6 order
% (q_x,q_y,q_z,k_x,k_y,k_z); the remaining vertical block is -I_2.
% Set G=S S^T, V=G_kk, B_phase=G_qk V^-1 and Sigma=epsilon V/2.
% All these matrices are FIXED under the following carrier derivatives.
% The evaluated carrier is k_c=(sqrt6,sqrt3,0). The implementation should
% expose the same moment function for any equatorial k_c in
% |k_c-k_*|<=1/4, since that function is needed at the eventual return root.
% First derivatives are with respect to its first two components, at fixed
% z=k-k_c and fixed Sigma/B_phase. No finite differences are needed.
%
% Required outputs, including both first carrier derivatives, are:
% Z=E_g[(1+n_c.k/p)/2], normalization factor Z^-1/2,
% mu=Z^-1[B_phase E_g(z(1+n_c.k/p)/2)
%                         -(epsilon/4)E_g((n_c cross z)/p^2)],
% m_z=<M_z>=[24Z]^-1 E_g[(k_x k_y^3-k_y k_x^3)
%                                  (p^-3+(n_c.k)p^-4)].
% Here n_c=k_c/|k_c| and E_g is the ACTUAL normalized cut scalar density.
% Also output q_c=e_1-mu for the reference target center, its carrier
% derivatives, and the raw native centroid derivative epsilon*m_z in the
% z direction. The other native-centering components vanish by parity.
% The reflected preparation is obtained by the already prescribed exact
% reflection, not a separate replacement preparation. No finite-width
% expectation is replaced by its central-carrier value.
%
% Representation: no spatial/angular grid or tensor quadrature is needed.
% For j=1,...,4 replace p^-j in these moment functions by
% L_j(k)=kappa^-j sum_{l=0}^{39} binom(-j/2,l)t^l,
% kappa=|k_c|, t=(2k_c.z+z.z)/kappa^2.
% The polynomial expectations under the UNCUT Gaussian are then evaluated
% exactly as finite Gaussian-moment series, with interval roundoff below.
% Their numerators have degree at most five in k; their high powers of t
% need not be expanded into 83-degree monomials.
%
% A concrete 40-term generating-function implementation is as follows.
% Let Tmat=2Sigma/kappa^2, ell=2k_c/kappa^2 and use a formal variable xi.
% D(xi)=I-xi Tmat,
% C(xi)=D(xi)^-1 Sigma,
% m(xi)=k_c+xi D(xi)^-1 Sigma ell.
% The Gaussian quadratic-form generating function is
% F(xi)=E exp(xi t)
% =det D(xi)^(-1/2)
%   *exp[(xi^2/2)ell^T D(xi)^-1 Sigma ell].
% Only coefficients through degree39 are needed. In particular
% [xi^n]log F=tr(Tmat^n)/(2n)
%       +1_{n>=2} ell^T Tmat^(n-2)Sigma ell/2,
% F_0=1, F_n=(1/n)sum_{l=1}^n l[log F]_l F_{n-l}.
% Let W_alpha(xi) be the raw Gaussian moment with mean m(xi) and covariance
% C(xi), computed recursively by W_0=1 and, for beta=alpha-e_i,
% W_alpha=m_i W_beta+sum_j beta_j C_ij W_(beta-e_j).
% Then E[k^alpha t^l]=l![xi^l](F W_alpha).
% The 56 multi-indices |alpha|<=5 and these 40-term series suffice for
% every required output. First carrier derivatives can use value/d_x/d_y
% jets through exactly these algebraic recursions; differentiate kappa,
% n_c, Tmat, ell, the numerator coefficients and the final quotients.
% The shape matrices remain fixed. The finite coefficients are real.
%
% Certified truncation/cutoff target: final absolute error <=10^-20 for
% every reported scalar value and first derivative. The following explicit
% conservative budget is sufficient, uniformly on the stated neighborhood.
% Write c=104/121 and
% T_40=(13/484)^40 sqrt[Gamma(83/2)/Gamma(3/2)]<6*10^-37.
% On the actual support r=|z|<=1, for j=1,...,4,
% |p^-j-L_j|<=123(cr)^40,
% ||D_{k_c}(p^-j-L_j)||<=6300(cr)^40.
% To check the derivative constant, |D_{k_c}t|<=0.9r and the coefficient
% bound |binom(-j/2,l)|<=l+1 give, with L=(1-c)^-1,
% sum_{l=40}^infinity(l+1)c^(l-40)=41L+cL^2,
% sum_{l=40}^infinity l(l+1)c^(l-40)
%                       =1640L+81cL^2+c(1+c)L^3.
% Multiplication by kappa^-j and differentiation of that factor yield
% the stated constants. Numerator values are bounded by 700 and their
% carrier gradients by 1800; the phase matrix has norm <=8192.
% With Z>0.965, |mu|<0.05 and the proved derivative bounds, direct product
% and quotient rules bound ALL output/derivative truncation errors by
% 10^12 T_40<6*10^-25. For example the largest unnormalized mu-gradient
% error is below 2*10^4(1800*123+700*6300)T_40<10^11 T_40.
% The larger 10^12 constant includes the finite quotients and normalization.
%
% Removing the scalar cutoff from these POLYNOMIAL expectations costs less
% than 10^-42 in the same values/derivatives. One direct bound uses the
% already proved degree<=83 Gaussian tails, evaluates the absolute
% polynomial coefficients at |z|=1/2 (where the absolute t polynomial is
% <=48/121), and includes the cutoff normalization. A very loose coefficient
% bound 10^12 times sqrt27 exp(-128) is already below 10^-42.
% Thus enlarging each final finite-polynomial interval by 2*10^-24 covers
% these analytic errors as well as its interval roundoff, subject to the
% width test below. This is a rigorous input allowance, not a propagation
% error estimate or an assumption that the return root lies at k_*.
%
% Roundoff/implementation acceptance: use directed interval arithmetic with
% at least 100 decimal digits for pi, square roots, cosh/sinh, matrix
% products, the 3-by-3 inverse and all series/jet operations. A cofactor
% inverse with a determinant interval strictly above zero is sufficient.
% The exact symmetries V=V^T and B_phase=B_phase^T must be enclosed; do not
% replace a failed enclosure by an unexplained symmetrization. Before the
% analytic enlargement, every output interval must have width <=10^-25.
% Report the enlarged endpoints and maximum width, and verify that they
% obey the already proved Z>0.965, |mu|<0.05, ||D_{k_c}mu||<0.34 and
% |m_z|<=17/48 bounds. Check F_0=1 and the zero/first/second Gaussian
% moments against 1,k_c,k_c k_c^T+Sigma, retaining their interval residuals.
% These are finite algebraic checks, not empirical grid extrapolation.
%
% Bounded resources: at most 56 moment series of length40, each carrying
% three real jet components, plus their matrix and temporary series.
% Direct truncated convolutions require fewer than about 2*10^6 interval
% multiply-adds; no dense 88k/102k input basis is formed. The estimate is
% tens of seconds to at most a few minutes in ordinary high-precision
% Python arithmetic; use a HARD 180-second, 512-MiB, one-CPU, nice10,
% no-GPU cap. Exceeding the cap or any enclosure test is a failed component
% check, not permission to broaden the run. Save the exact producer,
% parameter intervals, output intervals, analytic budget, runtime and peak
% memory in the existing task artifacts. Full propagation remains open.
% SOURCES: exact selected-member definition and checked MP21--MP23 input
% formulas; independent Gaussian quadratic-form integral, raw-moment
% recurrence and inverse-power/quotient estimates above. The actual complete
% RED MP26--MP27/SOURCES were read: its independent completed-square mean
% D^-1 k_c equals k_c+xi D^-1 Sigma ell, and its generating coefficients,
% first-derivative and cutoff budgets agree. The scientific specification is
% cleared within the original pair; actual producer review precedes execution.
% mpmath1.3.0 iv is an acceptable finite-operation implementation. Actual
% primary sources read were
% https://raw.githubusercontent.com/mpmath/mpmath/1.3.0/mpmath/ctx_iv.py
% and
% https://raw.githubusercontent.com/mpmath/mpmath/1.3.0/mpmath/libmp/libmpi.py .
% They implement precision control and directed interval endpoints for the
% required arithmetic, pi, exp and sqrt. Use (exp(pi)+/-exp(-pi))/2 for
% cosh/sinh, integer interval divisions for rationals, and exact dyadic
% endpoint floor/ceiling at 110 decimal places for outward serialization.
% No ordinary floating-point constant or decimal endpoint rounding may
% silently replace the interval enclosure.
% The actual complete declared_exterior_packet_input.py (446 lines) and
% RED MP28/SOURCES were read before execution. The fixed shape, all first
% jets, logarithmic Gaussian generating coefficients, 56 raw moments and
% inverse-power falling products agree with the specification. Its
% continuum-centering contraction retains the Berry cross product, and
% the native m_z numerator retains both degree4 and degree5 terms. The
% interval operations, exact-zero skipping, positive determinant and
% symmetry/moment checks preserve the prescribed enclosures. The integer
% dyadic floor/ceiling serialization is outward for either endpoint sign.
% Every output jet is enlarged only after its raw width is checked; the
% declared process caps match the bounded component. No implementation
% repair is needed. Original-pair clearance is for this component alone;
% actual output acceptance, propagation, return root and finite-frequency
% coefficient are not inferred from the implementation read.
% The parent corrected two interval-left Jet multiplications before any
% execution, since mpmath's conversion raises rather than dispatching to
% the Jet reverse operation. The complete corrected producer was reread:
% both epsilon products and mixed dot contractions now put Jet first.
% This changes Python dispatch only, with no mathematical formula change.
%
% ACTUAL INPUT-COMPONENT RESULT [N, preparation inputs only].
% The entire899-line actual
% results/declared_exterior_packet_input_20261015_1335/result.json
% was read, including every raw/enlarged output jet, fixed-shape parameter,
% determinant, moment/symmetry residual, analytic budget and execution field.
% Its producer hash is the reviewed
% 2a4f6055ce6b3f00a53e9d1b56d9528215e2098aedcd9ee48a49defbb33ed822.
% The single run passed in2.34819seconds with27648KiB peak RSS, below the
% unchanged180-second/512-MiB cap. All raw output widths satisfy10^-25;
% all enlarged value/first-derivative intervals have diameter below10^-20.
% The exact-shape symmetry and Gaussian moment residuals enclose zero,
% and the covariance determinant interval is strictly positive.
% At the specified reference carrier, approximate central values are
% Z=0.99999990464754005262654,
% Z^-1/2=1.00000004767623338322136,
% mu=(-7.71558021515884*10^-12,
%      1.09388058371010*10^-11,-3.89983379826127*10^-14),
% m_z=-0.03928359850809473903212,
% epsilon*m_z=-2.34148493457405204*10^-9.
% The full stored endpoints, not these abbreviated central values, are
% the certified inputs. In particular they give |mu|<1.339*10^-11 and
% ||D_(k_cx,k_cy)mu||_F<1.083*10^-10 at this carrier, and
% D_(k_cx,k_cy)m_z approximately(-0.0481128916024370,0.0453613788296950).
% These are point evaluations with the specified analytic allowances;
% the earlier uniform-neighborhood bounds remain the appropriate bounds
% away from the reference carrier. The accepted component supplies no
% propagated residual, return-root enclosure or finite-frequency coefficient.
% SOURCES: complete actual corrected producer, reviewed specification and
% complete actual result above. The interval endpoints include the proved
% truncation and actual-cutoff allowance; acceptance is not a comparison
% of rounded high-precision runs. No new execution is requested here.
%
% The actual RED MP24--MP25/SOURCES also supply useful sharper bounds for
% continued propagation work. Their exact position rows give a reference
% quadratic Gaussian position-frame norm <39 and ||C_xk||<4992. Gaussian
% integration by parts, with the real cutoff term retained, yields
% B_phase E_g[z s]=(epsilon/2)C_xk
%       [E_g grad_z s+E_0(s grad_z theta^2)/E_0 theta^2].
% Independent checking confirms |mu|<9*10^-7 and ||D_{k_c}mu||<4*10^-6
% for the SAME actual input throughout the stated carrier neighborhood.
% For instance ||D n(k)-D n(k_c)||<=2|z|/(p kappa) follows by splitting
% (I-nn^T)/p-(I-n_c n_c^T)/kappa and using the exact norm of the
% difference of the two rank-one projectors. This bounds the differentiated
% integration-by-parts term directly. The input specification keeps its
% already valid more conservative checks; no definition is changed.
% The same RED source proves a radius1/8 tail <15exp(-172) for its explicitly
% stated UNPROJECTED, UNCUT reference quadratic Gaussian and an exterior-flow
% degree79 coefficient approximation error <13exp(-86) on that Gaussian.
% The row norms, binomial remainder and incomplete-gamma estimates were
% checked independently. These are coefficient/comparison-state bounds;
% actual full-matrix propagation, amplitude/graph/label errors and the MP8
% return certificate remain outstanding. They do not justify replacing the
% chosen native state by a quadratic Gaussian at later times.
%
% EXACT LINEAR-CORE PROPAGATION COMPONENT [R, independent WIP].
% Saved before exchanging this new propagation calculation with the RED.
% This uses the ALREADY SELECTED core; it changes no state or detector.
% In scaled coordinates its inner region r<=4/5 has v=-gamma q, where
% gamma=sqrt(2250/2197). There the full continuum operator is
% H_lin=alpha.k-gamma(q.k+k.q)/2, k=-i epsilon partial_q.
% As a comparison operator on all space it has the explicit unitary
% momentum propagator
% (U_lin(t)psi)(k)=y^(3/2) exp[-i alpha.k(1-y)/(gamma epsilon)]psi(y k),
% y=exp(-gamma t). This follows from the momentum transport equation
% partial_t psi+gamma k.partial_k psi
%                  =-3gamma psi/2-i alpha.k psi/epsilon.
% It retains all four spinor components. Since P_sigma(k) is homogeneous
% of degree zero, [H_lin,P_sigma]=0 exactly; no additional adiabatic
% band leakage is generated by this linear segment. This is NOT an
% identity for propagation under the entire fixed, joined exterior.
%
% The same comparison has an exact native spacing derivative. Put
% delta=(a omega)^2, s_delta,i=sin(sqrt(delta)k_i)/sqrt(delta),
% Q_delta,i=sin(2sqrt(delta)k_i)/(2sqrt(delta)). The LINEAR-FLOW native
% generator in momentum space is
% alpha.s_delta-i epsilon gamma[Q_delta.partial_k+div Q_delta/2].
% On the low-cone chart its characteristic obeys
% tan(sqrt(delta)k_i(t))=exp(gamma t)tan(sqrt(delta)k_i(0)).
% The back map at fixed final k is
% k_0,delta=y k+delta y(1-y^2)k^(3)/3+O(delta^2),
% with square-root Jacobian
% y^(3/2)[1+delta(1-y^2)|k|^2/2]+O(delta^2).
% At intermediate s, write lambda=exp[-gamma(t-s)]; the kinetic vector
% along that backward characteristic is
% lambda k+delta(lambda/3-lambda^3/2)k^(3)+O(delta^2).
% The different Dirac matrices along it must be time ordered; the following
% integral performs their complete first derivative rather than dropping
% their noncommutativity.
%
% For p=|k|, Q=alpha.k/p, V=alpha.k^(3), define
% V_parallel=(sum_i k_i^4/p)Q, V_perp=V-V_parallel,
% beta=[1-2y+y^3]/(6gamma),
% J=(1/gamma)int_y^1(1/3-lambda^2/2)
%                  exp[2ip(lambda-y)/(gamma epsilon)]d lambda,
% I=beta V_parallel+Re(J)V_perp+i Im(J)Q V_perp.
% Here I is Hermitian, since Q anticommutes with V_perp. J is an elementary
% quadratic-polynomial oscillatory integral; its coincident-phase limit
% is taken by its integral or entire divided differences. Direct matrix
% Duhamel differentiation gives the exact propagator derivative
% (partial_delta U_lin,delta psi)|_0(k)
% =y^(3/2) exp[-i alpha.k(1-y)/(gamma epsilon)]
%   {[(1-y^2)p^2/2-i I/epsilon]psi(yk)
%        +[y(1-y^2)/3]sum_i k_i^3(partial_i psi)(yk)}.
% It includes the kinetic correction, endpoint two-step drift, half-density
% Jacobian, momentum-argument change and full spin transport. At small t
% its derivative is -it B_lin/epsilon, with
% B_lin=-alpha.k^(3)/6
%       +i epsilon gamma sum_i[2k_i^3 partial_i/3+k_i^2],
% exactly the ordered restriction of the original B. The actual native
% preparation derivative must still be propagated and added separately.
%
% A useful finite bound follows without resolving the rapid phase.
% Integrating J by parts and using f(lambda)=1/3-lambda^2/2 gives
% |J|<=epsilon[|f(1)|+|f(y)|+int_y^1|f'|]/(2p)
%      <=epsilon/(2p), for 0<y<=1.
% Since ||V_perp||=|k^(3)-k sum_i k_i^4/p^2|<=p^3/2,
% the off-diagonal spin part of I/epsilon has norm at most p^2/4.
% For P_sigma psi=psi, the derivative of its incoming projector under
% the back map contributes at most (1-y^2)p^2/12. Consequently the
% COMPLETE propagator-only opposite-free-band derivative satisfies
% ||P_-sigma(k)(partial_delta U_lin,delta psi)|_0(k)||
%       <=y^(3/2)(p^2/3)||psi(yk)||,
% and its L2 norm is at most y^-2|| |k|^2 psi||/3. There is no hidden
% epsilon^-1 loss in this off-diagonal estimate. The diagonal phase
% response and the actual initial P_2/centering term remain present.
%
% The lower reference characteristic also has an exact inner segment.
% Its initial q=e_1,k=(sqrt6,sqrt3,0) has e_-=-5 and angular momentum
% magnitude ell=sqrt3. With u(r)=|v(r)| and radial momentum k_r,
% p+u k_r=5 and
% D(r)^2=25-(1-u^2)ell^2/r^2,
% k_r=[-5u+/-D]/(1-u^2), dot r=minus/plus D/p.
% The plus momentum branch is incoming. It reaches r=4/5 without a
% radial turn, since D^2>=25-3/(4/5)^2>0 for 4/5<=r<=1. On that
% interval k_r>0, p<=5, so its arrival time is below
% (1/5)*5/sqrt[25-3/(4/5)^2]<0.223 in scaled units.
% Inside, k(s)=exp(gamma s)k_a and
% q(s)=exp(-gamma s)q_a-[1-exp(-gamma s)]n_a/gamma.
% Its turning radius is EXACTLY
% r_min=sqrt[3/(25+3gamma^2)], strictly between0.32 and0.34.
% At r_c=4/5 write u_c=gamma r_c and
% D_c=sqrt[25+3gamma^2-3/r_c^2]. The positive momenta on entry and exit are
% p_in=(5-u_c D_c)/(1-u_c^2),
% p_out=(5+u_c D_c)/(1-u_c^2),
% and the inner residence time is gamma^-1 log(p_out/p_in).
% These statements concern the reference lower characteristic, not the
% entire finite-width packet or an assumed measured return root.
%
% To use this exact component in the ORIGINAL propagation certificate,
% its matching error must be retained. If chi(t,q) is a smooth cutoff
% supported strictly inside the linear region and u_lin,w_lin solve its
% complete base/response equations, then the actual residuals of
% tilde u=chi u_lin, tilde w=chi w_lin are exactly
% r=i epsilon[partial_t chi+(alpha+v).grad chi]u_lin,
% r_1=i epsilon[partial_t chi+(alpha+v).grad chi]w_lin+[chi,B]u_lin.
% The commutator has order at most two and is supported on cutoff gradients.
% Its actual graph/weighted tails, joining pieces and parameter derivatives
% must be bounded before using the component. The explicit propagator does
% not establish those bounds, a sufficient finite propagated degree, an
% actual MP8 root or the requested absolute coefficient. No execution is
% specified by this analytic component.
% SOURCES: direct native momentum characteristics, half-density Jacobian,
% Clifford time-ordering integral, and conserved radial energy/angular
% momentum for the already fixed core. No comparison result was read before
% this independent save. The earlier complete RED MP29/SOURCES was read:
% its orthogonal free-band frame and coordinate connection agree, while
% its required cutoff/Gram controls remain necessary for a Galerkin use.
% The actual independently saved RED MP31/SOURCES was then read. Its
% initial Hermite-cutoff Gram estimate checks by the Mehler diagonal sum:
% with |nu|<=D and 0<rho<1,
% ||I-G_D||<=rho^-D(1-rho)^-3
%   tau_3(16 sqrt[(1-rho)/(1+rho)]).
% The factor (1-rho)^-3 includes all three-dimensional degeneracy. For
% rho=1/4 the bound34*4^D exp(-768/5) is below7*10^-16 at D=83 and
% below1.1*10^-5 at D=100. Exact free frames leave each scalar Gram block
% unchanged and opposite-band blocks zero, so the block operator norm
% needs no extra spin multiplicity. This certifies initial conditioning
% at those degrees only; neither degree is selected for propagation.
% Moving cutoffs, Gram derivatives, and actual residual/return control
% remain separate requirements.
% The complete actual RED MP32/SOURCES was then read. Its characteristic
% derivation and Clifford conjugation independently verify the full native
% first response, including the incoming-projector derivative in the1/3
% opposite-band bound. Its cutoff-commutator convention -[B,chi] agrees
% exactly with [chi,B] above. The additional reference-characteristic
% bounds also check: u_c^2=1440/2197 and
% D_c^2=325/16+6750/2197 imply3.91<u_c D_c<3.92, hence
% p_out/p_in<8.26, residence<2.12 and y^-2<69.
% These numerical inequalities follow from the displayed exact rationals
% and elementary logarithm bounds, not a propagation run. They remain
% reference bounds; no uniform finite-packet residence or matching-tail
% control is inferred. This original-pair component comparison has no
% mandatory repair and does not promote the open MP8 coefficient.
%
% EXACT LOCAL NONLINEAR GAUSSIAN COUPLINGS [R, independent WIP].
% This is a finite ingredient for the selected moving-phase propagation,
% derived before intake of the RED's new mean/covariance continuation.
% It is not a scalar replacement of the original matrix dynamics.
% In the rotating reference frame put q=e_1+x,
% k=k_*+z, k_*=(sqrt6,sqrt3,0), n_*=k_*/3, t=n_*.z,
% C=sqrt(2/3). The upper principal Hamiltonian is
% e_rot=|k|-C(q.k)|q|^(-3/2)
%                   -(q_x k_y-q_y k_x)/sqrt3.
% Its cubic and quartic Taylor terms are particularly small finite
% polynomials. Define
% f1=-3x_1/2,
% f2=-3|x|^2/4+21x_1^2/8,
% f3=21x_1|x|^2/8-77x_1^3/16,
% f4=21|x|^4/32-231x_1^2|x|^2/32+1155x_1^4/128,
% L0=sqrt6, L1=z_1+sqrt6 x_1+sqrt3 x_2, L2=x.z.
% These are the homogeneous pieces of
% (1+2x_1+|x|^2)^(-3/4) and q.k, respectively. Then
% h3=-t(|z|^2-t^2)/18-C(L0 f3+L1 f2+L2 f1),
% h4=-[|z|^4-6t^2|z|^2+5t^4]/216
%                         -C(L0 f4+L1 f3+L2 f2).
% The rotating angular-momentum term is quadratic and supplies no term
% here. All expressions use the SAME fixed core, which is exactly the
% exterior expression in a neighborhood of this reference circle.
%
% Removing the exact quadratic metaplectic motion gives, at the level of
% this principal scalar symbol, the interaction
% sqrt(epsilon) Op_W[h3(S_tau xi)]
%       +epsilon Op_W[h4(S_tau xi)]+higher terms,
% S_tau=exp[K(tau-pi sqrt3)], [xi_qi,xi_kj]=i delta_ij.
% This is an identity for the Taylor coefficients, not an error estimate
% for truncating their propagation. The fixed half-period squeeze is
% already included in S_tau; using an unsqueezed vacuum would change the
% chosen preparation.
%
% Every initial vacuum-to-Hermite coefficient of these operators is an
% explicit finite polynomial. Write
% C_Aj=(S_tau,Aq_j+i S_tau,Ak_j)/sqrt2,
% Sigma_AB=(S_tau S_tau^T)_AB/2.
% For any real polynomial f of the six canonical variables, define
% g_f(zeta)=[exp((1/2)Sigma_AB partial_A partial_B)f](C zeta).
% Repeated indices are summed over all six coordinates. Then exactly
% Op_W[f(S_tau xi)]|0>
%              =g_f(a^dagger)|0>,
% ||Op_W[f(S_tau xi)]|0>||^2
%              =sum_nu nu! |[zeta^nu]g_f|^2.
% The contraction factor follows from the Weyl generating exponential;
% for example q_i^2|0>=[(a_i^dagger)^2+1]|0>/2.
% Hence h3 has only degree1 and3 creation pieces (3+10 coefficients),
% while h4 has degrees0,2,4 (1+6+15 coefficients). No large carrier
% angular-momentum basis is needed to evaluate THESE finite couplings.
% The corresponding general matrix element is obtained from
% :[exp((1/2)Sigma_AB partial_A partial_B)f]
%                      (C a^dagger+conjugate(C)a): .
% The arguments inside the colons are commuting formal variables before
% normal-ordered substitution; no second contraction is performed.
% This keeps the number-preserving and annihilation terms required for
% subsequent Hermite amplitudes. A vacuum norm by itself is not an
% operator bound on an evolved finite space.
%
% The complete free-frame matrix propagation still contains its exact
% momentum-dependent frame connection, off-diagonal band terms, native
% preparation/centering response and differentiated frequency/return rows.
% They are not included in h3,h4 and cannot be discarded using these
% principal couplings. A finite propagated degree additionally requires
% bounds on the higher Taylor remainder, actual amplitude tails and
% parameter derivatives, followed by the residual/root certificate above.
% No sufficient degree, measured coupling, numerical run or selected
% member's return root is asserted here.
% SOURCES: direct Taylor expansion of the already selected exterior symbol,
% the independently derived constant quadratic generator K, and the exact
% Weyl generating exponential/creation-annihilation algebra. No new
% counterpart amplitude argument or numerical output was used.
% The actual RED MP33/SOURCES was then read. Its finite classical mean
% equation dot m=K m+(1/2)D^2 f:Gamma(t), with
% Gamma(t)=exp(Kt)G_* exp(K^Tt)/2, is precisely the contraction of this
% cubic Taylor term in the classical flow. It remains a principal moment
% component, not a matrix propagation replacement. Its initial frequency
% coefficient can be checked without a numerical matrix multiplication.
% In the horizontal order(q1,q2,k1,k2), write P=3K_h^2 as
% P=[[U,0],[L,U^T]], U=[[1,0],[-sqrt2,0]],
% L=3sqrt3[[0,1],[-1,0]]. With d=cosh(pi)-1 and s=sinh(pi),
% S_h=I+dP-sqrt3 s K_h. Direct block multiplication then gives
% (G_*)_q1,k2-(G_*)_q2,k1
%       =-(29sqrt3/3)s^2-(11sqrt6/6)d s.
% Since tr Hess e_rot=31/6 and S^T Hess e_rot S=Hess e_rot,
% the uncut scalar frequency coefficient is exactly
% E1=31/24-(29/6)s^2-(11sqrt2/12)d s.
% This confirms the need for the already prescribed frequency row; it
% is not the full spinor energy or its unevaluated Taylor remainder.
% The RED's exact reflected-pair kernel trace and double-commutator
% identity also check: for N_i=k_i/|k| and s0=(1+n_c.N)/2,
% <F>_pair=<g,{s0,F}g>/(2Z)
%               -sum_i<g,[N_i,[N_i,F]]g>/(8Z),
% before restoring the separately specified vertical translations.
% Here sum N_i F N_i=F-(1/2)sum[N_i,[N_i,F]] proves the factor and
% sign directly. Neither that exact identity nor traceless regular band
% corrections cancels the differentiated interband oscillation by itself.
% The actual cutoff covariance, scalar Taylor remainder, native frame and
% the full four-return response therefore remain in the certificate.
%
% COEFFICIENT-WEIGHTED MOVING-SPACE TAILS [R, independent continuation].
% A sufficient propagated degree must be tied to actual amplitudes, rather
% than inferred from the dimension of a polynomial input approximation.
% For scalar normalized Hermite columns g_nu in three whitened variables
% and any finite spin-vector coefficients c_nu, define
% A_rho(c)=sum_nu rho^(-|nu|)||c_nu||^2, 0<rho<1.
% Weighted Cauchy--Schwarz and the exact Mehler diagonal identity give
% int_{|xi|>R}||sum_nu c_nu g_nu(xi)||^2 dxi
%   <=A_rho(c)(1-rho)^(-3)
%                tau_3(R sqrt[(1-rho)/(1+rho)]).
% Unlike the worst-space rho^-D bound, this uses the actual amplitudes
% in high shells. The inequality also applies to time, spacing and launch
% derivative coefficient vectors separately. It is directly evaluable
% from a finite propagation, with certified coefficient errors added.
% It does not assert that their weighted norms stay small under evolution.
% Common real chirp/translation phases do not change this estimate.
% The exact free-band frame U_sigma(k) preserves the MOMENTUM density;
% opposite free-band columns are pointwise orthogonal even when their
% scalar Gaussian centers differ. A spatial tail for the full spinor is
% not obtained by silently treating that momentum-dependent frame as a
% constant: its derivatives or a controlled polynomial approximation are
% still required for coordinate/graph bounds.
%
% The complete actual RED MP34--MP35/SOURCES were read. The explicit Wick
% convention and midpoint contractions agree with the local algebra above.
% The fixed-whitened-cutoff construction in MP35 gives an exact useful
% moving domain. For a specified Gaussian momentum covariance V(t), take
% xi=V(t)^(-1/2)(k-k_c(t))/sqrt(epsilon),
% chi_fix(xi)=theta(|V_0^(1/2)xi|/2048)theta(|xi|/32).
% This second cutoff belongs to the approximation space; its omitted
% physical input is retained in the error. Change of variables makes its
% scalar Gram G_D exactly constant under center, covariance and chirp
% motion. The true W^dagger dot W remains; constancy of G_D does not
% remove the derivative of the columns. Multiplication by the exact free
% frames leaves G_D unchanged and the opposite-band Gram block zero.
%
% The proposed UPPER REFERENCE space has a checked momentum domain.
% With u=tau/sqrt3-pi, its full six-dimensional covariance factor obeys
% ||exp(sqrt3 u K)||_F^2
%  =122-116cosh u+(1561/12)sinh(u)^2+(88/3)sin(u)^2<128^2.
% In particular the vertical2x2 block has K_v^2=-I/3 and
% ||exp(sqrt3 u K_v)||_F^2=2+(88/3)sin(u)^2; it cannot be omitted
% at intermediate times. The full expression increases with |u| on
% [0,pi]. Symplecticity bounds its inverse as well. Thus both
% ||V^(1/2)|| and ||V^(-1/2)|| are below128. On the cutoff support
% |k-k_c|<=sqrt(epsilon)*128*64=2, so |k|>=1 at the reference
% carrier and >=3/4 if the chosen carrier changes by at most1/4.
% This is a domain of the approximation columns, not a proved tube for
% the actual finite-frequency state or its return labels.
%
% Its constant Gram defect is bounded by
% eta_D=34*4^D exp(-768/5)+68*4^D exp(-3072/5),
% so G_D^(-1/2) is well defined at the already stated input degrees.
% In coefficient tail estimates use c=G_D^(-1/2)b for orthonormal
% coordinates b; replacing c by b without this factor is not exact.
% No finite propagated degree is selected by this conditioning fact.
%
% One coordinate-cutoff contribution can now be bounded explicitly.
% The selected theta has |theta'|<=8, hence
% |grad_xi chi_fix|<=8*128/2048+8/32=3/4.
% That derivative is supported only where |xi|>16. If
% T_rho(16)=(1-rho)^(-3)
%            tau_3(16 sqrt[(1-rho)/(1+rho)]), then uniformly on the
% upper reference moving space
% ||epsilon(grad_k chi_fix)sum c_nu g_nu||
%       <=96sqrt(epsilon)[A_rho(c)T_rho(16)]^(1/2)
%       =(3/128)[A_rho(c)T_rho(16)]^(1/2).
% This is a full three-component derivative norm, not a per-coordinate
% estimate summed without its Euclidean factor. The inequality uses the
% matrix operator norm of V^(-1/2). The separate free-frame coordinate
% connection epsilon A_i, Gaussian top-shell term, normalization and
% temporal-cutoff derivatives remain in the physical coordinate/residual
% calculation. Their omission is not justified by this small shell bound.
%
% A convenient implementation may hold these reference columns independent
% of spacing and launch labels. Then W_delta=W_z=0 by construction;
% the actual initial projection and its derivatives still change with
% those labels and must be computed. If a different implementation adapts
% the columns to the labels, their W_delta and W_z terms return. Neither
% representation changes the fixed native preparation or MP8 conditions.
% SOURCES: independent weighted Cauchy--Schwarz/Mehler estimate followed
% by the complete actual MP34--MP35 original-pair intake, explicit vertical
% covariance calculation and cutoff chain rule. This supplies quantitative
% a posteriori controls for a future full-matrix propagation; no propagated
% coefficients, residual bound, sufficient degree or return root is asserted.
%
% FIRST MATRIX TERMS IN THE ROTATING HERMITE RECURSION [R, independent WIP].
% The local scalar couplings can be connected explicitly to the original
% four-spinor operator; a scalar Taylor polynomial alone is insufficient.
% Work at zero spacing in the rotating PHYSICAL spin frame, where
% H_rot=H-Omega J_z, Omega=1/sqrt3,
% J_z=q_x k_y-q_y k_x+epsilon Sigma_z/2.
% Use the exact free frames U_+=(I;sigma.n)/sqrt2 and
% U_-=(I;-sigma.n)/sqrt2, n=k/p. Define
% f(q,k)=v(q).k-Omega(q_x k_y-q_y k_x), F=grad_q f,
% C_j=sigma.(n cross e_j)/p.
% The Weyl symbol in that full free frame is
% diag(p,-p)+f I+epsilon h_1+O(epsilon^2),
% h_1=sum_j A_j F_j-(Omega/2)U^dagger Sigma_z U,
% A_j=(1/2)[[-C_j,C_j],[C_j,-C_j]].
% This follows by the two first-order Moyal products U^dagger#f#U;
% the coefficient is i U^dagger partial_k U.grad_q f, with no missing
% factor of two. The displayed remainder requires a smooth nonzero-p
% tube and its symbol/weighted bounds; it is not an established numerical
% error estimate for the entire evolution.
%
% The spin-rotation blocks are exactly
% (U^dagger Sigma_z U)_{++}=(U^dagger Sigma_z U)_{--}=n_z sigma.n,
% (U^dagger Sigma_z U)_{+-}=(U^dagger Sigma_z U)_{-+}
%                                      =sigma_z-n_z sigma.n.
% Consequently
% (h_1)_{++}=-(1/(2p))sigma.(n cross F)
%                                      -(Omega/2)n_z sigma.n.
% At the reference fixed point p=3, F=0 and n_z=0, so the diagonal
% subprincipal block vanishes THERE. Its first Taylor coefficient does
% not vanish. With A,B the previously specified q,q and q,k Hessian
% blocks of e_rot, and n_*=k_*/3, it is
% h_1,linear=-(1/6)sigma.
%                  [n_* cross(Ax+Bz)+Omega z_3 n_*].
% Therefore the first upper-band interaction after removal of the
% quadratic metaplectic motion is the two-spin polynomial
% sqrt(epsilon){h3(S_tau xi)I+h_1,linear(S_tau xi)}.
% It contains the thirteen scalar degree1/3 creation coefficients above
% and the additional degree1 spin matrix. Ignoring the latter would
% change the full-spinor amplitude recursion even though its trace is zero.
%
% The same fixed point has the NONZERO opposite-band block
% (h_1)_{-+}=-(Omega/2)sigma_z.
% Its first adapted-projector coefficient is
% D_{-+}=(h_1)_{-+}/(2p)=-sigma_z/(12sqrt3).
% A constant first-order block rotation with lower-left epsilon D_{-+}
% therefore puts the physical free-upper preparation into the opposite
% adapted band with coefficient -epsilon D_{-+} at this point. It must
% accompany the upper recursion; it is not a freely chosen new state.
% Its subsequent phase response remains part of the original native
% spacing derivative and cannot be omitted from the finite-frequency
% coefficient using its small undifferentiated norm.
%
% The initial upper amplitude also has a finite first correction. If
% sigma.n_c chi_eta=eta chi_eta, eta=+1 or-1 for the prescribed reflected
% chiral pair, the exact free-frame preparation coefficient is
% (I+eta sigma.n(k))chi_eta/(2sqrt(Z)).
% Its degree1 expansion about k_c is
% [chi_eta+eta sigma.[(I-n_c n_c^T)z]chi_eta/(2p_c)]/sqrt(Z)
%                                                     +O(|z|^2)/sqrt(Z).
% Since Z=1+O(epsilon), its normalizer does not change the subsequent
% first sqrt(epsilon) coefficient, but it is retained in this pointwise
% momentum expansion and in the actual input.
% At the reference point this gives the initial sqrt(epsilon) Hermite
% coefficient eta sigma.(I-n_*n_*^T)(S_0 xi)_k chi_eta/6.
% The normalization and previously computed actual centering are retained
% at their own orders; this local formula does not replace their certified
% exact input values. Reflection is kept as the stated two-state mixture.
%
% Thus the first base-amplitude equation has the explicit finite source
% h3(S_tau xi)I+h_1,linear(S_tau xi), plus the displayed input coefficient.
% It is a necessary part of a full recursion, not a completed truncation.
% Higher matrix symbols, opposite-band propagation, native P_2/centering,
% the spacing derivative of the rotating frame and all four return/frequency
% rows remain required. No finite error bound or execution specification
% follows from retaining only this first equation.
% SOURCES: direct physical rotation by total J_z, the exact full free
% frame from the already read MP29, first-order Weyl composition and
% elementary Pauli products. This block was derived before exchanging
% the new matrix-recursion calculation with the RED; no scalar result is
% promoted to a full microscopic propagation result.
%
% EXACT NATIVE-PROJECTOR RESPONSE GAUGE [R, independent continuation].
% A fixed physical rotation and the CONTINUUM free frame can be chosen
% independent of spacing. This avoids introducing an unnecessary moving
% native-band basis while retaining its effect exactly in the response.
% Put b_i(tau)=R_z(-Omega tau)e_i and p_i=b_i.k. In this fixed rotating
% frame the original spacing derivative is precisely
% B_rot(tau)=-(1/6)sum_i(alpha.b_i)p_i^3
%                      -(1/3)sum_i{b_i.v(q),p_i^3}.
% The rotation generator -Omega J_z has no spacing derivative. This
% operator formula includes the endpoint ordering and its spin matrices;
% it is not merely a Weyl scalar. Define the rotated cubic vector and
% quartic invariant
% w(tau,k)=sum_i b_i(b_i.k)^3, S4=sum_i(b_i.k)^4,
% n=k/p, d n=-w/(6p)+k S4/(6p^3),
% P_2=(alpha.d n)/2, P_0=(I+alpha.n)/2,
% T(tau,k)=[P_2,P_0]=i Sigma.(k cross w)/(12p^2).
% T is anti-Hermitian, [T,P_0]=P_2 and T P_0=P_2 P_0. Moreover
% ||T||<=p^2/24 by |k cross w|<=p^4/2. Its homogeneous expression
% is defined as zero at k=0; smooth symbol estimates are used on the
% already stated nonzero-momentum domains.
%
% For the exact full-spinor base u and spacing response w_resp, set
% z_resp=w_resp-T u. Direct differentiation, with no expansion in epsilon,
% gives
% i epsilon dot z_resp=H_rot z_resp+B_eff u,
% B_eff=B_rot+[H_rot,T]-i epsilon dot T.
% The kinetic commutator is exactly
% [alpha.k,T]=-p alpha.d n.
% It cancels the entire leading opposite-free-band kinetic part of B_rot;
% the remaining kinetic piece is -(S4/(6p))alpha.n.
% The rest is still a complete matrix operator. In particular the scalar
% flow F0=sum_j{v_j(q),k_j}/2 contributes exactly
% [F0,T]=(1/2)sum_j{[v_j(q),T(k)],k_j},
% and the total-rotation contribution -Omega[J_z,T] is retained along
% with -i epsilon dot T. No symbol remainder is hidden by this identity.
% On the rotating tensor T(tau)=U_rot(tau)T(0)U_rot(tau)^dagger,
% dot T=(i Omega/epsilon)[J_z,T]; therefore these two rotation terms
% cancel exactly. Equivalently
% B_eff=-(S4/(6p))alpha.n
%       -(1/3)sum_i{b_i.v,p_i^3}+[F0,T].
% Its explicit time dependence remains through the cubic lattice axes.
% This final equality uses the fixed reference angular velocity, not an
% assumed actual packet angular velocity or changed return period.
%
% The actual native input makes this gauge particularly useful:
% w_resp(0)=P_2 phi+i m.k phi=T(0)phi+i m.k phi,
% hence z_resp(0)=i m.k phi. Here m is the already certified exact
% native-centering coefficient, not a central-carrier substitution.
% The sizeable P_2 input is represented explicitly by T u; it has not
% been thrown away. For any self-adjoint detector Q,
% d_delta <Q>=2Re<u,Q z_resp>+<u,[Q,T]u>.
% For the native coordinate in this continuum response problem,
% [q_i,T]=i epsilon partial_{k_i}T exactly. At the initial free-upper
% preparation, P_0 i(partial_i T)P_0
%                  =P_0{-i[partial_i P_0,P_2]}P_0.
% Thus the explicit detector term recovers the earlier epsilon m_i
% shift and cancels its prescribed translation at launch. At later times
% the full i partial_i T is retained, including its opposite-band blocks.
%
% The initial physical frequency row must use
% <B+[H,T]>+2Re<phi,H z_resp(0)>, with the PHYSICAL H and B.
% It is not obtained by substituting a dynamically phase-subtracted source
% or by omitting the explicit commutator. The mean-radius, half-time and
% final coordinate detectors use the displayed [Q,T] term; changing to
% this exact response variable changes none of the MP8 conditions.
%
% The remaining commutator [F0,T] is an exact operator of first symbol
% order in epsilon, but a quantitative graph/tube bound is still needed
% before using that order in an error budget. This identity alone does
% not certify a propagated degree, the base interband leakage, launch
% derivative controls or an actual return root. It supplies a full-matrix
% response formulation that removes the leading kinetic interband source
% without an epsilon^-1 norm estimate or a change of physical preparation.
% SOURCES: direct differentiation of the native free projector, the exact
% fixed total-angular-momentum rotation and operator commutator algebra.
% This response-gauge block was independently saved before new counterpart
% intake. All previously declared state, frequency and detector terms remain.
%
% EXPLICIT GLOBAL FLOW-COMMUTATOR FORM BOUND [R, independent continuation].
% The remaining gauge commutator admits an elementary quantitative bound;
% no unquantified pseudodifferential constant or global spectral gap is
% needed for this step. For a scalar real f with bounded derivatives let
% M0=||f||_infinity, M1=||grad f||_infinity,
% M2=sup||Hess f||_op, M3=sup_{x,|e|=1}|D_e^3 f(x)|.
% Write p=|k| as an operator multiplier. For every R>0 and psi in H^1,
% ||[f,T]psi||<=epsilon a_f(R)||p psi||
%                                      +epsilon^2 b_f(R)||psi||,
% a_f=(5sqrt3/12)M1+(35/16)M2 R,
% b_f=(35/8)[M3 R/6+M0/R^2+M1/R].
% These constants are uniform in the rotated cubic axes.
%
% Here is a direct kernel proof. Each component H4_i(k) of k cross w
% is a homogeneous HARMONIC polynomial of degree4: the two Laplacians
% in k_j k_l^3-k_l k_j^3 cancel. Since
% H4_i(partial)(1/|x|)=105 H4_i(x)/|x|^9, the finite-part kernel of T is
% K_T(y)=i epsilon^2 [105/(48pi)] Sigma.H4(y)/|y|^9.
% There is no local degree2 polynomial term with this degree4 angular
% harmonic. The bound |H4(y)|<=|y|^4/2 therefore gives the angularly
% integrated majorant (35/8)epsilon^2 r^(-5).
% The exact first Taylor term of the commutator is
% i epsilon sum_j(partial_j f)(partial_{k_j}T)(k)psi.
% Direct differentiation of T gives ||D_e T||<=5p/12 for |e|=1;
% Euclidean Cauchy--Schwarz over the three multiplier derivatives bounds
% this term by (5sqrt3/12)epsilon M1||p psi||.
%
% For the remainder on |y|<R, split psi(x-y)=psi(x)+[psi(x-y)-psi(x)].
% The latter difference has L2 norm at most |y| ||p psi||/epsilon.
% Multiplication by the second Taylor remainder of f gives
% (35/16)epsilon M2 R||p psi||. In the term multiplying psi(x), the
% quadratic Taylor polynomial integrates to zero against the degree4
% angular harmonic, so its third remainder costs only
% (35/8)epsilon^2 M3 R||psi||/6. On |y|>=R retain the elementary bound
% 2M0+M1|y| after the first Taylor subtraction. Its two radial integrals
% give (35/8)epsilon^2(M0/R^2+M1/R)||psi||. These absolutely convergent
% remainders justify the finite-part identity by approximation from
% smooth compact states and then its extension to H^1.
%
% For EACH component f=v_j of the selected fixed core, the already proved
% Cartesian constants imply M0<1, M1<61, M2<=24300 and M3<2*10^7.
% Choosing the purely analytical splitting radius R=1/1000 yields
% a_f<100 and b_f<4.7*10^6. Thus this exact commutator has the explicit
% global estimate
% ||[v_j,T]psi||<=100epsilon||p psi||
%                                    +4.7*10^6 epsilon^2||psi||.
% The splitting radius is not a new state, source-core radius, smoothing
% length or numerical grid. All original functions and epsilon are fixed.
%
% A form estimate suffices for the dual residual pairing and avoids asking
% for an unproved fourth derivative of the joined core. On chi in H^1
% and psi in H^2, the exact identity
% [F0,T]=(1/2)sum_j{[v_j,T],k_j}
% and Cauchy--Schwarz give
% |<chi,[F0,T]psi>|
% <=(sqrt3/2){100epsilon[||chi||||p^2 psi||
%                                  +||p chi||||p psi||]
%       +4.7*10^6 epsilon^2[||chi||||p psi||+||p chi||||psi||]}.
% This is explicitly a FORM bound, not an L2 norm bound for [F0,T].
% It can be inserted in a primal/dual error pairing with the stated
% momentum moments; no hidden fourth flow derivative or omitted band
% contribution is required. A forward L2 source estimate would need its
% own additional argument. The bound quantifies this one remaining gauge
% term but still supplies no actual propagated amplitude/root certificate.
% SOURCES: direct harmonic-polynomial Coulomb kernel, exact Taylor
% cancellation and translation estimate in H^1, with the selected core's
% already independently checked M1--M3. This bound was saved before any
% new counterpart commutator estimate or scientific execution.
%
% PROPOSED BOUNDED FIRST MATRIX-AMPLITUDE COMPONENT [R, specification].
% PURPOSE AND USE. Compute the actual finite first coefficient a1 in the
% selected upper moving-Hermite recursion, including BOTH spin components,
% the prescribed reflected input and its first preparation coefficient.
% This coefficient is used by the next amplitude source V3 a1+V4 a0 and
% by the residual/tail estimates for that same recursion. It is not a
% new observable, a scalar replacement, a changed input run, or an error
% estimate for keeping only a0+sqrt(epsilon)a1. The full lower-band and
% native-response equations, the actual root and the absolute coefficient
% remain required. The finite calculation below has no time integrator,
% no fitted parameter, no angular lattice grid and no physical model change.
% Here V4 is the FULL matrix next coefficient: besides the scalar quartic
% term it retains the second Taylor term of h1 and the adapted second
% symbol Q2_eff. The opposite band is propagated in its own equation.
%
% FIXED INPUTS. Use the same K, h3, h1_linear, epsilon=2^-24 and
% reference carrier n_*=(sqrt(2/3),1/sqrt3,0). Let
% K=[[B^T,C],[-A,-B]] in the order(q1,q2,q3,k1,k2,k3), where
% A=[[-3/2,3sqrt2/2,0],[3sqrt2/2,3,0],[0,0,3]],
% B=[[sqrt(2/3)/2,-1/sqrt3,0],
%    [1/sqrt3,-sqrt(2/3),0],[0,0,-sqrt(2/3)]],
% C=[[1,-sqrt2,0],[-sqrt2,2,0],[0,0,3]]/9.
% For commuting x,z in R^3 put t=n_*.z, Cflow=sqrt(2/3),
% f1=-3x1/2, f2=-3|x|^2/4+21x1^2/8,
% f3=21x1|x|^2/8-77x1^3/16,
% L0=sqrt6, L1=z1+sqrt6 x1+sqrt3 x2, L2=x.z. The exact inputs are
% h3=-t(|z|^2-t^2)/18-Cflow(L0 f3+L1 f2+L2 f1),
% h1_linear=-(1/6)[sqrt3 x3 sigma_x+(-sqrt6 x3+z3)sigma_y
%       +((3sqrt3/2)x1+(sqrt6/2)x2+(sqrt2/6)z1-z2/3)sigma_z].
% K_h and K_v mean the horizontal and vertical submatrices of THIS K.
% These definitions reproduce the earlier independently checked entries.
% Set
% u=tau/sqrt3-pi, -pi<=u<=pi, and S(u)=exp(sqrt3 u K).
% Store its EXACT five-frequency representation
% S(u)=S_0+S_+ exp(u)+S_- exp(-u)+S_i exp(iu)+S_-i exp(-iu).
% In the horizontal order(q1,q2,k1,k2), its coefficients are
% S_0=I-P, S_+=(P+sqrt3 K_h)/2, S_-=(P-sqrt3 K_h)/2,
% P=3K_h^2. In the vertical order(q3,k3), they are
% S_i=(I-i sqrt3 K_v)/2, S_-i=(I+i sqrt3 K_v)/2.
% Unspecified blocks of these embedded6x6 matrices are zero. In particular
% this S_0 denotes a frequency coefficient here, not the initial S(-pi).
% The initial scalar envelope and all state definitions remain unchanged.
%
% Use the Pauli matrices and the two carrier spinors
% chi_+=(1,(sqrt2+i)/sqrt3)/sqrt2,
% chi_-=(1,-(sqrt2+i)/sqrt3)/sqrt2.
% Their sigma.n_* eigenvalues are eta=+1 and-1, respectively. Retain
% them as separate two-spin coefficient arrays before the stated mixture.
% For each eta and each normalized Hermite index nu with |nu|=1 or3,
% define the source coefficients v_eta,nu(u) by
% {Op_W[h3(S(u)xi)]I+h1_linear(S(u)xi)}chi_eta|0>
%                   =sum_nu v_eta,nu(u)|nu>,
% where |nu>=(a^dagger)^nu|0>/sqrt(nu!).
% Calculate them using the exact Wick polynomial
% [h3+(1/2)Sigma:partial partial h3](C zeta)chi_eta
%                         +h1_linear(C zeta)chi_eta,
% C_Aj=(S_Aqj+i S_Akj)/sqrt2, Sigma=S S^T/2.
% Multiply each zeta^nu coefficient by sqrt(nu!). The transpose in
% Sigma is the ordinary transpose, not a coefficientwise Hermitian
% conjugate: S(u) is real on the real u axis. This fixes every factor.
%
% The first input coefficient has only |nu|=1 components, defined by
% a1_eta(-pi)=eta sigma.(I-n_*n_*^T)(S(-pi)xi)_k chi_eta/6 |0>.
% Evaluate its creation part with the SAME C(-pi). The exact accepted
% normalizer is still used by the actual preparation; Z=1+O(epsilon)
% contributes only after this first sqrt(epsilon) coefficient. Its
% actual centering starts at the following relevant order as already
% retained in the exact input and response equations. Nothing here resets
% those quantities or omits their later contributions.
%
% EXACT FINITE METHOD. The source is a finite sum
% v_eta,nu(u)=sum_lambda v_eta,nu,lambda exp(lambda u),
% lambda=m+i n, m,n integers, |m|+|n|<=3. Use exact integer frequency
% labels and polynomial coefficients in the commuting zeta variables;
% there are13 Hermite indices,25 possible frequencies and2 spin entries
% per eta. Retain zero coefficients as certified zeros when available.
% The first amplitude is obtained coefficientwise, with no quadrature,
% a1_eta,nu(u)=a1_eta,nu(-pi)
%       -i sqrt3 sum_lambda v_eta,nu,lambda I_lambda(u),
% I_lambda=(exp(lambda u)-exp(-lambda pi))/lambda for lambda!=0,
% I_0=u+pi.
% Store a1 as a finite exponential sum plus its explicit linear-u term.
% This solves i partial_tau a1=V3(tau)a0 exactly for the specified finite
% coefficient, including its spin source and preparation term.
%
% REQUIRED OUTPUTS. Return both carrier arrays of v_lambda, initial a1,
% and the exponential/linear representation of a1; evaluate a1 at
% u=-pi,0,pi. Also return the exact finite-integral intervals
% int_0^(2pi sqrt3)||V3 a0||^2 d tau,
% int_0^(2pi sqrt3)A_(1/4)(V3 a0)d tau,
% int_0^(2pi sqrt3)||a1||^2 d tau,
% int_0^(2pi sqrt3)A_(1/4)(a1)d tau,
% separately for each eta, with A_rho=sum rho^(-|nu|)||coefficient_nu||^2.
% A certified uniform bound for this first coefficient follows directly:
% sup_tau sqrt[A_(1/4)(a1(tau))]
% <=sqrt[A_(1/4)(a1(u=-pi))]
%   +sqrt[(2pi sqrt3) int A_(1/4)(V3 a0)d tau].
% Report this bound without identifying it with the actual propagated
% state's amplitude tail. Its potentially large size is relevant input
% to selecting or rejecting a later truncation.
%
% The integrals require only quadratic products of the finite arrays.
% For zeta=m+i n define J_j(zeta)=int_-pi^pi u^j exp(zeta u)du, j=0,1,2.
% For zeta!=0, with E_+=exp(pi zeta), E_-=exp(-pi zeta), use
% J0=(E_+-E_-)/zeta,
% J1=[pi zeta(E_++E_-)-(E_+-E_-)]/zeta^2,
% J2=[pi^2 zeta^2(E_+-E_-)-2pi zeta(E_++E_-)
%                                      +2(E_+-E_-)]/zeta^3.
% At zeta=0 use J0=2pi,J1=0,J2=2pi^3/3 exactly. Multiply every time
% integral by sqrt3. For integer n, E_+=(-1)^n exp(m pi) and
% E_-=(-1)^n exp(-m pi), so no trigonometric interval library is needed.
% Conjugate frequencies and coefficients in squared norms, not merely
% their real parts before multiplication.
%
% ANALYTIC COMPONENT SCOPE. These are polynomial-Gaussian coefficient
% arrays. Their use for the actual fixed scalar cutoff includes a separate
% omitted-tail norm, rather than pretending that a cut state has only13
% Hermite modes. In whitened momentum variables, the cutoff differs from
% one only at |xi|>16; its derivatives through order3 have norms below2.
% Gaussian differentiation and the m<=3 tail integrals give the loose
% bound10^8 exp(-128) for every word of at most3 canonical operators
% applied to the omitted normalized scalar input. One can see this by
% bounding each unrotated word by64(1+|xi|)^3 times the Gaussian on that
% tail, and then allowing the norm-preserving U(3) Gaussian-coordinate
% rotation; its degree3 coefficient1-norm is at most6sqrt6.
% The normalization correction is smaller and is included in this bound.
% The scalar h3 coefficient1-norm is below60, and every row of S has
% coefficient1-norm below sqrt6*128<314; the linear spin source has
% coefficient1-norm below2. Therefore integrating this first source for
% at most2pi sqrt3<11, and adding the initial a1 cutoff discrepancy from
% its explicitly displayed degree1 preparation operator, gives an
% omitted-cutoff contribution below10^-35 to the first-amplitude L2 norm.
% This covers this finite coefficient only,
% not higher amplitude, launch, spacing or graph remainders. Keep the
% bound separate from the finite polynomial arrays in the saved output.
%
% ARITHMETIC AND ACCEPTANCE. Use directed intervals at100 decimal digits,
% with complex pairs of real intervals and exact integer frequency/multi-
% index labels. The same exact dyadic outward-decimal serialization as
% the accepted input producer is suitable. Each reported finite real/
% imaginary coefficient, endpoint value and integral must have raw width
% <=10^-25. Check the differential and initial identities coefficientwise,
% with residual intervals containing zero. Check S^T J S=J by its finite
% frequency coefficients; at u=0 S=I. The independently derived midpoint
% scalar and spin coefficients in actual RED MP34 and MP37 must agree
% with the directly generated Wick source, including the nonzero spin
% contribution221/864 for either equatorial carrier. Squared-norm
% integrals must be real, with imaginary interval containing zero and
% nonnegative real enclosure. State all parameters and the component
% scope with the output; no residual/root or physical coefficient passes
% merely because these finite arithmetic checks pass.
%
% RESOURCE BOUND. This is at most13 two-spin coefficient functions
% (26 scalar entries) per carrier, each with25 frequencies and one linear
% term, plus degree3 finite polynomial algebra. It needs no retained angular-mode array,
% large matrix exponential or propagator time steps. Use the existing
% strict180-second/512-MiB cap, one CPU, nice10 and no GPU. Reject/stop
% within that cap if the finite interval implementation exceeds it; there
% is no automatic refinement or larger run. The actual producer must be
% read against this specification by the same original author/RED pair
% before execution. No implementation or execution is authorized by this
% author's save alone; original-pair scientific review remains necessary.
% SOURCES: the independently saved full first rotating matrix recursion,
% the checked exact Wick representation and the selected S spectral
% projectors. This is the next finite coefficient of that same propagation
% construction, with explicit lower-band, native-response and final-root
% work still outstanding. No approximate orbit coefficient is its output.
% ACTUAL PRODUCER READ. Read all597 lines of
% declared_exterior_packet_amplitude.py against this repaired specification.
% The literal matrices/symbols, cubic single contractions, Hermite factorials,
% two eta preparations, frequency integration and initial/differential
% identities agree. The J0--2 quadratic products conjugate both coefficients
% and frequencies. Directed complex intervals, exact outward endpoint
% serialization, raw-width/independent-midpoint checks and resource caps
% preserve the stated arithmetic scope. No mandatory implementation repair
% was found. This clears only the specified finite component for the same
% pair's bounded execution; actual output acceptance remains separate.
%
% ACTUAL MP36--41 COMPARISON [R, same original pair].
% Read the complete actual RED blocks and their SOURCES. The complex-ball
% lower bound Re(k.k)>=kappa^2/6 and the homogeneous Cauchy remainder
% agree with the stated lower-degree INPUT construction; they do not bound
% the propagated degree. The free-frame response and one-position bounds
% retain the accepted exact normalizer, native centering and cutoff.
% Direct Weyl composition also checks MP38's second symbol:
% Q2=q2[[I,-I],[-I,I]], q2=[tr F0_qq-n^T F0_qq n]/(4p^2).
% Formal local block elimination adds |n cross F0_q|^2/(8p^3) to the
% upper diagonal. At the reference these terms are5/72 and1/72, giving
% Q2_eff,+ = I/12. The same transformation has the second coordinate
% term epsilon^2(F0_q-n(n.F0_q))/(4p^3); it cannot be omitted when that
% effective representation is used for native centroids. These are local
% continuum epsilon coefficients, not evaluated spacing responses.
%
% MP39's gauge identities and exact residual reconstruction agree:
% for w_tilde=z_tilde+T u_tilde, the original response residual is
% s_w=s_z+T r, with ||T r||<=||p^2 r||/24. Thus changing the response
% variable does not remove the base H^2 residual from the certificate.
% Its H^1-by-H^2 form insertion bound agrees with the actual kernel
% estimate, including the rounded second constant0.281 after division
% by epsilon at the selected value. This is not an L2 source estimate.
% The three bilinear error pairings in a dual insertion remain required.
% MP40 independently derives the same finite first-amplitude equation,
% and MP41 independently checks the actual repaired specification above,
% including the initial-cutoff term, full matrix V4 and100-digit finite
% arithmetic. The complete current repaired specification was self-read.
% No additional source/model premise or achieved propagation is imported.
% SOURCES: actual redteam_declared_exterior.tex MP36--41, full blocks and
% SOURCES; own Clifford/Weyl, residual and finite-coefficient checks.
%
% NEXT FULL MATRIX AMPLITUDE COEFFICIENT [R, independent derivation WIP].
% This identifies the actual use of the finite a1 component without
% claiming an error estimate for truncating at that order. Keep the same
% local adapted upper-band gauge as in the independently checked second
% symbol, the same physical rotating frame and the actual prescribed input.
% Write z_perp=z-n_* t, t=n_*.z, F1=A x+B z, F2=partial_x h3.
% Taylor expansion of
% h1_+=-sigma.(n(k) cross grad_q f)/(2p)
%                                  -(Omega/2)n_z(k)sigma.n(k)
% at p=3, grad_q f=0 gives the explicit quadratic spin polynomial
% h1_quadratic=-sigma.{ (n_* cross F2)/6
%              +[(z-2t n_*) cross F1]/18
%              +Omega z3(z-2t n_*)/18 }.
% In particular both the derivative of1/p and the derivative of n(k)
% are present; retaining only n_* cross partial_x h3 would be incomplete.
% After removal of the action and exact quadratic motion, the next
% upper interaction coefficient is
% V4(tau)=Op_W[h4(S_tau xi)]I
%             +Op_W[h1_quadratic(S_tau xi)]+I/12.
% The scalar I/12 is the full local second adapted symbol just checked,
% including the contribution from eliminating the first off-band term.
% The actual opposite-band input and its evolution remain separately in
% the reconstructed four-spinor state and physical coordinate; this
% diagonal coefficient is not a replacement for that evolution.
%
% For the uncut coefficient calculation, let Y=(S(-pi)xi)_k and
% p_c=3, P_perp=I-n_*n_*^T, V=(S(-pi)S(-pi)^T)_kk. All components of
% Y commute. The exact normalized free-frame preparation has expansion
% a(0)=chi_eta|0>+sqrt(epsilon)a1_eta(0)+epsilon a2_eta(0)+...,
% with the already specified a1 and
% a2_eta(0)={-eta sigma.[(n_*.Y)P_perp Y
%                              +n_*|P_perp Y|^2/2]/18
%                  +tr(P_perp V)I/144}chi_eta|0>.
% Indeed Z=1-epsilon tr(P_perp V)/72+O(epsilon^2), so the last term
% is the required normalization contribution. It yields the exact order
% identity2Re<a0,a2(0)>+||a1(0)||^2=0. In the same gauge the adapted
% upper input has no additional order-epsilon diagonal term, while its
% lower input starts with -epsilon D_-+ chi_eta|0>,
% D_-+=-sigma_z/(12sqrt3). The latter has NOT been set to zero.
% For the even scalar density the exact continuum centroid displacement
% starts at epsilon^2: grad s(k_c)=0, E[z s]=O(epsilon^2), and
% epsilon E[n_c cross z/p^2]=O(epsilon^2). After removal of the irrelevant
% constant translation phase, its amplitude starts at epsilon^(3/2).
% The accepted finite-epsilon normalizer/centroid remain the actual state
% data; these expansions do not redefine them. Cutoff and higher-order
% input errors still require their separately stated budgets.
%
% The second coefficient satisfies the ordered two-spin recursion
% i partial_tau a2=V3(tau)a1(tau)+V4(tau)a0.
% Exact Wick matrix elements of V3 acting on the stored a1 keep all
% annihilation and number-preserving terms. There are only the50 Hermite
% indices of degrees0,2,4,6 in this finite coefficient. Its functions are
% finite sums of u^j exp[(m+i n)u], |m|+|n|<=6, j<=2. The source uses
% j<=1, and the extra power is needed when its zero frequency is
% integrated. Matrix multiplication remains ordered; V3 at different
% times is not assumed to commute. These facts give a concrete next
% finite use of the first-amplitude arrays, but no next execution is
% specified here. Higher cut derivatives, omitted matrix symbols, both
% propagated bands, launch/native derivatives and an actual root bound
% remain necessary before a full finite-frequency result can be asserted.
% SOURCES: independent direct Taylor expansion of the actual h1 block,
% exact free-frame preparation and normalization, and ordered finite
% Weyl/Wick recursion. The I/12 local input is the explicitly compared
% MP38 second-symbol derivation, not a native orbit coefficient.
%
% ACTUAL FIRST-AMPLITUDE OUTPUT ACCEPTANCE [N, necessary component only].
% Read the actual result.json and final_status.txt in
% results/declared_exterior_packet_amplitude_20261015_1507/, following the
% complete actual597-line producer read. Inspected every parameter/scope
% field, all literal matrices and S coefficients, both carriers' complete
% source/amplitude arrays (including exact-zero frequency entries), every
% initial/linear/endpoint entry, all507 identity residuals, all integrals,
% uniform bounds and resource fields. Large arrays were inspected through
% compact numerical interval displays; the decisive endpoint and all
% integral/acceptance intervals were also read in their actual full decimal
% enclosures. Those displays are not substituted for the audited directed
% arithmetic or outward serialization. No scientific calculation was rerun.
%
% The accepted finite arrays use exactly the prescribed13 Hermite indices,
% 25 frequency labels,100-digit intervals and both eta carriers. The507
% residual intervals contain zero; the imaginary norm integrals contain
% zero and their real enclosures are positive. The maximum raw width is
% below2.807*10^-87, much smaller than10^-25. Both carriers give
% int||V3a0||^2 d tau =3660024150.23599545775...,
% int A_(1/4)(V3a0)d tau =102662411754.77817590868...,
% int||a1||^2 d tau =2424582900.15537379959...,
% int A_(1/4)(a1)d tau =64218275370.87003059388... .
% Their initial A_(1/4) is6.39646376919207962043... and the specified
% uniform sqrt(A_(1/4)) bound is1057004.94304156245727... .
% The nonzero spin midpoint norm is221/864 as independently required.
% The run reports0.631472 seconds and28800KiB peak RSS, within its fixed
% CPU/memory limits. The fixed-cutoff L2 allowance10^-35 remains separate;
% it is not a weighted, launch, spacing, graph or full-propagation bound.
% Arithmetic acceptance of these finite first-coefficient arrays is given.
%
% A CHECKED FAILURE OF THE FIRST TRUNCATION is a distinct consequence.
% At u=pi, eta=+, normalized Hermite mode(1,0,0), first spin component,
% the actual interval implies
% -28593.631689186481630879 < Im a1 < -28593.631689186481630877.
% This conservative displayed enclosure contains the full saved interval.
% Thus at the SELECTED epsilon=2^-24, epsilon||a1||^2>48, since
% 28500^2>48*2^24. All a1 modes are odd and a0 is the unit vacuum, so
% ||a0+sqrt(epsilon)a1||^2=1+epsilon||a1||^2>49.
% Its distance from every unit spinor is therefore greater than6 by the
% reverse triangle inequality. The same scalar-cutoff allowances leave
% this strict gap intact. This is the reference-circuit first unnormalized
% quadratic-frame truncation; it cannot provide a small state-error bound
% for the selected member. It does not exclude a reorganized propagation,
% determine the matched return root, or evaluate the absolute MP8 spacing
% coefficient. Merely writing the finite a2 identities above is not a
% proof that keeping another order is sufficient. The input, epsilon,
% reference, four returns and observable remain unchanged. No rerun or
% new execution is requested by this acceptance or by this limitation.
% SOURCES: actual approved producer, actual complete saved output fields
% and status, the independent Wick/spin specification, and the elementary
% parity/norm argument above. Same original RED actual output review is
% the paired comparison; no external coefficient agreement is inferred.
% FINAL SAME-PAIR READ: the complete actual RED MP42--44 and their SOURCES
% were read, including the last saved visible core. The actual producer
% and result retain their fully inspected approved hashes. MP43's input
% uses q_perp=tr(P_perp V)/2 and eta sigma.n chi_eta=chi_eta, giving the
% same normalization and quadratic spin terms as the independent a2
% calculation above. Its explicit common-phase removal is also retained.
% MP44 independently agrees with the finite-array acceptance and strict
% norm witness; neither identifies this first truncation with the exact
% propagated state or the matched return. No mandatory repair remains
% for this necessary-component handoff; the full MP8 result remains open.
%
% EXACT UNITARY EXTRACTION OF THE EXTREME CUBIC TERMS [R, independent WIP].
% The accepted first-amplitude failure requires reorganizing the SAME
% propagation, rather than changing epsilon or asserting that the next
% unnormalized order is sufficient. This calculation was saved before
% reading a new RED reorganization. It uses the fixed rotating reference,
% u=tau/sqrt3-pi and canonical whitened xi=(q1,q2,q3,p1,p2,p3),
% [q_i,p_j]=i delta_ij. These q,p are fluctuation operators, not the
% macroscopic coordinates or a new physical preparation.
%
% The horizontal K has eigenvalues +/-1/sqrt3 and a two-dimensional
% zero eigenspace. Direct multiplication gives the right eigenvectors
% v_s=(1,-sqrt2,0; s*3sqrt3-3sqrt6/2,-3sqrt3,0), s=+/-1.
% Their symplectic product is v_+^T J v_-=-6sqrt3. Write
% Q_+=(1/2+sqrt2/4)q1+q2/2+(p1-sqrt2 p2)/(6sqrt3),
% Q_-=(1/2-sqrt2/4)q1-q2/2-(p1-sqrt2 p2)/(6sqrt3).
% Then P_s xi=v_s Q_s for the exact rank-one spectral projectors,
% [Q_+,Q_-]=-i/(6sqrt3), and
% (S(u)xi)_q1=exp(u)Q_++exp(-u)Q_-.
% Substituting v_s into the actual cubic polynomial gives exactly
% c_s=h3(v_s)=1+s*33sqrt2/8.
% For example t=n_*.z=-6+s*3sqrt2, |z|^2-t^2=27/2+s*9sqrt2;
% the kinetic and flow contributions are respectively
% 3/2+s*3sqrt2/4 and -1/2+s*27sqrt2/8. This checks the sign and factor.
% Consequently the ACTUAL Weyl cubic has the exact decomposition
% Op_W[h3(S(u)xi)]
%   =c_+ exp(3u)Q_+^3+c_- exp(-3u)Q_-^3+R3(u).
% R3 is defined by this equality and retains every remaining frequency,
% including zero/vertical/mixed terms. The spin-linear h1 is also retained.
% No claim that R3 is small follows from removing the two extreme labels.
%
% Set kappa=sqrt(3epsilon). Extract the negative-frequency term on
% [-pi,0] and the positive-frequency term on[0,pi], without overlap:
% theta_-(u)=kappa*c_-[exp(3pi)-exp(-3u)]/3, -pi<=u<=0,
% theta_+(u)=kappa*c_+[exp(3u)-1]/3, 0<=u<=pi,
% W(u)=exp[-i theta_-(u)Q_-^3] on the first half,
% W(u)=exp[-i theta_+(u)Q_+^3]exp[-i theta_-(0)Q_-^3] thereafter.
% Each Q is self-adjoint on its standard linear-coordinate domain;
% these are exact unitary spectral functions and preserve Schwartz space.
% W is continuous at0. Its piecewise differentiability gives the exact
% transformed equation i b'=W^dagger(G-G_ext)W b for a=W b,
% whenever i a'=G a and i W'=G_ext W. No delta term appears at the join.
% This is an operator identity; neither the exact state nor its norm is
% replaced by the failed first Taylor expansion of W.
%
% Nonlinear conjugation MUST retain its quantum ordering term. For
% Q=c.xi, D_Q=c^T J partial_xi and U=exp(-i theta Q^3), the exact Weyl
% identity for any polynomial f is
% U^dagger Op_W(f)U
%  =Op_W{ [exp(theta D_Q^3/4)f](xi+3theta Jc Q^2) }.
% The exponential terminates at floor(deg(f)/3), and D_Q Q=0 makes its
% two differential generators commute. It follows directly from
% i[Q^3,Op_W(f)]=Op_W[-3Q^2 D_Q f+D_Q^3 f/4].
% In particular U^dagger p^3 U has Weyl symbol
% (p-3theta q^2)^3+3theta/2 when Q=q. Classical substitution alone
% misses that constant. The two factors of W are applied in their actual
% order, and ordered Pauli coefficients are unchanged by this scalar
% conjugation. Thus the retained cubic/spin and full V4 terms can be
% represented exactly by finite Weyl polynomials after this extraction.
%
% An explicit moment illustrates what the extraction does and does not
% control. In the unit vacuum put V_s=<Q_s^2>=47/144+s*sqrt2/8,
% a=theta_-(0)/(2sqrt3), b=theta_+(pi)/(2sqrt3),
% A=Q_+-a Q_-^2. At the final reference time,
% W^dagger Q_+ W=A,
% W^dagger Q_- W=Q_-+b A^2,
% <A>=-a V_-, <A^2>=V_++3a^2 V_-^2.
% Hence the radial Cartesian fluctuation of this unitary control is
% sqrt(epsilon){-exp(pi)a V_-+exp(-pi)b[V_++3a^2 V_-^2]}.
% Odd Gaussian moments make the cross term in <A^2> vanish. This is an
% exact control moment for W|0>, not the selected physical centroid or
% the matched MP8 radius. It cannot replace the exact prepared input,
% frequency adjustment or retained R3/spin/full-matrix evolution.
%
% For the actual four-spinor certificate this is only a CHANGE OF
% REPRESENTATION. On the already specified common localized free-band
% frame, extend W as diag(W I2,I2), retain the opposite-band input and
% all off-diagonal couplings, and keep every frame/localization remainder.
% Use the same reference W for all launch labels and for the spacing
% derivative: there is no artificial W_delta source. Both physical
% observables and the actual initial and driven states are conjugated.
% Time derivatives in the return conditions still differentiate W and
% its transformed observables. The physical mean-frequency row retains
% its actual Hamiltonian, not the extracted interaction generator.
% Exact unitary conjugation preserves residual norms in the appropriate
% representation; physical weighted/graph norms require the displayed
% polynomial shear factors. No Gaussian tail is silently inherited after
% a cubic phase, and no propagated degree or actual root follows yet.
% The next missing step is a quantitative residual/tail/launch certificate
% for the RETAINED full equation, not another unchecked finite-order sum.
% No new scientific execution or substitute detector is specified here.
% SOURCES: independent eigenspace substitution into the saved literal K
% and h3, exact one-coordinate spectral calculus, Weyl commutator and
% Gaussian moment identities. No new counterpart reorganization was read
% before this save. The accepted actual amplitude output motivates the
% need for this representation but supplies none of these operator steps.
%
% ACTUAL MP45 COMPARISON [R, residual component only]. After the preceding
% independent save, read the entire actual RED MP45 and its SOURCES.
% Its projectors, c_s, gamma=1/(6sqrt3), quadrature norms and ordered
% conjugations agree. Its candidate is DIFFERENT from the piecewise W:
% U_app(u)=U_+(u)U_-(u) with both factors evolving on[-pi,pi],
% b_+=kappa c_+[exp(3u)-exp(-3pi)]/3 and
% b_-=kappa c_-[exp(3pi)-exp(-3u)]/3.
% Direct differentiation gives its extreme-term residual
% [U_+ H_- U_+^dagger-H_-]U_app|0>, H_-=kappa c_-exp(-3u)Q_-^3.
% Conjugating first by U_+, then by U_-, gives exactly the ordered
% polynomial norm used in MP45. Its three d-power integrals follow by
% t=exp(3u), d=gamma kappa c_+(t-exp(-3pi)); they agree term by term.
% With I1<4e-6,I2<8e-7,I3<1.3e-6, the Wick word bounds give
% S1<20,S2<800,S3<16500, hence the integrated component is below0.023.
% I independently checked these loose bounds using the displayed
% c4,...,c12 margins, ||l_+||<1.01, ||l_-||<0.55 and |B|<2.62.
% This bound belongs only to THAT continuously ordered U_app and this
% one uncut-vacuum residual component. It is neither a bound for the
% piecewise W above nor for the remaining R3, spin, full V4, exact cut
% preparation, opposite band, native spacing or launch/root response.
% The finite ordered products correctly include the quantum term that
% appears explicitly in the Weyl formula above. No new execution or
% achieved full-state/observable error is inferred. The positive role of
% the calculation is to replace an unbounded first-truncation error by
% an exact unitary candidate with one explicitly controlled contribution.
% SOURCES: actual redteam_declared_exterior.tex MP45, complete derivation
% and SOURCES, versus the independent rank-one/Weyl derivation above and
% direct elementary integral and Gaussian-word checks.
%
% FULL-BAND CENTROID RESPONSE BOUND [R, independent certificate component].
% This estimate addresses the actual matrix readout, not an upper scalar
% substitute. Keep the exact native-projector gauge already derived:
% z=w-Tu, T=i Sigma.(k cross k^3)/(12p^2),
% q_delta=2Re<u,q z>+<u,[q,T]u>. Here q=i epsilon partial_k and p=|k|.
% Work first on the common localized momentum tube p>=kappa0>0, retaining
% its explicit localization errors separately. The exact free-band unitary
% has q in its coefficient representation equal to q I+epsilon A, with
% A=(1/2)[[-C,C],[C,-C]], C_i=sigma.(n cross e_i)/p.
% In the Euclidean vector readout norm, every A block obeys
% ||A_block psi||<=||psi||/(sqrt2 kappa0). This follows from
% sum_i C_i^dagger C_i=2I/p^2, not from treating the components as scalars.
%
% Write u=(u_+,u_-), z=(z_+,z_-), ||u||=1, and let
% eta=||u_-||, eta_x=||(q_i u_-)_(i=1)^3||,
% Z_+=||z_+||, Z_-=||z_-||.
% The difference between2Re<u,(q I+epsilon A)z> and its COMPLETE
% upper block2Re<u_+,(q I+epsilon A_++)z_+> is bounded by
% 2 eta_x Z_-+(sqrt2 epsilon/kappa0)
%                         [Z_-+eta(Z_-+Z_+)].
% The first term is the lower diagonal position contribution; the other
% three are the lower Berry block and both off-diagonal coordinate blocks.
% Thus a small occupation eta alone never controls a spacing response:
% its Z_- and position-weighted eta_x must also be bounded. In particular
% eta=O(epsilon) with Z_-=O(1) contributes at order epsilon, consistently
% with the earlier differentiated-interband warning. No cancellation of
% that term has been assumed for the reflection mixture.
%
% The direct native-gauge readout term has an explicit global form bound.
% Put s=k cross k^3. Then |s|<=p^4/2 and
% ||grad s||_F<= (sqrt3+3)p^3, since
% partial_i s=e_i cross k^3+3k_i^2(k cross e_i).
% Hence ||grad(s/p^2)||_F<6p and
% ||([q_i,T]psi)_i|| <= (epsilon/2)||p psi||.
% The inequality includes the spin matrices, whose real vector contraction
% has squared norm equal to the vector's Euclidean square. On a further
% p<=P tube, retaining only the upper block of this term leaves error
% at most epsilon P(eta+eta^2/2). Alternatively keep actual momentum
% norms instead of the upper cutoff; the homogeneous derivative bound
% extends through p=0 as a first-order form and does not need kappa0.
% The full native centroid response therefore retains the upper direct
% commutator expectation and all the explicit lower-band bounds above.
%
% For an actual approximate two-band solution, apply the same block
% decomposition to its errors as well as to its amplitudes; bound each
% bilinear difference with one exact-minus-approximate factor. This is
% compatible with the previously derived full-matrix dual residual identity.
% The exact cubic unitaries are representation changes inside these blocks;
% physical eta_x and graph norms must include their polynomial shears.
% A norm estimate for their uncut-vacuum ordering component supplies none
% of eta,eta_x,Z_+ or Z_- for the actual coupled matrix evolution.
%
% The same coordinate inequality applies to an initial-launch tangent,
% with z replaced by that tangent and no native T term. Root columns in
% R,p_r,L therefore require the corresponding differentiated amplitudes.
% The time column remains the actual native velocity at the readout times.
% The mean-frequency row can be evaluated at preparation: the background
% is fixed and static, so its exact full native mean energy is conserved.
% It must use that full initial Hamiltonian and prepared covariance,
% including normalization/centering; no extracted interaction generator
% replaces it. The time-mean norm of the centroid is1-Lipschitz in the
% centroid path before the already stated return-root correction. These
% observations give a quantitative FULL-MATRIX route for the specified
% readout, but no values for the required propagated norms or actual root
% are asserted. They are the remaining certificate inputs, not new state,
% frequency, detector or execution choices.
% SOURCES: independent exact free-band coordinate connection, native gauge
% identity already saved, Cauchy--Schwarz in the vector readout space and
% direct homogeneous-symbol differentiation. No new RED block on this
% estimate was read before saving it; no scientific run was used.
%
% GLOBAL ANGULAR CONTROL FOR THE BAND COORDINATE [R, independent WIP].
% The inverse-momentum coordinate connection need not rely only on a
% presumed momentum tube. This is an evaluated lower-angular-tail bound
% for the ACTUAL selected input and a small explicit launch neighborhood;
% it is not a bound on its upper angular dimension or on the response tail.
% Use physical scaled Jz=q1 k2-q2 k1+epsilon Sigma_z/2. Before P0 and
% the scalar cutoff, the translated squeezed Gaussian is metaplectic,
% with center q_c=R e1-mu and carrier k_c. In its unit-vacuum coordinates,
% Jz=L_c+sqrt(epsilon) l.S xi+epsilon Q_S+epsilon Sigma_z/2,
% L_c=(q_c cross k_c)_z,
% l=(k_c2,-k_c1,0;-q_c2,q_c1,0),
% Q_S=(S xi)_q1(S xi)_p2-(S xi)_q2(S xi)_p1.
% The displayed products commute within each term and are Weyl ordered.
% At the reference l.S=l and ||l||=sqrt10, as follows from l.K=0.
%
% Keep |R-1|<=0.01, |k_c-k_*|<=0.01 and the already proved
% |mu|<9e-7. Then L_c>1.7 and
% ||l.S||<=sqrt10+128 sqrt(0.01^2+0.010001^2)<5.
% On the Hermite space of total degree<=n a real quadrature has norm
% at most ||l||sqrt(2(n+1)) into degree<=n+1. Each ordered product
% of two quadratures has norm at most2(n+2) times their row norms.
% The sum of the two row-norm products in Q_S is less than
% ||S||_F^2/2<8192. Consequently on degree<=n,
% ||(Jz-L_c)psi|| <= [5sqrt(2(n+1)epsilon)
%                    +16384(n+2)epsilon+epsilon/2]||psi||.
% Each application raises degree by at most2. At epsilon=2^-24 and
% n<=14 the bracket is below0.023. Iterating8 times from the vacuum
% therefore gives ||(Jz-L_c)^8 g chi||<(0.023)^8, without a code run.
% The spin term is retained; chi need not be a Sigma_z eigenvector.
% By the spectral theorem,
% ||1_{Jz<0.8}g chi|| <(0.023/0.9)^8<2.1e-13.
%
% P0 commutes EXACTLY with Jz, so the same bound holds after projection,
% divided by the actual sqrt(Z). Uniformly on this neighborhood,
% 1-n(k).n_c<=2|k-k_c|^2/|k_c|^2 implies
% Z>=1-epsilon tr(V)/(2|k_c|^2)>0.9999 for the uncut Gaussian;
% the fixed scalar-cut correction is smaller than10^3 exp(-128) in
% normalized L2 and does not affect that margin. Apply the angular
% spectral projector's contraction property to the cut/uncut difference:
% no unproved eighth derivative of the cutoff is needed for this L2 step.
% The exact normalized, centered preparation thus has
% ||1_{Jz<0.8}phi_0||<3e-13, for both reflected members. The reflection
% preserves Jz. Central continuum H commutes with Jz, so this actual
% low-angular BASE-state norm bound persists for every propagation time.
% The finite-a Hamiltonian does not have that symmetry; this theorem is
% for the continuum base whose native spacing derivative is being computed.
%
% On a free-band coefficient state whose total Jz is at least j_min,
% each scalar spin component has orbital magnetic number at least
% j_min/epsilon-1/2. The cylindrical angular kinetic form therefore gives
% p^2 >= (j_min-epsilon/2)^2/r^2.
% Inverting positive forms yields the useful GLOBAL inequality
% ||p^(-1)psi|| <= ||r psi||/(j_min-epsilon/2).
% It follows first on smooth finite-angular compact radial cores and then
% by form closure; it is not a global spectral gap of the massless Dirac
% operator. The free-band unitary intertwines physical rotations, so its
% coefficient Jz is exactly orbital Jz+epsilon sigma_z/2 in each band.
% This justifies applying the form bound to the same band-coordinate map.
% Here r=|q| is the COEFFICIENT-position radius; the physical full-spinor
% coordinate in this frame is q I+epsilon A, which remains in the readout.
%
% In particular the preceding response inequality can replace
% (sqrt2 epsilon/kappa0)[Z_-+eta(Z_-+Z_+)] by
% (sqrt2 epsilon/L_min)[Zr_-+eta(Zr_-+Zr_+)],
% Zr_+=||r z_+||, Zr_-=||r z_-||,
% whenever all retained state/response components have
% Jz>=L_min+epsilon/2. This keeps the full matrix coordinate without
% assuming their momenta stay in a chosen compact tube.
% The native B and T have angular ranks0 and4, so one exact spacing
% insertion shifts Jz by at most4epsilon. On the retained base sector
% above, L_min=0.79 is a safe common lower value. A launch derivative
% needs its own prepared tangent/tail bound. Likewise the discarded
% 3e-13 base tail cannot simply be multiplied by an UNBOUNDED B: its
% graph-weighted input and propagated response tails must still be proved.
% No such response bound, upper-j cutoff, dimension, root or coefficient
% is asserted here. The new result removes one unnecessary momentum-
% localization premise from the exact coordinate estimate and supplies
% an actual persistent base-sector bound for the fixed preparation.
% SOURCES: independent metaplectic polynomial action, spectral Markov
% inequality, exact central rotational symmetry and cylindrical kinetic
% quadratic forms. The existing exact preparation/centering and ||S||_F
% bounds are used with their stated domains. No new RED angular result
% was read before this save, and no scientific calculation was executed.
%
% GRAPH-WEIGHTED LOW-ANGULAR RESPONSE TAIL [R, independent WIP].
% This upgrades the preceding ACTUAL input estimate to the graph norm
% needed by the full native B. It selects no upper angular truncation and
% changes neither preparation nor detector. Use the same launch box,
% exact normalization/centering, Jz and L_c, and let A=Jz-L_c.
% On total Hermite degree n<=62 the preceding bracket is less than0.078:
% 5sqrt(126)/4096+64/1024+2^-25<0.078.
% Applying A32 to the uncut vacuum thus produces a spinor polynomial of
% degree<=64 with norm less than(0.078)^32. This is the 32nd power of A,
% not a new state or a different semiclassical parameter.
%
% Let T_j=||p^j psi||. On any translated/metaplectic polynomial of degree
% <=64, the momentum row bound gives, for j<=3,
% T_j<=4^j||psi||: each multiplication has norm at most
% |k_c|+sqrt(epsilon)||S_k||_F sqrt(2(64+j))<4.
% P0 commutes with p and is a contraction, so
% T_j(P0 psi)<=4^j||psi||, with the UNPROJECTED polynomial norm on the
% right, without differentiating its momentum symbol. Write b=div(v)/2 and
% use b0<=52.5, V1^F<=105, V2^F<=24300sqrt3,
% b1^F<=12150sqrt3, b2^F<=17100000. Direct ordered differentiation gives
% ||H psi|| <=2T1+epsilon b0 T0,
% ||H^2 psi||<=4T2+420epsilon T1+45000epsilon^2 T0,
% ||H^3 psi||<=8T3+1890epsilon T2
%                           +500000epsilon^2 T1+80000000epsilon^3 T0.
% For example T_s(Hpsi)<=2T_(s+1)+epsilon b0 T_s
% +sum_(j=1)^s binom(s,j)[epsilon^j Vj^F T_(s-j+1)
%                         +epsilon^(j+1) bj^F T_(s-j)]
% proves these three bounds by composition. They imply
% E3(P0 psi)<600||psi|| on the unprojected polynomial input just specified.
% This is an UPPER
% bound on H powers; the earlier elliptic K_j bounds had the reverse use.
%
% H and P0 both commute strongly with Jz (P0 need not commute with H).
% The spectral theorem on Pi_low=1_(Jz<0.8), whose distance from L_c is
% greater than0.9, therefore gives for the normalized uncut input
% E3(Pi_low phi_uncut)<600(0.078/0.9)^32/sqrt(Z)<7e-32.
% This uses H^j A32 P0 g=H^j P0 A32 g and strong commutation of H with
% the angular spectral multiplier; no bound on H applied to a bare L2
% tail is being assumed.
%
% The actual cutoff can be restored in this GRAPH estimate directly in
% momentum space. With y=V^(-1/2)u, its removed tail has |y|>16;
% there p<=|k_c|+|y|/32<=|y|/4. For j<=3,
% ||p^j 1_(|y|>16)g||^2
% <=4^(-2j) Gamma(j+3/2,256)/Gamma(3/2)<90000 exp(-256).
% Normalizing the projected cut and uncut states, with both projected
% norms greater than sqrt(0.9999), consequently gives
% T_j(phi_cut-phi_uncut)<1000 exp(-128), j<=3.
% The same fixed actual translation q_c is used in this comparison.
% The three displayed H-power estimates imply
% E3(phi_cut-phi_uncut)<20000 exp(-128)<6e-52.
% There is no cutoff-derivative assumption in this step: T_j are momentum
% multipliers, and the coefficient derivatives of H^j were already bounded.
% Thus, conservatively, for the EXACT selected input and stated launch box,
%                         E3(Pi_low phi_0)<1e-31.
% Strong commutation and static unitary propagation preserve this bound
% for every time. The smooth globally subcritical core is essential to
% the graph constants; it has not been replaced by a global exterior.
%
% The actual native initial derivative is
% w0=(P2+i m.k)phi_0, m_i=<phi_0,M_i phi_0>.
% The previously derived |k cross k^3| bound gives ||P2(k)||<=p^2/24;
% the component formula for M gives |m|<1/2 on the actual cut support.
% Define an APPROXIMATION to the first spacing-response equations by
% replacing phi_0 with Pi_high phi_0, UNNORMALIZED, and using the same full
% coefficient m in its initial derivative. The omitted initial response
% is then exactly (P2+i m.k)Pi_low phi_0, with norm bounded by
% (K2/24+K1/2)E3(Pi_low phi_0)<2e-30.
% This retains one spacing insertion, not finite angular dimension;
% Pi_high still contains infinitely many angular labels.
% P2 has rank4 under simultaneous rotations, while m.k has rank1, so the
% retained response has Jz>=0.8-4epsilon>0.79+epsilon/2. This justifies the
% preceding global band-coordinate inverse-momentum bound on that part.
%
% For the DISCARDED full four-spinor response, exact Duhamel and the
% previously established ||B psi||<1180 E3(psi) now give, for 0<=tau<=12,
% ||w_discarded(tau)||
% <=2e-30+(12/epsilon)1180(1e-31)<3e-20.
% This bound retains the endpoint-ordering term, both bands and the actual
% native projector/centering derivative. It is a calculated persistent
% L2 response-tail bound for the specified low-angular omission, not an
% evaluated retained propagation or a smallness claim based only on its
% initial occupation. A coordinate-weighted BASE-tail bound, launch-label
% graph tails, an upper-j/dimension choice, the actual retained solution
% and its return-root certificate remain to be supplied. No run follows
% from this component alone.
% SOURCES: independent ordered-tensor differential estimates and the
% 32nd angular spectral moment, using the earlier checked selected-core
% constants. No new RED response-tail derivation was read before this
% save. The original-pair comparison and all remaining errors are pending.
%
% ACTUAL MP46--50 COMPARISON [R, same original pair].
% Read the complete current RED MP46--50 and their SOURCES after saving
% the independent full-band/base-angular/graph-tail routes above. Direct
% substitution into h3 gives the RED's complete horizontal polynomial
% X^3+(45sqrt2/8)X^2D-(3sqrt2/2)D^3
% +Y[(9sqrt2/2)(X^2+D^2)-9XD]
% -P[(sqrt3/2)(X^2+D^2)+(sqrt6/2)XD]
% +Y^2[-3X/2-(9sqrt2/4)D]+sqrt3 DYP,
% with X=e^u Q_++e^-u Q_-, D=e^u Q_+-e^-u Q_-, Y=Q0,P=P0.
% This checks its c_sigma,d_sigma,L_sigma,M_sigma. The remaining vertical
% terms are exactly sum_sigma e^(sigma u)Q_sigma V_sigma(u)
% +(sqrt3 P-3sqrt2 Y)(9Z^2-C^2)/18; the saved h1 third component reduces
% to sqrt3 X+(sqrt6/2)D. Thus the no-constant-frequency assertion and the
% stated finite frequencies are operator identities after Weyl ordering.
% This absence alone supplies no small normal-form error.
%
% The full scalar sign-group control in MP47 includes q^2p transport and
% conditional transverse quadratic evolutions. Its characteristic sign,
% half-density factor and denominator>0.728 on |q_in|<=16 check exactly.
% A global saturation outside |q|<=32 is only an approximation-space
% control; its actual tail difference and cross-group residual remain.
% Neither the original Gaussian tail nor MP45's0.023 component bound is
% transferred to an ordered product containing this enlarged control.
% For MP48 the conditional vertical Hessian is
% [[3sqrt3+delta lambda,-sqrt2],[-sqrt2,1/sqrt3-delta lambda/9]].
% Its determinant1-delta^2 lambda^2/9 and the final matrix
% cos(2pi omega_lambda)I+sin(2pi omega_lambda)M_lambda/omega_lambda
% check, including the sign. The norm<0.00037 on |lambda|<=32 is a
% conditional QUADRATURE-matrix defect, not a full state or root bound.
%
% MP49 independently checks the complete full-band coordinate and inverse
% angular-form estimates. Its coefficient-radius precision is now explicit
% above. MP50 was saved before reading the actual graph-tail block; its
% independent32nd-moment route uses R_tail=4e-34 for degree<=4 inputs.
% I checked its bracket0.081, scalar p^2/p^4/r^2 norms, exact angular
% commutations and momentum Hardy inequality. In particular
% ||q P0 Pi g||<=(1+sqrt2)||q Pi g||,
% since ||(partial_i P0 psi)_i||=||p^-1 psi||/sqrt2 and
% ||p^-1 psi||<=2||q psi||/epsilon in three dimensions.
% This is physical full-spinor q before the band-frame change. It yields
% ||q Pi phi_0||<5R_tail and, by the bounded vector velocity,
% ||q exp(-itH/epsilon)Pi phi_0||<77R_tail for0<=t<=12.
% Its independent loose graph constant832R_tail and response<9e-20
% agree with the sharper independent bound above. Exact bilinear expansion
% with ||q u||<38 and ||w||<4e14 then gives a raw fixed-time centroid-
% response omission below1e-16, with base centroid error below1e-31.
% All actual centering coefficients are retained and neither retained
% component is renormalized. These are bounds for the low-angular error
% contribution only, not for the retained propagation or matched root.
% The independent RED has also read the actual graph-tail source and
% checked its420/45000,1890/500000/80000000,600 and cutoff constants;
% no formula repair was requested. Both methods remain within one pair[R].
% SOURCES: full actual redteam_declared_exterior.tex MP46--50/SOURCES;
% direct hand polynomial expansion, canonical transport,2-by-2 quadratic
% evolution, Hardy and ordered energy-graph inequalities as displayed.
%
% LAUNCH-DERIVATIVE LOW-ANGULAR CONTROL [R, independent continuation].
% The actual fixed threshold Pi=1_(Jz<0.8) is independent of launch
% labels. Its derivative is therefore ZERO; the label-dependent L_c
% only centers the moment estimate. Use coordinates (R,k_cx,k_cy) on
% the same explicit launch box. The final MP8 labels (R,p_r/omega,
% L/(omega r_*)) are recovered by k_cy=[L/(omega r_*)]/R.
% The covariance and scalar cutoff are held fixed exactly as prescribed.
%
% Before normalization and P0 the input is
% f=e^(-i q_c.k/epsilon) g(k-k_c) chi(k_c), q_c=R e1-mu(k_c).
% P0 is label independent. For the uncut comparison each first launch
% derivative, including the derivative of its actual projected norm,
% is P0 times a degree<=1 polynomial on this same translated Gaussian.
% The common actual q_c and D mu are retained in the cut/uncut comparison.
% At the equatorial carrier the fixed spinor gauge has ||D chi||<1.
% The exact Gaussian identity
% ||A_gauss V^(1/2)||^2=||G_qq||<39^2
% gives ||partial_(k_cj)g||<39/(sqrt(2epsilon))<1.2e5.
% The phase contribution is bounded by
% ||D mu|| ||p g||/epsilon<4e-6(3.1)/epsilon<209.
% For R the raw tangent is -i k1 f/epsilon, with norm<3.1/epsilon.
% Differentiating the positive projected normalizer costs at most another
% copy of each raw norm divided by its norm>sqrt(0.9999). Consequently
% every normalized uncut tangent is P0 f_j with degree(f_j)<=1 and
% ||f_j||<2e8. This loose common bound deliberately retains any harmless
% global phase derivative; no phase convention changes the preparation.
%
% For completeness the first differentiated actual-cut comparison is
% also quantitative. On |y|>16, p<=|y|/4 and the Gaussian derivative
% satisfies |partial_k g|<=1.6e5 |y| |g|. The first two cutoff derivatives
% obey ||D_u theta||<=8/2048 and
% ||D_u^2 theta||<136/2048^2. The label phase coefficients are below1e8;
% after one physical q derivative, use epsilon||A_gauss||<=2^-9 and
% the same two cutoff derivatives. Direct Gaussian radial moments up to
% Gamma(11/2,256), with normalization differentiated as well, give the
% deliberately loose simultaneous bounds
% T_s(partial_j phi_cut-partial_j phi_uncut)<1e16 exp(-128), s<=3,
% ||q(partial_j phi_cut-partial_j phi_uncut)||<1e16 exp(-128).
% One way to check this margin is to bound a raw tail tangent by
% 1e7(1+|y|)|g|, its one-q derivative before P0 by
% 2e7(1+|y|^2)|g|, and p^s by(|y|/4)^s. Projection's one-position
% derivative is bounded by the same Hardy factor1+sqrt2 used in MP50.
% The projected normalization and its derivative add less than a factor10
% to these tail norms. All are far below the displayed1e16 margin.
% No high angular derivative or additional cutoff smoothness is assumed.
%
% Apply the MP50 degree<=4 moment estimate with R_tail=4e-34 to f_j
% and r^2 f_j. For E3 apply D32 to the degree-one f_j FIRST, then the
% T_s<=4^s bound to its degree<=65 polynomial. Odd |p|^s f_j need not
% be polynomial. The r^2 f_j input has degree<=3 and norm below2||f_j||
% by the position-row bound. The independent H-power estimates give
% for e_j(0)=Pi partial_j phi_0
% ||e_j(0)||<1e-25, E3(e_j(0))<5e-23, ||q e_j(0)||<5e-25.
% The common fixed H is independent of the launch labels, so e_j(t)
% propagates homogeneously. Hence ||e_j(t)||<1e-25 and
% ||q e_j(t)||<5e-24 for0<=t<=12. These are actual tangent-tail bounds,
% not derivatives of a parameter-dependent angular projection.
%
% Combining them with ||q e||<77R_tail, ||partial_j u||<2e8 and
% ||q u_h||<39 gives the raw full-spinor centroid-column error
% |2Re<partial_j u,q u>-2Re<partial_j u_h,q u_h>|<3e-23.
% The bounded-velocity time column has error below1e-32. Initial-energy
% and its first launch derivatives are controlled by the same graph
% bounds; a common bound1e-21 is sufficient for their omission. The
% coordinate change to the stated dimensionless MP8 launch labels has
% norm below4 on this box. Passing from centroid coordinates to radius
% and its derivatives STILL needs the actual nonzero-centroid margin;
% if |<q>|>=1/2 throughout, these tail contributions to the radius
% columns are below1e-20. Explicitly MP50's bilinear estimate gives base
% centroid error<6.3e-32; the normalization map has Lipschitz constant
% below2.01 on the joining segment. The final column bound is at most
% 4[3e-23+2.01(6.3e-32)(2*39*2e8)]<1e-20.
% That radius margin and the retained inverse Jacobian
% are not presumed proved by the angular analysis.
%
% Thus the selected low-angular omission is controlled in the fixed-time
% native response AND in the baseline launch/time/energy rows needed by
% the four-return protocol. The retained full-matrix propagation, its
% upper angular size, an actual nonzero-centroid/return-root certificate
% and the finite-frequency coefficient remain open. No mixed spacing-
% launch derivative is claimed, and no execution specification follows.
% SOURCES: independent differentiation of the exact prepared input after
% the actual MP50 comparison, using its already checked degree4 angular
% bound and Hardy estimate. The normalized tangent, actual centering and
% fixed projection were kept explicitly; no new RED launch calculation
% was read before saving this extension. The actual RED read checked the
% constants and requested the D32-first and explicit radius-map precisions;
% those two proof details are now retained above.
%
% EXACT SPIN FACTORS FOR THE SAME CUBIC REPRESENTATION [R, independent WIP].
% This extends the checked scalar control representation, not the model or
% the state. Write zeta=sqrt(3epsilon), to distinguish this coefficient
% from the native spacing-response parameter delta=(a omega)^2. The first
% interaction generator in u is zeta[h3(S(u)xi)I+h1_linear(S(u)xi)].
% All following factors act in the upper two-spin amplitude; embedding
% them in the full four-spinor representation retains and transforms the
% opposite-band Hamiltonian/couplings, native response and physical readout.
% No full-matrix residual is removed by the scalar or spin identities alone.
%
% The horizontal spin term belonging to the scalar sign group sigma is
% -zeta b_sigma exp(sigma u)Q_sigma sigma_z/6,
% b_sigma=sqrt3+sigma sqrt6/2. It can be included EXACTLY in MP47's
% conditional transport. With q(u)=q_in/D_sigma(u),
% D_sigma=1+gamma zeta d_sigma[exp(sigma u)-exp(-sigma pi)]q_in,
% its spin phase along the same characteristic is
% Theta_sigma=-b_sigma log(D_sigma)/(6sigma gamma d_sigma).
% Indeed partial_u log D_sigma=sigma gamma zeta d_sigma exp(sigma u)q(u).
% The factor exp(-i Theta_sigma sigma_z) therefore solves the spin part
% along that characteristic and commutes with its scalar and transverse
% quadratic evolutions. On |q_in|<=16, the checked denominators give
% |Theta_+|<0.096 and |Theta_-|<0.016 throughout the reference circuit.
% Outside this controlled characteristic set, any transport saturation
% must retain its corresponding spin difference in the actual residual.
% Products of opposite sign groups still have their full noncommuting
% residual; no error estimate from the scalar-only product is transferred.
%
% The remaining vertical spin term also has an exact operator resolution.
% Define the canonical pair
% Q=3^(1/4)q3, P=(p3-sqrt6 q3)/3^(1/4), [Q,P]=i,
% a=(Q+iP)/sqrt2, N=a^dagger a. The vertical quadratic generator is
% H_v=(P^2+Q^2)/2=N+1/2,
% while K_v=-3^(1/4)(Q sigma_x+P sigma_y)/6.
% The time-dependent vertical spin term is
% zeta exp(iu H_v)K_v exp(-iu H_v), exactly the saved h1_linear vertical
% block. Its propagator from -pi is
% W_spin(u)=exp(iu H_v)exp[-i(u+pi)(H_v+zeta K_v)]exp(i pi H_v).
% The middle Hamiltonian is self-adjoint on D(N), since a and a^dagger
% are relatively bounded with arbitrarily small bound. It has the exact
% invariant pairs {|n,down>,|n+1,up>}, n>=0, and uncoupled |0,up>.
% On a pair its matrix is
% [[n+1/2,-g sqrt(n+1)],[-g sqrt(n+1),n+3/2]],
% g=zeta 3^(1/4)/(3sqrt2), with eigenvalues
% n+1 +/- sqrt[1/4+g^2(n+1)]. Thus no Fock-space truncation is needed
% to define this factor or its exact action on arbitrary input coefficients.
%
% At u=pi the free factors conjugate each pair matrix by i tau_z.
% Since both free eigenphases over2pi are -1, spectral calculus gives
% ||(W_spin(pi)-I)psi||<=2pi g^2||(N+1)psi||.
% This is valid for arbitrary spin input in D(N), not merely for an
% unpolarized expectation. The zero-mode uncoupled block contributes zero.
% For the uncut reference (q3,p3) vacuum, which is a squeezed vacuum for
% (Q,P), n_bar=5/(2sqrt3)-1/2 and
% ||(N+1)psi||=sqrt(3 n_bar^2+4 n_bar+1)<2.73.
% With g^2=epsilon sqrt3/6 this yields a control-factor state defect
% below3e-7 after the reference circuit. This evaluated number belongs
% only to this factor on that stated reference vacuum. The actual
% projected/cut input and a preceding factor need their own (N+1) norm;
% neither is identified with the reference vacuum.
%
% The conditional scalar vertical term in MP48 can be retained alongside
% this spin coupling without discarding either. In the spectral fiber
% T0=lambda the exact middle operator becomes
% H_v+zeta lambda(9q3^2-p3^2)/18+zeta K_v.
% For |zeta lambda|<3 its quadratic part is positive, with determinant
% 1-zeta^2 lambda^2/9, and the linear spin term is relatively bounded on
% its oscillator domain. The resulting factor is a well-defined unitary
% oscillator-spin evolution. The simple two-state blocks above no longer
% diagonalize it when lambda!=0; its counter-rotating terms and their
% actual coefficient errors cannot be dropped. This is a constructive
% exact representation, not an evaluated joint propagation or a new run
% specification. Cross-group terms, full V4/higher symbols, both bands,
% the native response, actual propagated weights and matched root remain.
% SOURCES: independent canonical shear/rescaling, differentiated
% characteristic logarithm and exact two-state oscillator-spin matrices,
% using the already compared complete h1 in RED MP47. No new RED spin
% factor or scientific execution was used before this save.
%
% ACTUAL MP51--53 AND CONDITIONAL VERTICAL GRAPH CONTROL [R, WIP].
% Read the full actual original RED MP51--53 and SOURCES. MP51 confirms
% the independently saved graph-tail constants; MP52's two launch proof
% precisions are included above. In MP53, D'=sigma(D-h) and
% f=(D-h)/(gamma d D) give int f=log D/(sigma gamma d). Direct matrix
% exponentiation of K0^2=3I then gives its exact zero-pair powers D^beta,
% beta=18/d. Changing variable to D in the affine force integral yields
% exactly Psi_nu and its stated logarithmic limits. With
% D_eta=exp[-i eta^T J(Y,P)], direct differentiation gives the residual
% scalar phi_force'-(eta.l_force)/2; hence the saved phase derivative
% phi_force'=(eta.l_force)/2 is necessary and has the correct sign.
% The pure cubic phase follows by integrating
% (D-h)^2/D^3=1/D-2h/D^2+h^2/D^3, confirming every sign in Phi_c.
% The vertical regular-singular equation and endpoint rotations also
% check. These elementary formulas describe a continuous conditional
% unitary; they do not replace q_in by finitely many samples.
%
% The following independently derived bounds evaluate its remaining
% vertical matrix size and first three q_in derivatives without a run.
% Use the canonical oscillator variables Q=3^(1/4)q3,
% P=(p3-sqrt6 q3)/3^(1/4) from the spin block. In these variables the
% homogeneous vertical oscillator matrix is the antisymmetric J. The
% sign-group quadratic Hessian becomes
% H_sigma^osc=[[-5A_sigma/6,B_sigma],[B_sigma,A_sigma/3]],
% A_sigma=2sqrt3-sigma sqrt6>0,
% B_sigma=(7sqrt6-4sigma sqrt3)/6.
% This follows by substituting p3=3^(1/4)(P+sqrt2 Q) in the complete
% V_sigma, including its symmetrized ZC term. Put
% c_sigma^v=sqrt[B_sigma^2+(7A_sigma/12)^2],
% h_sigma^v=c_sigma^v+A_sigma/4=||H_sigma^osc||.
% The symmetric part of J H_sigma^osc has norm c_sigma^v, giving
% c_+^v<1.81,h_+^v<2.07; c_-^v<5.30,h_-^v<6.79.
% These are distinct from the cubic coefficients c_sigma.
%
% For S'= [J+f(u)J H_sigma^osc]S, define
% I_sigma=int_-pi^pi |f(u)|du,
% I_(sigma,j)=int_-pi^pi |partial_(q_in)^j f(u)|du.
% At fixed q_in with |q_in|<=16, f has the sign of q_in, so the exact
% characteristic logarithm gives
% I_sigma=|log D(pi)|/(sigma gamma d_sigma),
% I_(sigma,j)=(j-1)! a_sigma(pi)^j/
%                 [(sigma gamma d_sigma)D(pi)^j], j>=1.
% The denominator sigma gamma d_sigma is positive for both signs.
% Its safe lower bounds are1.66 and1.08; a_+(pi)<0.017,
% a_-(pi)<0.011, D_+>=0.728 and D_->=0.824. Elementary substitution
% gives the simultaneous bounds
% (I,I1,I2,I3)_+ <(0.192,0.0141,0.000330,0.0000154),
% (I,I1,I2,I3)_- <(0.180,0.0124,0.000166,0.0000045).
% The same bounds hold with any earlier upper time, since the defining
% integrands are nonnegative. The q_in=0 formulas use their continuous
% limits; no division by a vanishing q_in is needed.
%
% Euclidean logarithmic norm, with J's contribution exactly zero, yields
% ||S(u,s)||<=exp[c_sigma^v int_s^u|f|]. Thus for the full factor
% ||S_+||<1.42 and ||S_-||<2.60 on this domain. Differentiated Duhamel
% has products of interval propagators whose exponents ADD to this same
% total bound, not its square at every insertion. For j=1,2,3 it gives
% ||partial_q S||<=exp(c I) h I1,
% ||partial_q^2 S||<=exp(c I)[h I2+h^2 I1^2],
% ||partial_q^3 S||<=exp(c I)[h I3+3h^2 I1 I2+h^3 I1^3].
% Consequently safe rounded bounds are
% (||S||,||S_q||,||S_qq||,||S_qqq||)_+ <(1.42,0.042,0.0023,0.00018),
% (||S||,||S_q||,||S_qq||,||S_qqq||)_- <(2.60,0.220,0.022,0.0025).
% The known endpoint oscillator rotations have norm1 in these variables
% and are q_in independent, so these bounds also hold in the conditional
% interaction frame. Conversion back to(q3,p3) uses the written fixed
% canonical matrix; it is not an orthogonal change and its norms must
% be retained in any physical weighted estimate.
%
% These values control the exact conditional classical matrices and their
% parameter derivatives used to construct a metaplectic factor. They do
% not by themselves bound the unbounded metaplectic-generator derivatives
% on an arbitrary state; those require its oscillator moments. They also
% do not diagonalize the joint oscillator-spin middle operator, establish
% the incoming q_in domain after another factor, or bound the full cross-
% group residual. The physical residual still keeps V4/higher symbols,
% both bands, actual cut/centering/native response and launch/root errors.
% No finite dimension or execution follows solely from these constants.
% SOURCES: actual RED MP51--53/SOURCES plus independent oscillator metric,
% exact characteristic integrals and differentiated matrix Duhamel bounds.
% No new RED bound on these matrix norms was read before saving this block.
%
% ACTUAL MP54--55 AND INITIAL VERTICAL GRAPH WEIGHTS [R, independent WIP].
% Read the full actual RED MP54--55/SOURCES after the preceding independent
% matrix bounds were saved. The logarithmic derivative has the correct
% sign: iU^dagger U_q=Q_A, A=-S^T J S_q, so U_q=-iU Q_A. Its third
% derivative contains iQ_A^3-2Q_A Q_A'-Q_A'Q_A-iQ_A'', giving precisely
% the stated8/12/2 norm coefficients. The characteristic half-density
% pulls p_final back to{D^2,p_in}/2; the conditional metaplectic factor
% adds -D^2 Q_A. This confirms MP55's0.20/1.60 momentum-connection
% constants, with its actual oscillator weights retained. MP54 likewise
% checks the exact spin factors and the two saved launch precisions.
%
% Here is a quantitative graph bound for the ACTUAL baseline input of
% that representation, not a new state or numerical component. Let
% xi_v=(Q,P) be the oscillator pair in the preceding block and define
% L_s(psi)^2=sum_(i1,...,is in{1,2})||xi_i1 ... xi_is psi||^2, L_0=||psi||.
% On the common Schwartz core, followed by graph closure,
% L_s(psi)^2=<psi,F_s(N)psi>,
% F_0=1, F_(s+1)(n)=(n+1)F_s(n+1)+n F_s(n-1).
% This follows by writing Q,P as creation/annihilation operators in
% sum_i xi_i F_s(N)xi_i. In particular
% F_1=2N+1, F_2=4N^2+4N+3,
% F_3=8N^3+12N^2+26N+11.
% All these graph norms are nondecreasing with s. For a conditional
% metaplectic U(q), ordered words transform with S(q)^{tensor s}, hence
% L_s(U(q)psi)<=||S(q)||^s L_s(psi) EXACTLY, without commuting the
% individual Q,P factors. The same inequality integrates over q and
% all spectator variables. It requires the actual |q_in|<=16 domain;
% it is not an estimate for an unrestricted preceding-factor output.
%
% At the prescribed initial time the vertical block S(-pi)=-I. Thus
% k3=-sqrt(epsilon)p3 and commutation of any xi_v component with a
% momentum multiplier is a derivative in k3 with vector coefficient
% of norm c_v sqrt(epsilon), c_v=sqrt(3sqrt3)<2.28.
% In the upper free frame the exact multiplier before the scalar cutoff
% is F(k)exp(i mu.k/epsilon),
% F(k)=[I+(sigma.n(k))(sigma.n_c)]/2 acting on the fixed carrier spinor.
% Its real-momentum norm is<=1. The selected launch box has |k_c|>=2.99
% and the actual cutoff has |k-k_c|<=1, so |k|>=1.99 throughout its
% support. On the complex k3 disc of radius|k|/4,
% |(k+z e3).(k+z e3)|>=(7/16)|k|^2 and |k+z e3|<=5|k|/4.
% Consequently ||F||<2 there, and Cauchy's formula gives
% ||partial_(k3)^j F||<=2 j!(4/|k|)^j<2 j! (2.02)^j.
% No global massless gap or analyticity across k=0 is asserted.
%
% Retain the actual centering phase. The previously checked uniform
% |mu|<9e-7 gives c_v|mu3|/sqrt(epsilon)<0.00841, while
% c_v sqrt(epsilon)*2.02<0.001125. Their sum is<0.01.
% Therefore the vector norm of each j-fold ordered commutator with
% F exp(i mu.k/epsilon) is<=2 j! (0.01)^j, j>=1.
% The j=0 bound remains1. This statement keeps the full projected
% spin multiplier and its phase; the two reflected carrier components
% satisfy the same bound. The lower free coefficient is initially zero.
%
% Moving the multiplier through s ordered oscillator coordinates yields
% L_s(a_initial)<=c_cut Z^(-1/2)[L_s(g_0)
%   +2 sum_(j=1)^s binom(s,j)j!(0.01)^j L_(s-j)(g_0)]
% plus the explicit cutoff-boundary contribution. Here g_0 is the SAME
% uncut Gaussian in the initial quadratic frame. The scalar cutoff
% normalizer c_cut=(E_(g_0)theta^2)^(-1/2) obeys c_cut-1<1e-108 by
% the already checked scalar tail. Z is the ACTUAL projected norm of
% this normalized cut scalar input, not reset to1. Since Z>0.9999 and the
% L_s(g_0) increase, for s<=6 the bracket ratio is below
% 1+2[0.06+0.003+0.00012+0.0000036+0.000000072+0.00000000072]
% <1.12625. This is an upper estimate for the actual cut input; no
% approximation to the prescribed state has been made in the definition.
%
% The boundary term is negligible with a direct finite-derivative bound.
% For b(t)=exp(-1/t),0<t<=1, its explicit derivative polynomial gives
% |b^(j)|<=j!100^j for j<=6. The switch denominator is>=exp(-2),
% so the reciprocal derivative recursion gives
% |partial_t^j[1/(b(t)+b(1-t))]|<=8 j!1700^j.
% The product rule then bounds the unit switch derivative by9 j!1700^j.
% Radial differentiation on |u|>=2048 gives the sufficient bound
% |partial_(u3)^j theta(|u|/2048)|
% <=9(j!)^2(6800/2048)^j<1e10, 1<=j<=6.
% Indeed the radial derivative polynomials r^(l-j)P_(j,l)(u3/r)
% have combined coefficient norm at most4^j j!: differentiation adds
% at most3j+1 times that norm, which is below4(j+1).
% Its support has whitened |y|>16. The finite product rule, oscillator
% matrix norm<2.37, the above multiplier bounds and Gaussian radial
% moments through degree12 bound the total six-word boundary contribution
% by1e28 exp(-128)<3e-28. Lower word orders obey the same loose margin.
% Thus, with a common convenient roundoff-free margin,
% L_s(a_initial)<1.13 L_s(g_0), 0<=s<=6.
% The inequality includes all exact normalization/cutoff/spin/centering
% factors. It concerns this initial vertical graph, not propagated input
% into the next noncommuting sign factor.
%
% A finite numerical bound follows by elementary Gaussian moments alone.
% For the uncut q3,p3 vacuum, the oscillator mean occupation is
% n_bar=5/(2sqrt3)-1/2<1, and the anomalous moment squared is
% n_bar(n_bar+1)<2. Wick contraction gives the falling moments
% <N^(underline j)>
% =sum_(l=0)^floor(j/2) (j!)^2 n_bar^(j-2l)
%      [n_bar(n_bar+1)]^l/[(j-2l)!(l!)^2 4^l].
% For j=0,...,6 this is bounded by(1,1,4,24,204,2220,29520).
% The ordinary moments of orders1,...,6 are consequently below
% (1,5,37,377,4921,78365). Expanding (N+6)^6 gives a bound809453<900^2.
% The recurrence above gives F_6(N)<=[2N+11]^6, hence
% L_6(g_0)<7200 and L_6(a_initial)<8200.
% A single conditional vertical factor therefore has L_6 below
% 1.42^6*8200 for plus or2.60^6*8200 for minus, ON its written incoming
% q-domain. These deliberately loose but actual-input graph estimates
% can be used in the connection/cross-residual calculation. They neither
% establish that domain after another factor nor justify transferring
% Gaussian tails to its output. The complete scalar/zero-pair/spin and
% opposite-band evolution, response, return root and coefficient remain.
% SOURCES: actual complete RED MP54--55/SOURCES; independent ordered-word
% oscillator recursion, complex-disc derivative bounds, cutoff reciprocal
% recurrence and finite Wick moments. No new RED actual-input graph bound
% was read before saving this extension, and no scientific code was run.
%
% ACTUAL MP56 RESPONSE GAUGE AND EXPLICIT DISPLACEMENT [R, WIP].
% Read the entire actual RED MP56/SOURCES after saving the preceding
% actual-input graph extension. Its exact full-four-spinor identities
% check directly: T_g=-i Op_W(g)/epsilon gives
% B_g=B_f+[H_f,T_n+T_g]-i epsilon partial_t(T_n+T_g),
% while the physical frequency vertex omits the last time derivative.
% The detector adds[Q_f,T_n+T_g], with
% [Q_f,T_g]=Op_W(partial_k g)+epsilon[A,T_g]. The latter term is not
% removed by calling g scalar. The transformed residual reconstructs as
% s_physical=s_g+(T_n+T_g)r, retaining weighted base-residual costs.
% The same scalar g acts in both bands; no relative phase or reference
% state has been changed. The matched launch direction is merely added
% and subtracted in the original implicit-function response formula.
%
% Expanding partial_t g+{g,h_+}, with g=c+l.xi+xi^T G xi/2, gives
% c'=b_*, l'+K^T l=grad b_*,
% G'+G K+K^T G=b_*''+D^3h_+[J l,.,.].
% In particular the third-derivative sign follows from
% l^T J D^3h=-D^3h[Jl]. Then d=Jl satisfies the forced Hamiltonian
% variation d'=K d+J grad b_*. This independently verifies MP56's
% sign and its full-matrix versus principal-jet distinction.
%
% The constant and linear coefficients need no new integration run.
% Put theta=tau/sqrt3,0<=theta<=2pi, and use the same principal
% fixed-frequency orbit already independently checked in both sources.
% Define the following explicitly evaluated functions:
% X=-57/4+(77/68)cos4theta+(11sqrt2/17)sin4theta,
% Y=69/8-(9/136)cos4theta-(45sqrt2/68)sin4theta,
% L1=-45/8+(1/8)cos4theta+2sqrt2 sin4theta,
% Atheta=135/8-(755/136)cos4theta-(25sqrt2/17)sin4theta,
% Wtheta=135theta/8-(755/544)sin4theta
%                         -(25sqrt2/68)(1-cos4theta).
% Here Wtheta is the FIXED-TIME angular displacement, not an extra
% preparation parameter. The radial Hamilton equation gives
% Y=(3sqrt2/2)X'-(5/2)X-F1-(3/2)F1_y.
% Homogeneity gives F1_y+F1_L=3F1, hence the angular-speed change
% Atheta=3F1-sqrt2 X' and Wtheta'=Atheta, Wtheta(0)=0.
% Thus the complete rotating-coordinate displacement is
% d_q=(X,Wtheta,0),
% d_k=(sqrt6 Y-sqrt3 Wtheta,
%                  sqrt3(L1-X)+sqrt6 Wtheta,0), l=-Jd.
% This is a full canonical displacement in the same physical Cartesian
% rotating frame; retaining only X would miss the secular angular terms.
% Its initial value agrees exactly with MP56's
% d(0)=(-223/17,0,0,291sqrt6/34,259sqrt3/34,0).
%
% The constant phase, with c(0)=0, is
% c(tau)=sqrt3[45theta/8-sin4theta/32
%                              +(sqrt2/2)(cos4theta-1)].
% It satisfies dc/dtau=F1=45/8-cos4theta/8-2sqrt2 sin4theta.
% The full-period launch/time direction therefore is
% (-223/17,291sqrt6/34,-11sqrt3/2,-135pi sqrt3/4)
% in the ORIGINAL (R,p_r/omega,L/(omega r_*),tau_T) labels. The time
% component cancels Wtheta(2pi) in the fixed-frame return section;
% it acts through the original endpoint/half-time/mean-time rows, not
% through a changed fixed-time source or changed averaging definition.
% Elementary suprema give |X|<16.3, |Y|<9.63, |L1-X|<11.55,
% |Wtheta|<108.5, |d_k1|<212, |d_k2|<286, ||l||<375 and |c|<64.
% These are bounds for an EXACTLY specified response gauge coefficient,
% not for the actual packet response or return-root displacement.
%
% G(0)=0 keeps the gauge's initial quadratic part zero. Its forcing is
% completely fixed by the native principal symbol in the rotating axes:
% b(tau,q,k)=-sum_i[(R_theta k)_i]^4/(6|k|)
%    +(2sqrt(2/3)/(3|q|^(3/2)))
%                  sum_i(R_theta q)_i[(R_theta k)_i]^3.
% Take its Hessian at(q,k)=(e1,(sqrt6,sqrt3,0)), and add the displayed
% D^3h_+[d] term, with h_+'s actual cubic Taylor polynomial already
% saved above. The unique matrix solution is
% G(t)=exp(-K^T t)[int_0^t exp(K^T s)
%       (b_*''(s)+D^3h_+[d(s)])exp(Ks)ds]exp(-Kt).
% It is finite-dimensional and has only the known exponential,
% trigonometric and polynomial frequencies. No bound on its actual
% entries or on the resulting cubic symbol/graph residual is inferred
% from this representation. That quantitative step, the retained full
% lower/spin/endpoint source, actual weighted propagation and the actual
% four-return Jacobian remain necessary. No new run is specified here.
% SOURCES: complete actual original RED MP56/SOURCES, independently
% checked by exact CAR-frame gauge algebra and Hamiltonian ray equations;
% direct trigonometric reduction of the previously checked X,F1,F1_y.
% All original finite-frequency/native-input/root limits remain explicit.
%
% BOUNDED PARTICULAR QUADRATIC RESPONSE GAUGE [R, independent continuation].
% The G(0)=0 choice above is convenient but is not a physical preparation
% requirement: g is added and subtracted from the EXACT response. A
% nonzero quadratic G(0) is allowed only with its full initial term
% -T_g(0)u(0) retained. Here is an explicit alternative which avoids the
% exponentially growing homogeneous part of G without changing the state.
% No reduction of the TOTAL response error is inferred before that initial
% compensation and the retained source have both been propagated.
%
% Let lambda=1/sqrt3, M=diag(J_z,J_z), J_z q=e3 cross q, and let
% W=-J M be the symmetric Hessian of L_z=q1 k2-q2 k1, ||W||=1.
% The secular part of d above is a t M z_*, a=135/(8sqrt3).
% Rotational invariance of the actual continuum rotating h_+ gives
% D^3h_+[M z_*]=-M^T h_+''-h_+'' M.
% For the Sylvester operator L(G)=G K+K^T G this is precisely L(W):
% M commutes with the canonical J and K=J h_+'', so
% L(W)=M h_+''-h_+'' M.
% Thus the quadratic forcing has the exact finite form
% C(t)=C0+a t L(W)+Cplus exp(4i lambda t)+Cminus exp(-4i lambda t),
% with Cminus the complex conjugate of Cplus and all matrices symmetric.
% C0,Cplus,Cminus are fixed by the literal b_*''+D^3h_+[d] above.
% In particular its secular coefficient has ZERO projection onto the
% zero eigenspace of L; an artificial t^2 resonant term is unnecessary.
%
% The required finite spectral maps are explicit. The horizontal K_h
% obeys K_h^3=K_h/3, with projectors
% P_h0=I_h-3K_h^2,
% P_hplus=(3K_h^2+sqrt3 K_h)/2,
% P_hminus=(3K_h^2-sqrt3 K_h)/2.
% The vertical block has projectors
% P_vplus=(I_v-i sqrt3 K_v)/2, P_vminus=conj(P_vplus).
% Their eigenvalues are0,+lambda,-lambda,+i lambda,-i lambda.
% Extend each by zero to the complementary coordinates and denote them
% by P_r. For any symmetric G define Q_(r,s)(G)=P_r^T G P_s.
% Summing all ordered pairs gives the identity and diagonalizes L with
% eigenvalues r+s. The symmetric subspace is retained by each sum below:
% Z0=sum_(r+s=0) Q_(r,s),
% D1=sum_(r+s!=0) Q_(r,s)/(r+s),
% Rplus=sum Q_(r,s)/(r+s+4i lambda), Rminus=conj(Rplus).
% There is no harmonic resonance: |r+s+/-4i lambda|>=2lambda.
%
% Put Cbar=C0-a W. An explicitly bounded particular solution is
% G_b(t)=a t W+D1 Cbar+t Z0 Cbar
%       +exp(4i lambda t)Rplus Cplus
%       +exp(-4i lambda t)Rminus Cminus.
% Direct differentiation gives G_b'+L(G_b)=C(t), with no asymptotic
% approximation. Its generally nonzero initial value is written by the
% same formula at t=0. The original zero-initial solution is exactly
% G_zero(t)=G_b(t)-exp(-K^T t)G_b(0)exp(-K t).
% This identity identifies, rather than silently discards, the growing
% homogeneous term. Using G_b transfers that term into the exact initial
% response and its subsequent full unitary propagation.
%
% The projector bounds can also be checked by finite hand algebra. The
% rank-one horizontal projectors have squared norm1561/48. For P_h0,
% direct substitution in the coordinates(q1,q2,k1,k2) gives
% (0,sqrt2 q1+q2,-3sqrt3 q2+sqrt2 k2,3sqrt3 q1+k2),
% whose squared Frobenius norm is60. Each vertical rank-one projector
% has squared norm25/3. Therefore sum_r||P_r||<25 and
% ||Z0||<60+2(1561/48)+2(25/3)<143,
% ||D1||<625sqrt3, ||Rplus||=||Rminus||<625sqrt3/2.
% These are real/complex matrix norms induced by the same Euclidean
% phase-space norm; no orthogonality of these nonnormal projectors is
% assumed. Consequently on0<=t<=12,
% ||G_b(t)||<117+2800||Cbar||
%                  +542(||Cplus||+||Cminus||).
% This is a finite explicit bound once the already specified three
% Hessian coefficient matrices are evaluated; it has no exp(4pi)
% factor. The evaluation and the resulting weighted operator/source
% error still have to be carried out on the actual retained state.
%
% This alternative gauge does not alter the reflected mixture, exact
% cut/projector/centering, frequency row or four-return conditions.
% It retains the same scalar g in both bands, the full physical detector
% connection and the T_g r residual term. In particular the new initial
% contribution is i Op_W(xi^T G_b(0)xi/2)u(0)/epsilon; dropping it would
% change the preparation's spacing derivative. This remains an exact
% representation and quantitative matrix control, not a propagated error
% or a numerical coefficient. No scientific execution is specified.
% SOURCES: independent rotational differentiation of the actual h_+,
% explicit K_h/K_v projectors and direct Sylvester substitution, after
% the actual MP56 comparison. No new RED bounded-particular calculation
% was read before this save; all lower-band/native/root limits persist.
%
% ACTUAL MP57--58 AND COMPLETE PRINCIPAL QUARTIC RESONANCE [R, WIP].
% Read the entire actual original RED MP57--58/SOURCES. Direct canonical
% calculation confirms its cubic homological identity and conjugation
% sign: U=exp(-i zeta G3) gives U^dagger H U=H2+
% zeta^2[W4+(i/2)[G3,V3]]+..., since i[H2,G3]=V3. The bounded-cutoff
% control and exact conjugation remainder are still needed for a global
% unitary; no essential self-adjointness of an unrestricted cubic is used.
% Its fixed-L radial calculation also checks independently. In tortoise
% coordinate x, U=f/r^2=1/3-x^2/9+2x^3/81+2x^4/81+... . The cubic
% Poisson bracket has resonant coefficient-30 alpha^2/w, giving the
% squared-energy coefficient23/54 and then5sqrt3/(108L) for E's I^2.
% The dilation characteristic function and its fourth derivative agree
% with the displayed covariance det15/2 and c^2=121/192. No sign or
% normalization repair was found in those radial principal coefficients.
%
% MP58's actual-input extension also checks: the three-mode recurrence
% gives L_(0,...,4)=(1,sqrt3,sqrt15,sqrt93,sqrt705), and its0.67
% tensor commutator constant includes both the projected spin multiplier
% and actual centering phase. The normalized/cut L4<160 and
% ||I_h^2 a_initial||<2804 follow without replacing the input by the
% Gaussian. The resulting<9e-5 unitary-factor defect belongs ONLY to
% that radial resonant scalar factor on the actual REFERENCE launch.
% It is not propagated through another factor or assigned to the unknown
% matched launch. Its full source/readout/domain qualifications remain.
%
% A further classical structure identifies which other PRINCIPAL
% quartic resonances can occur. It does not supply the matrix/ordering
% terms of that order. Work near the same equatorial circular orbit,
% with L=|q cross k|>0, L_z>0. The exact central Hamiltonian depends on
% radial(r,p_r) and L only. Its angular action can be chosen as
% J_v=L-L_z, with a smooth local Cartesian oscillator chart at J_v=0.
% One explicit angular-momentum chart is
% J_z=L-J_v, J_x+iJ_y=sqrt(2L-J_v)(Q_v+iP_v)/sqrt2,
% J_v=(Q_v^2+P_v^2)/2,
% with {Q_v,P_v}=1, so {J_x,J_y}=J_z as required.
% A conjugate orbital phase completes the angular Darboux chart. The
% usual action-angle singularity at zero amplitude is absent in Q_v,P_v.
% This is a local canonical-coordinate statement; a single angular
% phase chart over every winding is not assumed.
%
% The radial normal-form transformation can preserve L exactly; its
% L-dependent generating function then changes the conjugate orbital
% phase as required by canonicity. Let I be the resulting radial
% hyperbolic action, with quadratic part I_h. The checked energy form is
% E(L,I)=L/sqrt3+I/sqrt3+5sqrt3 I^2/(108L)+O(I^3).
% Subtracting the ORIGINAL rotating generator L_z/sqrt3 gives exactly
% H_rot=(I+J_v)/sqrt3+5sqrt3 I^2/(108L)+O(I^3).
% Count I and J_v as degree2, and L-L_* as degree1 with L_*=sqrt3.
% Therefore its ENTIRE principal quartic normal form is
% (5/108) I_h^2: no J_v^2, I_h J_v or neutral-coordinate quartic term
% is present in this action chart. The dependence of1/L first enters
% degree5. This eliminates an artificial exponential estimate on every
% principal quartic component; it does not make the required canonical
% transforms or the resulting remainder disappear.
%
% This conclusion is about the principal classical symbol only. Weyl
% ordering already distinguishes Op_W(I_h^2)=Ihat_h^2-1/4. The cubic
% commutator's third-derivative contraction adds its own scalar term;
% spin-linear commutators, the full second band symbol and its matched
% coordinate, subprincipal transport, and the continuum epsilon^2
% matrix terms are still retained. They cannot be inferred from the
% scalar radial action or deleted because the principal normal form is
% simple. Nor are the prescribed packet, native spacing vertex or four
% return conditions transformed without their exact conjugate terms.
% The result organizes the necessary retained propagation and its error
% calculation; it is not a propagated quantum normal-form theorem or an
% evaluated finite-frequency coefficient. No new run is specified.
% SOURCES: full actual RED MP57--58/SOURCES, independently checked by
% finite Poisson/Weyl algebra, radial expansion and oscillator moments;
% independent local angular-action and L-preserving radial canonical
% argument for the complete PRINCIPAL quartic scope. No new RED full
% angular normal-form derivation was read before saving this extension.
%
% EXACT RADIAL MATRIX RESOLVENT AND QUANTUM BARRIER NORMAL FORM [R, WIP].
% This is a full two-component radial reorganization of the fixed continuum
% propagator used in MP8, not a new probe or a change of its angular weights.
% Work in the same scaled variables, epsilon=2^-24, and retain the actual
% smooth subcritical core. In a native continuum angular block,
% H_kappa=sigma_y p_r+L sigma_x/r+{v,p_r}/2,
% p_r=-i epsilon partial_r, L=epsilon kappa. Changing the sign of kappa
% interchanges the partner channels below; write L>0 for their formulas.
% The regular full-Dirac center domain is retained. Let
% f=1-v^2>0, dx/dr=1/f, g'(r)=v/f,
% M(r)=f^(-1/4)exp[-atanh(v)sigma_y/2],
% W(x)=L sqrt(f(r(x)))/r(x), K_L=sigma_y p_x+W sigma_x.
% Direct substitution gives the EXACT stationary intertwining identity
% (H_kappa-E)M exp[-i E g/epsilon]
%   =exp[-i E g/epsilon]f^(-1)M^(-1)(K_L-E).
% It includes the endpoint half-density term -i epsilon v'/2. Indeed
% M^2=(I-v sigma_y)/f and M'/M=vv'/(2f)-v'sigma_y/(2f);
% these cancel precisely the connection in the original first-order ODE.
% Dropping that connection would not give the stated partner barriers.
%
% Consequently the radial Green kernel, with dr on the native side and
% dx on the static side, is exactly
% R_PG(E;r,r')=M(r)exp[-i E(g(r)-g(r'))/epsilon]
%              R_static(E;x(r),x(r'))M(r').
% The source factor f M cancels dx'=dr'/f, so no missing Jacobian is
% hidden in this formula. It holds with matched regular-center and outgoing
% resolvent boundary conditions; boundary values can then enter Stone's
% spectral formula. For nonreal E the phase and resolvent must be treated
% together, rather than separately bounded on the whole half-line.
% The fixed core gives ||M||<1/sqrt(1-sqrt(5/6)) and
% ||M^-1||<sqrt(1+sqrt(5/6)). This energy-dependent stationary map is
% NOT asserted to be a pointwise unitary map of equal-time initial data.
%
% Squaring the static matrix is exact:
% K_L^2=p_x^2+W^2-epsilon sigma_z W',
% R_static(E)=(K_L+E)diag[(A_+-E^2)^(-1),(A_--E^2)^(-1)],
% A_s=p_x^2+W^2-s epsilon W', s=+1,-1.
% The (K_L+E) numerator and both signs of E retain the full matrix band
% content. No lower-band preparation or spacing response is discarded.
% In the original MP8 calculation these kernels must still act on the
% EXACT projected/cut/centered initial angular coefficients, all retained
% angular coherences, the physical coordinate blocks and the native rank
% zero/four spacing insertion. The displayed identity does not change
% any of those inputs or supply their evaluated propagation by itself.
%
% An independent Weyl calculation now gives the first QUANTUM radial
% normal form of this exact scalar resolvent. Center x=0 at r=1 in the
% Schwarzschild exterior. The previously checked principal expansion is
% U=f/r^2=1/3-x^2/9+2x^3/81+2x^4/81+O(x^5),
% W=L/sqrt3[1-x^2/6+x^3/27+5x^4/216+O(x^5)].
% Thus, with w=L/3, c3=c4=2L^2/81,
% b1=L/(3sqrt3), b2=L/(9sqrt3),
% A_s-L^2/3=p_x^2-w^2 x^2+c3 x^3+c4 x^4
%                   +s epsilon b1 x-s epsilon b2 x^2+O_weight(5).
% Here x,p have weight1 and epsilon weight2. This is a local formal jet
% with an explicit analytic exterior coefficient; a controlled cutoff/
% conjugation remainder on the actual propagated state remains required.
%
% Rescale x,p by sqrt(epsilon) for the finite Weyl algebra. With
% a=(p+w x)/sqrt(2w), b=(p-w x)/sqrt(2w), [a,b]=i,
% I=(ab+ba)/2, the quadratic generator is2w I. Put
% alpha=c3/(2w)^(3/2), beta=b1/sqrt(2w),
% V3=alpha(a-b)^3+s beta(a-b) in Weyl order.
% The solution of i[2w I,G]=V3 is
% G3=alpha/(6w)(a^3+b^3)-3alpha/(2w)(a^2 b+a b^2),
% G1=s beta(a+b)/(2w), G=G3+G1,
% with every polynomial still Weyl ordered. Conjugation by
% exp[-i sqrt(epsilon)G] gives at the next order
% V4+(i/2)[G,V3]. Its principal resonant I^2 coefficient is23/54.
% The mixed linear/cubic resonant I coefficient cancels EXACTLY:
% s[b2/w-6alpha beta/w]=0. The linear/linear constant is
% beta^2/(2w)=1/12. The cubic Moyal contraction is also indispensable:
% G3 Lambda^3 V3=96alpha^2/w, so (i/2)[G3,V3] contributes the additional
% scalar2alpha^2/w=1/81. Hence the resonant Weyl symbol is
% (23/54)(ab)^2+31/324.
% Since Op_W((ab)^2)=I^2-1/4, its OPERATOR form is
% (23/54)I^2-7/648. Both partner signs give exactly the same constant.
%
% Restoring epsilon and the unscaled dilation Ihat, this gives the formal
% local operator normal form
% A_s=L^2/3+(2L/3)Ihat+(23/54)Ihat^2-7epsilon^2/648+O_weight(5).
% The positive square-root spectral branch therefore has
% E_+=L/sqrt3+Ihat/sqrt3+5sqrt3 Ihat^2/(108L)
%                         -7sqrt3 epsilon^2/(1296L)+O_weight(5).
% This computes a quantum constant that the principal radial action alone
% cannot determine. It is not obtained by inserting the earlier local
% free-frame I/12 symbol as though it were this transformed operator.
% The angular rotation, exact spectral measure, energy-dependent boost/
% phase, initial state and readout must all be retained when using it for
% MP8. In particular no absolute return coefficient or error is inferred
% from this formal normal form, and no additional finite-order run is
% proposed. Its purpose is an exact full-matrix propagation route with
% a computable barrier remainder, rather than another unnormalized
% state truncation known to fail for the selected member.
% SOURCES: independent direct stationary first-order Dirac substitution,
% Green-kernel source/Jacobian calculation and finite Weyl commutator
% algebra, after the actual MP57 radial principal comparison. No new RED
% exact radial resolvent or quantum-constant derivation was read before
% this save; the original sine/endpoint spacing insertion is unchanged.
%
% ACTUAL MP59--60 AND COMPLETE RADIAL PROPAGATOR ASSEMBLY [R, WIP].
% Read both complete actual original RED blocks and all SOURCES after
% saving the radial barrier calculation. The Cartesian triad identity
% checks directly: for n=k/|k|, L=q cross k and b=3c p_r/(4r^(5/2)p^2),
% dot n=(2bL-Omega e_z)cross n and dot Lhat=-Omega e_z cross Lhat.
% Thus its local spin lift solves precisely the stated subprincipal
% transport; it cannot be omitted from transformed input or detectors.
% The five cubic Moyal sectors sum to-sqrt3/3888, and the spin scalar
% is-11gamma/24. With Q2=I/12 and the Weyl1/4 correction they give
% -7/1296 in tau units, independently matching the partner-barrier
% calculation above. These are named same-pair methods, still[R].
%
% The exact Cauchy-graph completion in MP60 also checks without a gap or
% band projection. Set x(r_*)=g(r_*)=0 and b0=sup|v|<1. Static solutions
% chi(t,x)=exp(-it K_L/epsilon)phi(x) have conserved density rho=chi^dag chi
% and current j=chi^dag sigma_y chi. The graph t=g(x), g_x=v, has flux
% (rho-vj)dx, which equals |M chi|^2 dr because M^2=(I-v sigma_y)/f.
% The same first-order regular-center domain gives zero endpoint flux.
% For compact data the graph intersects its causal cone in a compact set,
% since |g(x)|<=b0|x|; no flux at infinity is discarded. The forward and
% reverse spacelike Cauchy maps are isometries and inverse by uniqueness.
% Consequently the ENERGY-INDEPENDENT nonlocal map
% (W_L phi)(r)=M(r)[exp(-i g(r)K_L/epsilon)phi](x(r))
% is unitary and H_kappa W_L=W_L K_L. The variable-time evaluation is
% performed after forming the solution; it is not a pointwise exponential
% of a position-dependent scalar operator on initial data.
%
% Finite propagation gives its kernel support |x-y|<=b0|x|. Unitarity
% then bounds its outer projection by the input outside(1-b0)R. Integrating
% these tail inequalities proves
% |||x(r)|^m W_L phi||<=(1-b0)^(-m)|||x|^m phi||,
% |||x|^m W_L^dag psi||<=(1+b0)^m|||x(r)|^m psi||.
% Since |r-r_*|<=|x(r)|<=|r-r_*|/(1-b0^2), either direction between
% the native centered radial weight and static |x| has sufficient factor
% (1-b0)^(-m). At this core the first-weight factor is<11.478. Exact
% intertwining gives the same bounds after each corresponding power of H
% or K. This verifies the graph-domain claim, not a narrow-packet premise.
%
% Here is the COMPLETE fixed-member assembly in this representation.
% Let W be the direct sum over the original kappa,j,m blocks, retaining
% their coherent amplitudes and both signs. It depends only on the fixed
% continuum background. For each reflected input member and launch label
% y=(R,p_r,L,T), form EXACTLY
% phi(y)=W^dag u_in(y), phi_delta(y)=W^dag w_in(y),
% phi_j(y)=W^dag partial_(y_j)u_in(y), j=1,2,3,
% Btilde=W^dag B W, Qtilde_i=W^dag q_i W.
% Here u_in is the original normalized cut/P0/continuum-centered spinor;
% w_in includes P2 and the original native-centering correction, not a
% static spectral-band preparation. The factor W is independent of delta
% and the launch labels, so these derivatives have no discarded W_delta
% or W_j terms. Btilde is the full sine/endpoint derivative; it is not the
% rotational average of B or the scalar barrier's parameter derivative.
%
% Put V(t)=exp(-it K/epsilon). The base and response are
% u_t=V(t)phi,
% w_t=V(t)phi_delta-(i/epsilon)int_0^t V(t-s)Btilde V(s)phi ds.
% The original Cartesian centroid and its spacing derivative are exactly
% q_i(t)=<u_t,Qtilde_i u_t>,
% q_delta,i(t)=2Re<u_t,Qtilde_i w_t>.
% Average these expectations over the original reflection pair. Launch
% columns use V(t)phi_j; the T column uses the exact time derivative.
% The frequency row is
% E0=<phi,K phi>, E_delta=2Re<phi_delta,K phi>+<phi,Btilde phi>,
% E_j=2Re<phi_j,K phi>.
% In particular the physical frequency still contains the full Btilde;
% a scalar phase removed in a response gauge does not alter this row.
% With n(t)=q(t)/|q(t)|, r_delta=n.q_delta and r_j=n.q_j. Substitution
% in the ORIGINAL four rows F=(r(T/2)-R,r(T)-R,q_y(T),E0-1) and in
% T^-1 int r(t)dt gives the same matched-root derivative already stated,
% Rcal_delta=partial_delta Rcal-grad_y Rcal (D_y F)^(-1)F_delta.
% No detector, launch centroid or energy-matching convention is changed.
%
% This can be evaluated with the exact full matrix radial Green kernels,
% or with a norm-controlled approximation to the same K propagation.
% The partner factorization makes the latter explicit:
% exp(-it K/epsilon)=cos(t sqrt(K^2)/epsilon)
% -i K sin(t sqrt(K^2)/epsilon)/sqrt(K^2),
% with the continuous value at zero and BOTH scalar partners retained.
% It does not license replacing the original input by one partner.
% In a complete spectral representation, the insertion is the exact
% divided difference
% [exp(-itE/epsilon)-exp(-itE'/epsilon)]/(E-E') Btilde_(E,E'),
% with diagonal value-(it/epsilon)exp(-itE/epsilon). The earlier angular
% blocks specify its rank-zero/four links and Qtilde's rank-one links;
% the radial map preserves their tensor ranks and full coherences.
%
% Residual certification is likewise representation invariant. If trial
% tilde u,tilde w satisfy r=(i epsilon partial_t-K)tilde u and
% s=(i epsilon partial_t-K)tilde w-Btilde tilde u, unitarity gives
% ||e(t)||<=||e(0)||+epsilon^-1 int||r||,
% ||f(t)||<=||f(0)||+epsilon^-1 int(||s||+||Btilde e||).
% The saved physical estimate ||B psi||<=1180 sum_(j=0)^3||H^j psi||
% becomes the identical K-graph estimate for Btilde. The exact W weight
% bound above carries physical-coordinate residual estimates without
% substituting a radial mean detector. Apply the already saved dual and
% centroid error formulas to these trial states, including their actual
% initialization, normalization, launch and quadrature defects. A root
% certificate still needs the actual enclosed F, D_yF, mean and response
% residuals, plus the stated positive-radius and winding margins.
%
% This is a complete identification of state/source/readout and errors for
% the reorganized FULL continuum matrix calculation. It does not yet
% select a controlled finite propagated space, certify that root, or give
% its absolute coefficient. The formal normal-form agreement is one
% possible means to construct the needed trial propagation, with all
% conjugations and remainders retained. No bounded execution is requested
% from these identities alone and no accepted input check is rerun.
% SOURCES: complete actual original RED MP59--60/SOURCES, independently
% checked by Pauli/Moyal algebra, spacelike flux and tail integration;
% exact unitary pullback of the previously saved MP8 state, native
% response, Cartesian detector, frequency row and residual identities.
% No other-lane propagated result or coefficient is assumed.
%
% EXACT ONE-PARTNER FOLDING AND ACTION CHART [R, independent WIP].
% This block was derived before opening the RED's new MP61. Its announced
% potential-Morse chart motivated comparison, but no actual MP61 formula or
% coefficient was read. The following action chart has a different purpose:
% it retains the kinetic coefficient instead of treating a quadratic-looking
% potential as a completed parabolic propagator.
%
% For each L=epsilon|kappa|>0, first exchange radial components if necessary
% to use the static operator K=[[0,S^dag],[S,0]],
% S=epsilon partial_x+W(x), W=L sqrt(f)/r. Use the inherited regular full
% Dirac domain at the center, not an independently chosen momentum domain.
% Both ker S and ker S^dag are zero. Indeed their only formal zero solutions
% are exp[-int W/epsilon] and exp[+int W/epsilon]; the former behaves as
% r^(-|kappa|) at the center and fails its regular domain, and the latter
% behaves as r^(|kappa|) at infinity and is not in L2. The same endpoint
% statements follow from W dx=L dr/(r sqrt f). Consequently the polar
% decomposition S=V sqrt(A_+), A_+=S^dag S, has UNITARY V. There is no
% asserted spectral gap at zero and no bounded inverse of sqrt(A_+).
%
% The exact unitary
% F_L=(1/sqrt2)[[I,I],[V,-V]]
% gives F_L^dag K F_L=diag(sqrt(A_+),-sqrt(A_+)). Thus one scalar partner
% plus its polar map suffices for BOTH energy bands; dropping the second
% diagonal block is still forbidden for the actual free-band preparation.
% Compose this map with the saved Cauchy-graph W_L. The original inputs,
% P2/centering response, full B and Cartesian Q are all conjugated by this
% fixed unitary. Their off-diagonal energy-band entries are retained. This
% reduces the spectral representation without changing the MP8 observable,
% but does not yet give a controlled finite representation of V or A_+.
%
% An explicit stationary comparison chart is available on the entire
% native annulus .9<=r<=1.1, outside the selected core. Write
% f=1-2/(3r), U=f/r^2, w=L/3, s=r-1, F(r)=dx/dr=r/(r-2/3),
% b(r)=sqrt[3(r+2)/r^3],
% H(s)=2 int_0^1 t b(1+st)F(1+st)dt,
% y=s sqrt(H(s)), a(y)=dy/dx=b(r)/sqrt(H(s)).
% Here H is a positive ordinary function, not the matter Hamiltonian.
% Since U(1)-U(r)=(r-1)^2(r+2)/(3r^3), this map obeys EXACTLY
% L^2[U(1)-U(r)]=w^2 y^2 a(y)^2,
% y^2=(2/w) int_0^x L sign(t) sqrt[U(1)-U(r(t))]dt.
% It is smooth and orientation preserving across x=y=0, with a(0)=1.
% This is the action-Liouville chart, not the algebraic assignment that
% makes U directly quadratic while leaving a different kinetic function.
%
% The unitary coordinate map (T chi)(x)=sqrt(a(y(x)))chi(y(x)) gives
% T^dag p_x T={a,p_y}/2 and
% T^dag(A_s-L^2/3)T
%  =a(p_y^2-w^2 y^2)a-s epsilon W'(x(y))+epsilon^2 a^2 R(y),
% R=a_yy/(2a)-a_y^2/(4a^2), A_s=p_x^2+W^2-s epsilon W'.
% The symbol s in A_s denotes the partner sign +/-1, separately from the
% temporary radial difference r-1 above. The stated kinetic ordering follows
% from ({a,p_y}/2)^2=a p_y^2 a+epsilon^2(aa_yy/2-a_y^2/4).
% For a stationary equation the nonunitary half-density a^(-1/2) instead
% removes its first derivative. With Delta=E^2-L^2/3 one has exactly
% (A_s-E^2)[a^(-1/2)chi]
% =a^(3/2)[p_y^2-w^2y^2-Delta/a^2
%             -s epsilon W'/a^2+epsilon^2 R]chi.
% The factor Delta/a^2 and partner term must remain; neither is included
% in a constant-frequency inverted oscillator merely by this change of
% coordinate. An energy-dependent comparison map or an actual bounded
% retained perturbation is needed to remove them for propagation.
%
% Explicit real-annulus margins need no fitted constants. On .9<=r<=1.1,
% 2.53<F<3.86, |F'|<12.25, |F''|<105,
% 2.64<b<3.46, |(log b)'|<1.5, |(log b)''|<1.86.
% Therefore |b'|<5.19, |b''|<14.23. For c=bF,
% 6.67<c<13.36, |c'|<62.5, |c''|<546.
% The integral definition gives
% 6.67<H<13.36, |H'|<125/3, |H''|<273.
% With l=log a this implies |l_r|<4.63 and |l_rr|<42. In particular
% .72<a<1.34, |a_y|=|l_r/F|<1.84,
% |a_yy|=|(l_rr-l_r F'/F)/(F^2 a)|<14,
% |R|<12, |a^2 R|<22.
% These bound the EXACT coordinate ordering contribution by22epsilon^2
% as a multiplication operator on this annulus. For the selected epsilon
% it is<8e-14. They do not bound the other retained kinetic/partner terms
% or a cutoff commutator, and therefore are not a propagator error bound.
%
% The local expansions also provide a check independent of interval margins:
% U(x)=1/3-x^2/9+2x^3/81+2x^4/81+O(x^5),
% y=x-x^2/27-175x^3/5832+O(x^4),
% a(y)=1-2y/27-541y^2/5832+O(y^3), R(0)=-61/648.
% This coordinate-ordering constant is not the previously computed
% full barrier normal-form constant -7/1296. The latter also contains the
% partner, canonical-conjugation and square-root terms. The actual native
% input is not replaced by a Gaussian in y, and a later-time annulus or
% energy-window premise has not been inferred from initial localization.
% SOURCES: direct polar decomposition on the inherited regular radial
% Dirac domain; exact exterior factorization of U(1)-U; direct half-density
% differentiation and rational derivative bounds. No new numerical run or
% actual RED MP61 source was used before this independent save.

% ACTUAL MP61 INITIAL-TUBE COMPARISON [R, WIP].
% Read the entire actual MP61/SOURCES after saving the independent folding
% and action chart above. Its weighted polynomial input estimate checks:
% Gamma(l+3/2)/Gamma(3/2)<=[(l+2)/2]^l is the arithmetic-geometric mean
% bound for its l factors; the degree2j+1 and rho=1/4 Hermite weights give
% the stated ratio<.551 and total weighted squared norm A_(1/4)<9.
% The ACTUAL cut/projected normalizers and same centering translation are
% retained. Mehler then bounds the actual input outside .94<=r<=1.06 by
% 10^-24, uniformly on the saved launch box. For precision,13.85 is the
% first threshold BEFORE its sqrt(3/5) factor; its norm estimate is
% sqrt(64/3) exp[-(3/10)(13.85)^2]<5e-25. The final10^-24 margin holds.
% This is only an initial L2 tail, not a later graph or launch-tail bound.
%
% The inverse spacelike map's causal support gives the exact transformed
% input tail outside|x|<=.39 by the same10^-24: on the native annulus,
% |v|<11/13 and |x|+|g|=|x-g|<=(13/2)|r-1|. This keeps BOTH energy
% bands and coherences. Its potential-Morse coordinate
% zeta=sqrt3(r-1)sqrt(r+2)/r^(3/2) has
% d zeta/dx=sqrt3(3r-2)/(r^(7/2)sqrt(r+2)) and exactly the displayed
% partner potential. The Weyl symbol of ({a,p_zeta}/2)^2 includes
% +epsilon^2(a')^2/4, as checked by the second Moyal contraction. This
% is distinct from, and consistent with, the action chart above.
%
% The initial static tube lies in13/15<r<6/5, wider than the first real
% annulus used above. Extending that SAME action chart gives explicit
% margins on the wider annulus. These bounds were derived independently
% in reasoning before receiving the RED's matching extension:
% 2.25<=F<4.334, |F'|<16.667, |F''|<166.667,
% 2.35<b<3.64, |(log b)'|<1.56, |(log b)''|<1.94,
% |b'|<5.68, |b''|<15.93;
% 5.28<c=bF<15.78, |c'|<85.3, |c''|<866;
% 5.28<H<15.78, |H'|<170.6/3, |H''|<433.
% Hence |l_r|<6.95, |l_rr|<101, .59<a<1.59,
% |a_y|<3.1 and |a_yy|<52. These imply
% |R|<52/(2*.59)+3.1^2/(4*.59^2)<51,
% |a^2 R|<1.59^2*51<129.
% Thus the exact coordinate-ordering multiplier is<4.6e-13 throughout
% this initial tube at epsilon=2^-24. Initial localization is not used
% to delete the later cutoff residual or the retained detuning/partner
% and variable-kinetic terms. No propagated degree, root or run is claimed.
% SOURCES: complete actual original RED MP61/SOURCES, independently
% checked against Gaussian moments, Mehler, causal support and direct
% differentiation; the action-chart extension is the direct elementary
% derivative estimate stated here, not an imported propagator estimate.

% EXACT SEPARATED-ENERGY RESPONSE CONTROL [R, independent WIP].
% This uses the full continuum matrix operator and the original spacing
% insertion, without an adiabatic band identification or a radial detector.
% Let H be the fixed native continuum generator on the selected smooth
% core, and let P=1_[e0,E1](H), M=1_(-infinity,0](H), e0>0. Zero may be
% included in M; no global massless gap is assumed. Set H_P=H|P,H_M=H|M.
% The Cartesian velocity is exactly
% (i/epsilon)[H,q_i]=alpha_i+v_i, with norm<=1+sup|v|<2.
% Bounded monotone coordinate truncations with derivative<=1 have the same
% commutator bound. Their separated-spectral matrix blocks converge to the
% bounded off-diagonal position operator
% P q_i M=-i epsilon int_0^infinity exp(-sH_P)
%                   P(alpha_i+v_i)M exp(sH_M)ds,
% and hence ||P q_i M||<=2epsilon/e0. This follows by solving its exact
% Sylvester equation, not by dropping the Berry or opposite-band position.
% The bound is invariant under the exact Cauchy/polar maps already saved.
%
% For the original self-adjoint B, the saved graph estimate gives
% ||B P||<=1180 sum_(j=0)^3 E1^j=:b1. Thus X:M<-P solving
% H_M X-X H_P=M B P is the bounded operator
% X=-int_0^infinity exp(sH_M)M B P exp(-sH_P)ds,
% ||X||<=b1/e0. Starting from u(0)=P u(0), the part of the spacing response
% in M driven by this component has the EXACT form
% w_M,cross(t)=[exp(-itH_M/epsilon)X
%                  -X exp(-itH_P/epsilon)]u(0).
% The sign follows by differentiating exp[-i(t-s)H_M/epsilon]X
% exp[-isH_P/epsilon]. Its norm is<=2b1||u(0)||/e0 at every time; the
% off-diagonal source does not incur a t/epsilon norm bound here.
%
% Consequently its contribution to each Cartesian spacing-response
% expectation is bounded by
% |2Re<u_P(t),q_i w_M,cross(t)>|
%                    <=8epsilon b1||u(0)||^2/e0^2.
% At the selected epsilon, with the CONCRETE spectral block e0=1/2,E1=2,
% b1=17700 and this is<=566400epsilon||u(0)||^2
% <.033761||u(0)||^2 per component, or<.05848||u(0)||^2 in vector norm.
% Convex reflection averaging preserves the bounds. An initially given
% w_M(0) contributes at most4epsilon||u(0)|| ||w_M(0)||/e0 per component.
% These are evaluated bounds for the actual full-matrix interband source
% and detector, not an assertion that it vanishes or an asymptotic symbol.
%
% This decomposition does NOT replace the selected packet by P u_in.
% Its input outside[e0,E1], existing negative component, all their weighted
% response/launch terms and diagonal blocks still have to be retained or
% bounded. No estimate for those actual spectral tails is inferred from
% an initial spatial tail, nor is an inverse root-Jacobian bound supplied
% here. The calculation isolates one uniformly controlled contribution
% that a complete retained propagation can use; it does not evaluate the
% absolute matched-root coefficient. No new execution is requested.
% SOURCES: direct bounded-velocity commutator, separated self-adjoint
% spectral semigroups, exact Sylvester/Duhamel identity and the already
% proved full native1180 graph bound. This derivation was saved before any
% new counterpart separated-energy-response calculation was read.

% ACTUAL MP63--65 POLAR/RESPONSE COMPARISON [R, WIP].
% Read each complete actual original RED block and its SOURCES, including
% the saved plus-i conjugation repair. The relative form bound follows
% from K_L^2>=[L^2-70epsilon L]w^2 and square-root monotonicity; it does
% not replace w by its infimum or assert a massless energy gap. Inserting
% |K_L|^(1/2) in the resolvent quadratic integral gives pi/m_0, so the
% bounded-weight sign estimate and its separated-tau exponential follow
% without a spurious integral of1/|t|. The correct exact conjugation is
% exp(theta phi)p exp(-theta phi)=p+i epsilon theta phi_x. I requested
% this sign repair and reread its actual save; none of the norm estimates
% changes. The stated inherited endpoint domain and increasing saturated
% weights justify the polar maps and weighted-domain limits.
%
% MP64's r^n and<x>^n estimates follow from dr/dtau=r sqrt f and
% sqrt6 |x|(1+|x|)/(1+x^2)<3. The same weighted factor for the two-spin
% fold is exact Hadamard norm algebra, not a triangle bound with an extra
% sqrt2. The Neumann coefficients give n!(epsilon/.79)^n for the bounded
% iterated tau commutators. These statements apply to retained L>=.8;
% the previously bounded low-angular contribution stays in its ORIGINAL
% native representation. The full folded INPUT tail outside|x|=.49 is
% therefore<2e-24, retaining both signs. No later support or response
% graph bound is inferred from it.
%
% The MP65 vector improvement to the separated-energy result above is
% also exact. With V_i=alpha_i+v_i,
% sum_i V_i^dag V_i=(3+|v|^2)I+2alpha.v, hence
% c_v=sqrt[3+2sqrt(5/6)+5/6]<2.38. Spectral compression is contractive
% into the threefold sum, and bounded coordinate truncations preserve
% this velocity bound. Thus ||(Pq_iM psi)_i||<=epsilon c_v||psi||/e0,
% and the source-driven raw response vector is bounded by
% 4epsilon c_v b1||u0||^2/e0^2<.040175||u0||^2
% at e0=.5,E1=2,epsilon=2^-24. The initial opposite-sign response adds
% at most2epsilon c_v||u0|| ||w_M(0)||/e0. This improves the earlier
% component-to-vector conversion only; it does not supply the actual
% spectral tails, negative preparation, diagonal response or inverse
% four-return Jacobian. In particular it is not a final coefficient error.
% SOURCES: complete actual original RED MP63--65 and SOURCES, checked
% directly against the radial square, spectral quadratic integral,
% bounded-weight domain limit, exact Hadamard algebra and Clifford sum.
% The author's independent separated-energy derivation precedes this
% intake. No counterpart result, numerical execution or model change.

% EXACT FREE/PHYSICAL ENERGY-SIGN ANGLE [R, independent WIP].
% This controls the TOTAL negative-energy part of the ACTUAL free-band
% preparation; it does not replace that preparation by an energy projector.
% It was derived after reading the actual MP66 spatial-localization lemma,
% before reading any counterpart energy-sign-angle calculation. Use one
% regular radial block, with L=epsilon|kappa|>=.79, and write
% D_f=sigma_y p_r+L sigma_x/r, D=|D_f|, S=sign(D_f), A=v(r)sigma_y.
% A component exchange for negative kappa leaves every norm below unchanged.
% The native operator has the EXACT anticommutator identity
%                  H=D_f+{A,D_f}/2.                       (EA1)
% This includes v'/2; no radial momentum is independently made self-adjoint.
% The inherited regular full-Dirac domain is used throughout. The free
% square and Hardy bounds give
% ||p_r D^-1||<=1, ||r^-1 D^-1||<=[L(L-epsilon)]^-1/2.
% Inverse notation denotes the bounded closures of these products, not a
% bounded D^-1 or a gap at zero. The free kernel is zero in this domain.
%
% The existing global profile bound sup|rv'|<=91/3, |v|<=sqrt(5/6),
% suffices. Since D^2=p_r^2+L^2/r^2+epsilon L sigma_z/r^2,
% [D^2,A]=-i epsilon[{p_r,v'}sigma_y+2L v sigma_x/r^2].
% Therefore the antisymmetric relative form
% R=D^-1[D^2,A]D^-1 obeys
% ||R||<=epsilon[2(91/3)/sqrt(L(L-epsilon))
%                   +2sqrt(5/6)/(L-epsilon)]<80epsilon=:r_0. (EA2)
% Cauchy--Schwarz in the p_r/r graph proves this without differentiating
% v' again. All identities can first be cut to compact regular radial
% cores and spectral intervals of D, then extended by the displayed forms.
%
% Two exact divided-difference estimates are useful. First,
% M=D^-1/2[D,A]D^-1/2
%   =int_0^infinity D^1/2 exp(-sD) R exp(-sD)D^1/2 ds,
% so ||M||<=r_0/2 by int_0^infinity D exp(-2sD)ds=I/2.
% Second, define the self-adjoint relative anticommutator defect
% C=D^-1/2{D,A}D^-1/2/2-A.
% In the spectral representation of D it is the Schur multiplier
% C_(lambda,mu)=f(log lambda-log mu)R_(lambda,mu),
% f(t)=tanh(t/4)/[4cosh(t/2)]
%     =[csch(t/2)-2csch(t)]/4, with its removable value at t=0.
% With Fourier transform int exp(-ist)f(t)dt,
% fhat(s)=-(i pi/2)[tanh(pi s)-tanh(pi s/2)],
% and (2pi)^-1 int|fhat(s)|ds=log2/(2pi)<.111.
% Thus C=int fhat(s)D^(is)R D^(-is)ds/(2pi) is a norm-convergent
% operator integral and ||C||<.111r_0. This explicit decaying multiplier
% avoids assuming that the unbounded similarity D^1/2 A D^-1/2 alone
% has a small norm. No operator-Lipschitz assertion for an arbitrary
% discontinuous sign function is being used.
%
% Decompose(EA1) with the FREE projections P_0,+/-=(I+/-S)/2.
% Its diagonal operator H_d=P_0,+ H P_0,++P_0,- H P_0,- satisfies
% H_d,++>=d D_+, -H_d,-->=d D_-,
% d=1-sqrt(5/6)-.111r_0>.08699.
% Indeed its two relative diagonal forms are +(I+A_+++C_++) and
% -(I+A_--+C_--), respectively. The off-diagonal remainder is exactly
% H-H_d=D^1/2 E D^1/2, where E_+-=M_+-/2,
% E_-+=-M_-+/2. Hence E is self-adjoint and
% ||E||<=r_0/4=:eta<1.2e-6. In particular |H_d|>=dD and sign(H_d)=S.
%
% Apply the already proved quadratic-resolvent lemma with weight D,
% reference H_d and bounded middle perturbation E. Its two ingredients are
% ||D^1/2(H_d-it)^-1D^1/2||<=1/d and
% int_R||D^1/2(H_d-it)^-1u||^2dt<=pi||u||^2/d.
% The resolvent identity and eta/d<1 therefore give
% ||sign(H)-S||<=eta/(d-eta),
% ||P_E,- P_0,+||<=||P_E,--P_0,-||<=eta/[2(d-eta)]<7e-6.  (EA3)
% The argument is uniform at zero energy and all L>=.79; it uses no
% finite global gap. H_d is defined by its positive/negative closed forms,
% and the bounded relative-form resolvent identity identifies the perturbed
% operator with the actual self-adjoint native H on the common regular core.
%
% The prescribed input is exactly in the FREE positive band, including
% its scalar translation and cutoff. On J_z>=.8, all its radial blocks
% have L>=.8, so(EA3) applies to the full angular superposition and both
% reflection partners without discarding coherences. Its remaining native
% low-J_z norm is<8e-34. Consequently the FULL actual preparation obeys
%       ||P_E,- u_in||<7e-6,   ||P_E,-u_in||^2<4.9e-11.     (EA4)
% This probability is conserved by the exact fixed-background continuum
% propagator. The state is not renormalized or replaced by P_E,+u_in.
% Its off-diagonal and same-sign spacing-response contributions remain
% part of the binding observable; a small base probability alone does not
% bound a derivative containing the full B insertion or the root response.
%
% The actual MP66 spatial theorem, separately checked in full, is compatible
% with(EA4): |H_L|>=.068/r follows from its exact radial graph triangle,
% and the log-radius weighted resolvent gives EACH sign component a native
% initial tail<2e-24 outside[.9,1.1]. It supplies location rather than the
% norm in(EA4). No finite positive energy window or later-time trapping
% statement follows from combining these two inequalities.
% SOURCES: direct identity(EA1), free radial Hardy square and explicit
% spectral divided differences/Schur integral(EA2--3); the complete actual
% original RED MP66/SOURCES supplied and was checked only for the existing
% radial graph, relative-sign-resolvent and initial spatial-localization
% facts. The new angle estimate is independent and precedes counterpart
% intake. No numerical run, core/state change or coefficient claim.

% BOUNDED FREE-BAND COUPLING AND FIRST NEGATIVE GRAPH [R, independent WIP].
% The same actual profile gives a useful stronger domain statement for
% EA1--4. In addition to|v'|<=35 and|v/r|<=35, one has
% sup|r v''|<=7020: the core/join uses r<=13/15 and the saved Cartesian
% second-derivative bound8100 on an axis; the exterior formula is smaller.
% Reordering the exact Q=[D^2,A] from(EA2) gives
% Q=-2i epsilon v' p_r sigma_y-epsilon^2 v''sigma_y
%                                      -2i epsilon L v sigma_x/r^2.
% Thus, on L>=.79,
% ||Q D^-1||, ||D^-1 Q||
% <=epsilon[70+7020epsilon/sqrt(L(L-epsilon))
%                         +70sqrt(L/(L-epsilon))]<141epsilon.
% The second bound is the adjoint of the first. Three-lines interpolation
% of D^(-z)Q D^(z-1), first on positive compact spectral intervals and
% then by its uniform closures, gives
% ||D^-1/2 Q D^-1/2||<141epsilon.
% The exact Sylvester integral now implies the UNWEIGHTED bound
% ||[D,A]||<141epsilon/2,
% ||H-H_d||=||P_0,- H P_0,+||<141epsilon/4<36epsilon.        (EA5)
% In the first integral, Q=D^1/2 R_0 D^1/2 and
% int_0^infinity||D^1/2 exp(-sD)psi||^2ds=||psi||^2/2.
% No integration of a1/s operator norm is used. Equation(EA5) identifies
% H_d as the actual native H minus a bounded self-adjoint off-diagonal
% operator on the same regular domain, strengthening the form-domain
% description in(EA3).
%
% For the retained actual positive-free-band input phi, write
% chi=P_E,-phi. Energy-sign projection commutes with H, so exactly
% H chi=P_E,-P_0,+ H phi+P_E,-P_0,- H phi.
% EA3 and(EA5) therefore give
% ||H chi||<7e-6||H phi||+36epsilon||phi||.                  (EA6)
% The actually read MP67 bound ||(H-e_c)phi||<.0019 and e_c<1.02
% implies ||H phi||<1.0219, hence ||H chi||<9.4e-6 on the retained
% sector. Together with EA4 this gives its first energy graph<1.64e-5.
% This is much sharper than the negative-tail graph obtained from a
% single Chebyshev inequality. The saved low-angular first-graph bound
% is added separately when taking the full state; it is not assigned the
% retained radial constant. Higher graphs of chi or of the other spectral
% tails are still needed for the full cubic B response.
% SOURCES: independent ordered radial commutator, graph products and
% explicit interpolation/Sylvester integral; saved actual global profile
% derivative bounds. Only the final scalar input norm in(EA6) uses the
% complete actual original RED MP67, whose calculation was saved before
% the new author sign-angle route was read. No scientific execution.

% ACTUAL MP67 COMPLEMENTARY SPECTRAL-INPUT COMPARISON [R, WIP].
% Read the complete actual MP67 and SOURCES. Its identity for(H-e_c)phi
% keeps alpha.p phi=|p|phi exactly and retains the divergence, correlated
% delta-q/z term and actual centering. The mixed Gaussian second moment
% has the extra canonical six-term contribution; the ordered six-point
% bound and projection-derivative corrections also check. Directly, the
% radial exterior Hessian satisfies
% ||D^2v[d,d]||^2=(3/2)r^-5|d|^4[1+t^2-7t^4/4],
% t=(n.d)/|d|, whose maximum is(12/7)r^-5|d|^4. This confirms its
% .77 Taylor-remainder margin on r>.94. Its Mehler p and q-p tails
% control the exterior of the good coordinate box instead of extending
% that local Hessian bound through the core. The exact correlation
% ell_*K=0 retains the small linear-energy variance5epsilon/3 despite
% the squeezed packet. The sum<.0019, reference<.0008, and spectral
% window amplitudes<.004, reference<.0016, follow with the stated
% normalizers and cut allowances. The first energy graphs follow by the
% bounded spectral quotients. EA4 improves its separate negative norm;
% it does not improve the positive part outside[.5,2] or supply third
% graph tails. Both results concern the original state, not a filtered
% replacement, and neither evaluates the propagated response or root.
% SOURCES: complete actual original RED MP67/SOURCES, checked against
% exact free projection, ordered Gaussian pairings, the radial Hessian,
% Mehler and spectral theorem. Existing independent sign-angle EA1--6
% is the complementary author route; no extra run or comparison method.

% EXACT GAUGED OPPOSITE-SIGN RESPONSE ON FINITE ENERGY BLOCKS [R, WIP].
% This uses the already saved exact native-response gauge in the INERTIAL
% Cartesian frame, so H, T and B_eff=B+[H,T] are time independent. The
% physical preparation, native-centering term and explicit detector
% [q_i,T] remain as in that gauge. No moving-band or scalar replacement is
% made. Write B_eff=B_0+[F,T],
% B_0=-(S4/(6p))alpha.n-(1/3)sum_i{v_i,k_i^3},
% F=sum_i{v_i,k_i}/2, T=i Sigma.(k cross k^(3))/(12p^2).
%
% A direct free-sign resolvent estimate controls the scalar flow's
% opposite-free-band part. On retained angular sectors with L>=.79,
% |D_f|>=sqrt(L(L-epsilon))/r>.789/r. For f=v_i the exact commutator
% [D_f,f]=-i epsilon alpha.grad f has
% sup r|grad f|<=91/3<31: the radial derivative eigenvalues are rv'
% and v. The same quadratic-resolvent integral used in EA3 therefore
% gives ||[sign D_f,f]||<40epsilon between these buffered sectors, and
%                  ||P_0,- v_i P_0,+||<20epsilon.          (EG1)
% Here v_i has angular rank1 and k_i^3 rank<=3. The source changes J_z
% by at most4epsilon, so input J_z>=.8 and its source buffer stay above
% L=.79. The free resolvents commute with the angular projectors; the
% relative1/r estimate is applied on BOTH sides of the commutator form.
% No unjustified full-space angular Hardy inequality is used.
%
% Since every k_i commutes with the free projector, the free off-diagonal
% part of B_0 has the exact form bound
% |<chi,(B_0)_off psi>|
% <=12epsilon[||chi||T_3(psi)+T_3(chi)||psi||].              (EG2)
% For each ordered sign block this follows from(EG1), its adjoint and
% sum_i||k_i^3 psi||<=sqrt3 T_3(psi). Summing the two sign blocks uses
% two-component Cauchy--Schwarz and P_0 commutation with p^3, so no extra
% factor2 is required. The free kinetic term commutes with P_0 exactly.
%
% The full polynomial part also has the global L2 bound
% ||B_0 psi||<=.78T_3+61epsilon T_2+14030epsilon^2 T_1
%                                      +3800000epsilon^3 T_0. (EG3)
% This is its exact left-ordered expression: the leading scalar flow
% is bounded by(2/3)|v|p^3, while its three derivative coefficients are
% the saved M1,M2,M3. The kinetic term contributes at most p^3/6.
% For a vector in a physical spectral interval|H|<=2, the proved graph
% constants give T_1<34.434, T_2<922.25, T_3<22695, hence(EG3)<17800
% times its norm. These inequalities are used for estimates only; the
% actual physical cubic source is still retained.
%
% Let P=1_[.5,2](H) on J_z>=.8, and M=1_[-2,0](H) on its rank4
% angular buffer. For normalized psi=Ppsi and chi=Mchi, EA3 bounds the
% wrong FREE-sign component of each by7e-6. The two free-diagonal B_0
% pairings therefore cost at most2(7e-6)(17800)=.2492. Equation(EG2)
% costs at most24epsilon(22695)<.0325. The independently saved exact
% form bound for[F,T] gives, on the SAME two spectral blocks,
% (sqrt3/2){100epsilon(T_2+T_1^2)
%                         +4700000epsilon^2(2T_1)}<.011.
% It follows, as a bounded rectangular operator defined by this form,
%                         ||M B_eff P||<.3.                 (EG4)
% This is a full-matrix estimate including the flow commutator. It is
% much smaller than the un-gauged norm bound17700 because the exact P_2
% gauge cancels the leading opposite-free-band kinetic source.
%
% The SAME exact Sylvester/Duhamel calculation as before, now with(EG4),
% bounds the M-response driven by a P base at every physical time. The
% corresponding raw Cartesian response-vector contribution is
% <=4epsilon c_v(.3)/(.5)^2 ||P u_in||^2
%                    <7e-7||P u_in||^2, c_v<2.38.          (EG5)
% This is the z_resp contribution in the exact gauge, not the complete
% physical derivative. Its initial M z_resp(0), explicit [q,T] detector,
% negative base, all blocks outside[-2,2], the positive spectral tails,
% the diagonal response and the actual matched-root Jacobian still enter.
% In particular the actual input is NOT replaced by P u_in and the
% saved .004 spectral tail alone is not enough for its cubic-source graph.
% Equation(EG5) evaluates one necessary full-matrix component and does not
% complete a retained-propagation or root certificate. No run is requested.
% SOURCES: direct free-sign resolvent/Hardy form(EG1), exact scalar-flow
% anticommutator and ordered native source, EA3, the already derived
% global[F,T] form and full native graph bounds, and exact separated-
% energy Sylvester identity. This block was independently saved before
% reading a counterpart gauged finite-energy-response argument.

% GLOBAL TWO-BAND BASE-PROPAGATOR REORGANIZATION [R, independent WIP].
% This uses the exact bounded FREE-band coupling in EA5, not the failed
% first moving-Gaussian amplitude truncation. It supplies a genuine
% finite-time L2 operator estimate while leaving the retained scalar
% propagators, weighted response and return root to be evaluated. Fix a
% retained native angular sector L>=.79. In the free-sign decomposition,
% write H=[[A_+,V],[V^dag,-A_-]], with A_+,A_->=dD_+,dD_-, respectively,
% d>.08699, ||V||<36epsilon. Its off-diagonal relative form has norm
% eta<1.2e-6, and set rho=eta/d<1.38e-5. The symbol rho here is NOT a
% matter density or the earlier Gaussian Mehler parameter.
%
% The actual positive spectral subspace is the graph of X:P_0,+->P_0,-.
% The direct-rotation theorem here can be proved explicitly from the
% projection blocks: P_E,++ is invertible, X=P_E,-+(P_E,++)^-1, and
% ||X||<=||P_E,+-P_0,+||/sqrt(1-||P_E,+-P_0,+||^2).
% In fact EA3 gives the sharper projection bound
% 1.2e-6/[2(.08699-1.2e-6)]<6.9e-6, hence ||X||<7e-6.
% All blocks are on the regular native domain.
% Expand the sign resolvent about H_d=diag(A_+,-A_-). Its nth integrated
% Neumann term has norm<=rho^n: the two outer quadratic-resolvent
% integrals give pi/d and every middle relative resolvent gives1/d.
% Each insertion of V changes free sign. Hence odd terms are off-diagonal
% and even terms are diagonal, also without a zero-energy gap.
%
% Define the bounded Sylvester integral
% X_1=int_0^infinity exp(-sA_-) V^dag exp(-sA_+)ds,
% so A_-X_1+X_1A_+=V^dag and ||X_1||<=rho/2.
% Its integral is defined by the same relative-form quadratic estimate,
% not by integrating a divergent uniform small-energy resolvent bound.
% The first off-diagonal positive-projection term is exactly X_1. If
% a=rho^2/[2(1-rho^2)] and b=rho^3/[2(1-rho^2)], the parity expansion
% gives ||P_E,++-I||<=a and ||P_E,-+-X_1||<=b. Therefore
% ||X-X_1||<=[b+(rho/2)a]/(1-a)<3e-15.                    (ER1)
% This is a norm-convergent relative-coupling expansion of the ACTUAL
% spectral subspace, rather than a formal series in semiclassical order.
%
% For any bounded X define the EXACT unitary graph rotation
% U(X)=[[R_+,-X^dag R_-],[X R_+,R_-]],
% R_+=(I+X^dag X)^-1/2, R_-=(I+XX^dag)^-1/2.
% Elementary positive-resolvent differentiation gives, on this small ball,
% ||U(X)-U(X_1)||<=4||X-X_1||<1.2e-14.
% The actual U(X) diagonalizes H. Its positive effective generator is
% K_+=(I+X^dag X)^1/2(A_++VX)(I+X^dag X)^-1/2.
% The Riccati identity
% -A_-X-XA_++V^dag-XVX=0
% controls the otherwise unbounded normalization commutator. Put
% T=X^dag X, S_+=(I+T)^1/2 and B=VX. Then
% [T,A_+]=B^dag(I+T)-(I+T)B,
% Y=[S_+,A_+] solves S_+Y+YS_+=[T,A_+], and
% K_+-A_+=Y S_+^-1+S_+B S_+^-1.
% Thus it is a bounded self-adjoint operator on the A_+ domain. With
% x=||X|| and b_0=||B||, comparison to Y_0=(B^dag-B)/2 gives
% ||Y-Y_0||<=1.5 b_0 x^2 and
% ||K_+-A_+-(B+B^dag)/2||<=3||V||x^3.                    (ER2)
% The same proof holds for the negative effective block, with the
% corresponding minus sign. It keeps the normalization commutator and
% does not replace it by zero because X is small.
%
% Consequently the explicitly defined TWO-band operator
% K_2=diag(A_++C_+,-A_-+C_-),
% C_+=(V X_1+X_1^dag V^dag)/2,
% C_-=-(V^dag X_1^dag+X_1 V)/2
% is self-adjoint and satisfies
% ||U(X)^dag H U(X)-K_2||
% <=||V||[3x^3+||X-X_1||]<1.46e-13 epsilon.               (ER3)
% Both C blocks are bounded; their norms are<36epsilon(rho/2).
% Thus the interaction-picture Dyson series for K_2 about H_d has
% 12||diag(C_+,C_-)||/epsilon<.003. Retaining through THREE such bounded
% insertions has an L2 operator remainder at most
% exp(.003)(.003)^4/24<3.4e-12. These are finite integrals of the EXACT
% A_+,A_- propagators and C blocks; no temporal or radial discretization
% error is included or presumed zero. They offer a finite convergent
% evaluation once those retained propagators and X_1 are actually bounded.
% Duhamel and(ER1--3) now give the evaluated finite-member statement
% sup_(0<=t<=12)||exp(-itH/epsilon)
%       -U(X_1)exp(-itK_2/epsilon)U(X_1)^dag||<2e-12.       (ER4)
% This holds for ALL retained inputs, not only the central Gaussian, and
% does not assume that their later support stays in a local radial annulus.
% The actual low-angular input remains in its prior native tail treatment.
%
% No time-integrator or finite basis has been selected here. A_+ and A_-
% are the exact free-sign compressions of the original radial operator;
% their scalar radial propagators, X_1 and the bounded C blocks still need
% a controlled retained evaluation. The original input is transformed by
% U(X_1)^dag, so its lower component is present even though it began in
% the FREE upper band. The original Cartesian detector and complete B_eff
% are transformed by the SAME fixed unitary; native-centering and [q,T]
% remain. In particular(ER4) is an L2 BASE estimate. An unbounded-coordinate
% or cubic-source graph error, its launch derivatives, and a spacing-
% derivative estimate do not follow by differentiating that number. Those
% weighted/graph bounds and the original four-return root remain required.
% This reorganization supplies a convergent two-band base representation,
% not a band deletion or an evaluated absolute MP8 coefficient.
% SOURCES: independent parity-resolvent expansion of EA3, exact graph
% projection/Riccati algebra and bounded Sylvester identities, using the
% already proved EA5 off-diagonal norm. Actual input, core, spacing gauge,
% physical detector and return protocol are unchanged. This derivation
% was saved before reading any counterpart direct-rotation continuation;
% no new execution, model assumption or numerical error was imported.

% ACTUAL MP68--70 GRAPH/ROTATION COMPARISON [R, WIP].
% Complete actual original RED MP68--70/SOURCES were read after the
% independent EA5--6, EG and ER saves. The second radial Hardy identity
% retains its c(c-6epsilon^2) term; the resulting .708 and1.000001
% constants and the ordered p_rQ/(L/r)Q derivative margins check. The
% Sylvester graph recurrence D M=Q-MD and D^2M=DQ-QD+MD^2 therefore
% gives the stated(36,106,226)epsilon bounds for the actual free-band
% coupling through two D graphs. Its H^2/H^3 commutator expansion keeps
% both H^2C and HCH; using the actual Gaussian/projector moments yields
% the full negative input E3<.0029 and ||B chi||<3.43. These are genuine
% graph bounds, not an omitted-response tolerance. The physical/free
% weighted-sign commutators likewise give retained||r chi||<8.1e-6 and
% ||r U(t)chi||<.000209 through t12. The small-L part remains in its
% unsplit native treatment, exactly as in that source.
%
% MP69 removes the earlier form-only limitation for the exact gauge:
% [F,T]=sum_i[v_i,T]p_i-(i epsilon/2)sum_i[partial_i v_i,T].
% Apply ||Tpsi||<=T2/24 directly to the last commutator, rather than
% applying a three-derivative kernel estimate to partial_i v_i. The product
% rule then proves the actual H2->L2 estimate
% ||[F,T]psi||<=178epsilon T2+8150000epsilon^2 T1
%                                      +712500epsilon^3 T0.
% The coefficients follow from100sqrt3+35/8 and
% sqrt3(4700000+8100/8), with the stated M3 last term. This checks with
% only the already bounded first THREE profile derivatives. On|H|<=2
% it gives<.0098 and may replace EG's .011 form allowance without
% changing its rounded .3 source or7e-7 response-vector bound.
%
% MP70's independent direct rotation is exactly U(X) in ER, written as
% a projection formula. Its domain argument uses bounded[H,P0] and the
% convergent projection series. Its sharper equation
% N C_++C_+N=N V X+X^dag V^dag N
% gives ||C_+||<=||V||delta/(1-delta^2); thus its .003025 finite-time
% comparison omits the bounded diagonal correction but retains the exact
% rotation. ER instead approximates that correction by its actual
% Sylvester/parity term and bounds the remaining L2 error. The two
% estimates therefore have different retained operators and must not be
% compared as errors of the same approximation. Both keep both bands and
% leave weighted response/launch errors and scalar propagation/root open.
% SOURCES: complete actual original RED MP68--70/SOURCES, independently
% checked ordered radial integration by parts, commutators, matrix norms
% and Riccati/Sylvester identities. All are same-pair[R] controls, with
% no new numerical task, other-lane agreement or coefficient claim.

% EXPLICIT SCALAR RADIAL KERNELS FOR ER1--4 [R, independent WIP].
% This makes the retained diagonal operators literal without changing the
% actual packet or replacing its Cartesian detector by a radial one. For
% kappa>0, L=epsilon kappa, put nu=kappa+1/2 and, for p>0,
% b_nu(p,r)=sqrt(p r) J_nu(p r/epsilon)/epsilon.
% The Hankel closure is int_0^infinity b_nu(p,r)b_nu(p',r)dr=delta(p-p').
% The regular free radial eigenspinors are exactly
% e_+(p,r)=2^-1/2(b_nu(p,r),b_(nu-1)(p,r))^T,
% e_-(p,r)=2^-1/2(b_nu(p,r),-b_(nu-1)(p,r))^T.
% The Bessel recurrences give (epsilon partial_r+L/r)b_nu=p b_(nu-1)
% and its adjoint relation, so their eigenvalues are+p and-p. Each
% scalar Hankel map is unitary, hence the joint two-band transform is
% unitary as well; the positive/negative cross closure vanishes. Negative
% kappa uses the already stated regular radial component exchange.
%
% On smooth compact spectral test functions define the real-basis flow
% kernel by its exact symmetric differential form
% F_nu(p,p')=-(i epsilon/2)int_0^infinity v(r)
%   [b_nu(p,r)partial_r b_nu(p',r)
%        -(partial_r b_nu(p,r))b_nu(p',r)]dr.
% It extends on the inherited flow graph domain; the displayed generalized
% kernel need not be absolutely integrable at every pair(p,p'). Its
% integration-by-parts form includes v'/2 and has no center boundary term
% on the regular domain. Put F_d=(F_nu+F_(nu-1))/2 and
% V=(F_nu-F_(nu-1))/2. The operators used in ER1--4 are exactly
%                    A_+=p+F_d, A_-=p-F_d,                (RK1)
% and the off-diagonal block is V=V^dag. No fitted local dispersion or
% principal-symbol matrix element is substituted for these kernels.
%
% Since every b_nu is real, F_d and V are purely imaginary Hermitian
% operators in this spectral basis. For scalar conjugation C,
% A_-=C A_+ C, C V C=-V. Consequently
% exp(+itA_-/epsilon)=C exp(-itA_+/epsilon)C.               (RK2)
% Thus ONE scalar radial propagator and its conjugate supply both exact
% diagonal bands. This is an antiunitary symmetry, not omission of the
% negative component. In the native radial representation it is simply
% (sigma_z C)H(sigma_z C)^-1=-H for this massless real radial flow.
%
% The same symmetry is retained by the bounded ER correction. Transpose
% in the real Hankel basis gives A_+^T=A_-, V^T=-V and X_1^T=-X_1.
% Hence C_-=-C C_+ C and the two K_2 blocks are negatives of conjugates.
% The lower transformed input remains present, and every full Cartesian
% coordinate and anisotropic spacing-source matrix element retains its
% actual angular coupling and band indices. Neither their off-diagonal
% pieces nor the frequency/centering rows follow from a scalar radial
% expectation alone.
%
% The actual normalized input components in this representation are the
% integrals of e_+^dag and e_-^dag against its original spinor spherical-
% harmonic coefficients, with all j,m coherences and the prescribed
% reflection mixture. At the FREE input stage the negative coefficient is
% exactly zero; U(X_1)^dag then creates its specified small lower component.
% These identities do not select a radial/angular quadrature or discard
% large carrier labels. A controlled scalar evaluation of(RK1), the
% Sylvester integral, source/readout matrix elements and the original
% weighted/launch/root error still has to be supplied. No execution is
% requested from the kernels or their formal generalized integrals alone.
% SOURCES: direct regular radial Bessel recurrences and Hankel closure,
% integration by parts of the DECLARED symmetric flow, and explicit
% conjugation/transpose algebra. This construction was independently saved
% before a counterpart scalar-kernel derivation was read. The full-matrix
% ER reorganization, preparation and detector are retained without change.

% ACTUAL MP71--72 COMPARISON [R, WIP].
% The complete actual original RED MP71--72/SOURCES were read. MP71 checks
% ER's relative-resolvent parity, normalization commutator, regular domain
% and distinct optional Dyson remainder with no repair. MP72's quadratic
% norm was checked directly: tr(A+C)=31/6, ||A-C||_F^2=1001/36 and
% ||B+B^T||_F^2=6 give3395/576. Its third directional-flow derivative,
% cubic kinetic remainder and Gaussian moments give the displayed
% 34+18+10+.005 allowances in epsilon units. The total-spin contribution,
% actual cutoff/centering and the bounded36epsilon physical free-band
% coupling are all retained. Thus its reference-input rotating width
% ||(H-J_z/sqrt3)u||<101epsilon is valid. It is not a shifted-launch
% statement, a finite spectral support or a propagated-error certificate.
% SOURCES: actual complete MP71--72/SOURCES and the explicit independent
% contractions and Taylor-margin checks just stated. RK1--2 separately
% received the original RED's actual-source Hankel/domain/symmetry check;
% the source records that comparison in its own existing file.

% GRAPH-CONTROLLED FINITE ROTATION AND THE FULL RAW RESPONSE [R, WIP].
% The following strengthens the L2-only ER representation without changing
% the retained scalar problem. Both physical signs, every angular index,
% the actual preparation, native B and Cartesian readout remain. All
% statements use the buffered angular sectors L>=.79; the saved small-Jz
% input/launch tails are still treated separately in the native theory.
% Write H_0=diag(A_+,-A_-), mathcal V=H-H_0 and
% G=I+|H_0|=diag(G_+,G_-), G_+/-=I+A_+/-. Thus G is a positive graph
% weight, not Newton's constant. For either orientation of the off-diagonal
% block V, MP68 and the elementary graph comparisons below give
% ||G_-^j V^dag G_+^-j||, ||G_+^j V G_-^-j|| <=v_j,
%       (v_0,v_1,v_2)=(36,6000,1000000)epsilon.             (GR1)
% To verify these conservative numbers put Dbar=I+D. The existing native
% graph estimates give ||G Dbar^-1||<3 and ||Dbar G^-1||<13.
% The ordered H^2 estimate and the free-block compression give
% ||G^2 Dbar^-2||<6. Conversely, on G^-2,
% ||D G^-2||<=11.478[1/4+88.5epsilon]<2.870,
% ||H G^-2||<=1/4+36epsilon,
% ||H^2 G^-2||<=1+(36epsilon)^2+9epsilon
%                         +203epsilon(2.870)+1890epsilon^2.
% Consequently131.75 times the sum of the last two norms and1, followed
% by2||D G^-2||+1, is<305: ||Dbar^2 G^-2||<305.
% Here the lower H^2 block is V^dag A_+-A_-V^dag, so its203epsilon
% term is retained. Finally
% ||Dbar^j V Dbar^-j|| <=(36,142,474)_j epsilon.
% Thus3*142*13<6000 and6*474*305<1000000 prove(GR1).

% Introduce a COMPLEX coupling parameter z only for the convergent
% expansion H(z)=H_0+z mathcal V; it is neither the spacing derivative
% nor a physical parameter change. On |z|<=R=1000 the relative norm is
% <.0138. The same analytic sign-resolvent projection defines a graph
% X(z), odd in z. Its off-diagonal/diagonal parity estimates give
% ||X(z)|| <= [.0138/(2(1-.0138^2))]
%                 /[1-.0138^2/(2(1-.0138^2))]<.007.
% The Riccati identity is
% A_-X+X A_+=z[V^dag-X V X].                              (GR2)
% It also supplies graph regularity, not just a weak form identity.
% For S solving A_-S+S A_+=R_S,
% G_-S=S(2-G_+)+R_S, and ||(2-G_+)G_+^-1||<=1.
% Repeated application, with(GR1), gives the uniform analytic bounds
% ||G_-^j X(z)G_+^-j||<=M_j,
%                 (M_0,M_1,M_2,M_3)=(.007,.01,.37,69).
% Indeed x_(j+1)<=x_j+1000 v_j(1+x_j^2) for j=0,1,2.
% The same bounds hold for the opposite orientation. Spectral-domain
% cutoffs followed by this recurrence justify the closed graph extension;
% their uniform limits are holomorphic in these graph operator norms.

% Keep the FIVE odd coefficients through degree9,
% Y=X_1+X_3+X_5+X_7+X_9, with
% A_-X_1+X_1A_+=V^dag,
% A_-X_n+X_nA_+=-sum_(a+b=n-1) X_a V X_b (n=3,5,7,9).
% All a,b in the sum are positive odd indices. Their unique bounded
% relative-form solutions are the coefficients of(GR2), and Cauchy's
% inequality gives ||G_-^j X_n G_+^-j||<=M_j 1000^-n.
% This is an exact finite operator prescription; no scalar propagator or
% Sylvester integral has yet been discretized. Set U_9=U(Y), the EXACT
% unitary graph rotation with its full positive normalizers, and define
% H_9=U_9 diag(U_9^dag H U_9) U_9^dag. The discarded off-diagonal block
% is R_- mathcal R R_+, where R_+/- are those normalizers and
% mathcal R=V^dag-A_-Y-YA_+-Y V Y.
% Terms through total coupling degree9 cancel EXACTLY. The residual has
% only the pairs a+b>=10 among a,b=1,3,5,7,9. Hence
% ||G_-^j mathcal R G_+^-j||
% <=5v_j M_j^2 1000^-10/(1-1000^-2), j=0,1,2.
% The graph normalizers are controlled by their convergent power series:
% ||G^j U_9 G^-j|| <=(1+y_j)/sqrt(1-y_j^2),
% y_j=M_j/[1000(1-1000^-2)]. Including both outer rotations gives
% E=H-H_9 self-adjoint, bounded, and
% ||E||<5.4e-40,
% ||G E G^-1||<1.9e-37,
% ||G^2 E G^-2||<4.2e-32.                                (GR3)
% In particular H_9 has the original regular domain. It is an actual
% two-band self-adjoint operator, not a nonunitary truncated rotation.
% Its two diagonal blocks still obey RK's conjugation symmetry, since
% every X_n^T=-X_n. One scalar retained propagation therefore suffices,
% with the full lower input and all transformed observables preserved.

% The graph propagation constants do not require exp(t v_2/epsilon).
% For K=H_0+W with ||W||<=w_0 and ||G W G^-1||<=w_1, direct expansion
% of H_0^2=(K-W)^2 yields
% ||G^2 psi||<=||K^2 psi||+(2+w_0+w_1)||K psi||
%       +[1+2w_0+2w_0^2+w_1(1+w_0)]||psi||.
% Also ||K G^-2||<=1/4+w_0 and
% ||K^2 G^-2||<=1+w_1+w_0/4+w_0^2.
% Conservation of the K graphs shows ||G^2 exp(-itK/epsilon)G^-2||<3
% for BOTH K=H,H_9, with w_0<36epsilon+1e-39 and
% w_1<6000epsilon+1e-36. The triangle bounds give<2.502, so3 is strict.
% Duhamel, (GR3), and interpolation then prove, uniformly|t|<=12,
% eta_0=||U_H(t)-U_H9(t)||<1.1e-31,
% eta_2=||G^2[U_H(t)-U_H9(t)]G^-2||<1e-22,
% eta_1=||G[U_H(t)-U_H9(t)]G^-1||<4e-27.                (GR4)
% These compare exact propagators. They do NOT include the still absent
% controlled scalar evaluation of H_9 or its finite Sylvester coefficients.

% There is a useful unbounded-position-free readout of this same physical
% observable. With V_i=i[H,q_i]/epsilon=alpha_i+v_i,
% q_i(t)=q_i(0)+int_0^t <u(s),V_i u(s)>ds. It is an EXACT Heisenberg
% identity, not a new detector. Use the PHYSICAL V_i in the approximate
% U_H9 path as well. Since ||(V_i psi)_i||<2.38||psi||, (GR4) bounds
% the resulting base centroid-vector error by<6.3e-30 through time12.
% Its R,p_r,L columns follow by the same bilinear identity on the actual
% launch tangents: their saved norm<2e8 gives a conservative<3e-21
% raw centroid-column error. The physical velocity supplies the time row;
% the original mean-frequency and its spacing derivative remain exact
% INITIAL readouts, not expectations of the auxiliary H_9.

% The full NATIVE spacing response can likewise be read without discarding
% a band or a rapid phase. Keep the original w(0) and B, and set
% V_B,i=i[B,q_i]/epsilon=-alpha_i k_i^2/2-{v_i,k_i^2}.
% The derivative of the velocity identity contains the initial w(0)
% bilinear, int<u,V_B u>, and the exact double integral
% (2/epsilon) Im int_0^t ds int_0^s dt'
%   < U_H(s-t')^dag V_i u(s), B u(t') >.                 (GR5)
% The sign follows directly from w'= -i(Hw+Bu)/epsilon.
% Use the same B,V_B,w(0) in the H_9 approximation. A single integration
% by parts in the ordered B form proves
% |<chi,Bpsi>| <= .78 T1(chi)T2(psi)
%  +(sqrt3 M1/3)epsilon||chi||T2(psi)
%  +(2M1/3)epsilon T1(chi)T1(psi)
%  +(sqrt3 M2/3)epsilon^2 T1(chi)||psi||
%                    <3100||G chi||||G^2 psi||.
% The same elementary product bounds give
% ||(G V_i psi)_i||<100||G psi|| and
% ||(V_B,i psi)_i||<800||G^2 psi||.
% The actual launch-box input obeys ||G u(0)||<7,
% ||G^2 u(0)||<50, and ||w(0)||<.753. Thus the primal second graph is
% <150 and the vector dual first graph in(GR5) is<2100. Their H_9
% differences are respectively<5e-21 and<1e-23, using(GR4).
% The double-integrand vector error is therefore
% <3100[(1e-23)(150)+(2100)(5e-21)]<3.26e-14.
% Multiplication by144/epsilon, with the smaller direct/initial terms
% included, bounds the FULL raw centroid-spacing-response error from
% this rotation replacement by<8e-5 on0<=t<=12.            (GR6)
% This is a physical full-matrix response estimate; no small occupation
% argument or differentiated L2-only error was used. The cutoff/centering
% and both-band input are the actual ones, and buffers include all source
% and coordinate angular shifts. The saved low-angular errors are separate.
% Radius/return combinations can use these raw bounds once their actual
% radius and Jacobian enclosures are supplied. (GR6) is not the absolute
% MP8 coefficient or an error bound after multiplication by an unmeasured
% inverse return Jacobian. Finite scalar propagation, all transformed
% matrix elements and the actual four-condition root remain to evaluate.
% No new run is requested from this finite representation alone.
% SOURCES: independent analytic Riccati/Cauchy graph continuation of ER,
% MP68's already checked D-graph coupling bounds, explicit ordered native
% source forms and the exact physical velocity/Duhamel identities. This
% full graph/readout construction was saved before reading a counterpart
% higher-rotation or physical-velocity response certificate.

% EXACT CRITICAL-PHASE RADIAL CHART FOR RETAINED EVALUATION [R, WIP].
% This gives an explicit phase-adapted representation for the scalar/full
% radial evaluation still missing in GR. It is not a further truncation
% or an assertion that the propagated packet stays in the stated annulus.
% In the Schwarzschild annulus .9<=r<=1.1 let f=1-2/(3r),
% v=-sqrt(2/(3r)), L=epsilon|kappa|, E_c=L/sqrt3 and
% g(r)=sqrt((r+2)/r^3). The exact identity
%                   1/3-f/r^2=(r-1)^2(r+2)/(3r^3)
% gives TWO analytic critical radial momenta, indexed s=+1,-1,
% P_s(r)=E_c[-v+s(r-1)g]/f,
% e_s(r)=sqrt(P_s^2+L^2/r^2)
%       =E_c[1-s v(r-1)g]/f,
% e_s+vP_s=E_c,
% a_s(r)=v+P_s/e_s=s f(r-1)g/[1-s v(r-1)g].            (CP1)
% Thus a_s'(1)=s/sqrt3 exactly. These momenta are the full nonlinear
% critical branches; replacing them by their tangent quadratic phase
% would reintroduce the large first-amplitude effect already measured.
% Put S_s(r)=int_1^r P_s(t)dt. Multiplication by exp(iS_s/epsilon)
% is unitary, and shifts p_r to p_r+P_s EXACTLY.

% The full matrix spin connection is explicit. Set
% phi_s=atan(P_s r/L) and
% U_s=exp(-i phi_s sigma_z/2)exp(-i pi sigma_y/4).
% It obeys U_s^dag(P_s sigma_y+L sigma_x/r)U_s=e_s sigma_z.
% With c_s=L/(r e_s),
% A_s=vI+(P_s/e_s)sigma_z+c_s sigma_y,
% -i U_s^dag U_s'=phi_s' sigma_x/2. Therefore
% U_s^dag exp(-iS_s/epsilon)(H_kappa-E_c)
%            exp(iS_s/epsilon)U_s
%   ={A_s,p_r}/2+diag(0,-2e_s)
%                       +epsilon v phi_s' sigma_x/2.     (CP2)
% The last term follows from the anticommutator of the transformed
% velocity matrix with the spin connection; it is not zero because the
% flow is scalar. Both components and this term must be retained in an
% actual approximation. A principal elimination alone would yield an
% upper p_r^2 coefficient c_s^2/(2e_s), equal to1/18 at the reference,
% but CP2 does not replace the matrix operator by that scalar expression.

% The critical upper transport can be straightened without approximation.
% Let lambda=1/sqrt3, define
% Acal(r)=lambda/[f(r)g(r)],
% x_s(r)=(r-1) exp int_1^r
%      [(Acal(t)-1)/(t-1)-s lambda v(t)/f(t)]dt.
% The apparent quotient at1 is removable, with Acal'(1)=-2/3.
% Then x_s(1)=0, x_s'(1)=1 and a_s x_s'=s lambda x_s.
% The usual unitary half-density change therefore sends the upper
% diagonal transport {a_s,p_r}/2 EXACTLY to
%                         s lambda{x_s,p_x}/2.            (CP3)
% All other terms of(CP2), the source and the Cartesian detector transform
% along with it. In particular a switch of sign chart is its exact unitary
% change of coordinate, spin frame and phase, not a new packet preparation.

% Elementary real bounds on this annulus make the coordinate nonsingular:
% 1.526<g<1.995, .259<f<.394, |v|<.861 and
% |v(r-1)g|<.172. Thus1.21L<e_s<2.62L and the local matrix gap
% 2e_s>1.91 on L>=.79. The logarithmic derivative of Acal has the
% numerator3r^2+5r-10 over r(3r-2)(r+2), so Acal decreases on this
% interval. Its exact endpoint values give .95<Acal<1.12; also
% |(log Acal)'|<2.186, hence |Acal'|<2.449.
% The integrand defining x_s has absolute value<4.37; its exponent
% is<.437 in absolute value. Using
% x_s'=exp(the integral) Acal[1-s v(r-1)g]
% gives the conservative positive bounds .5<x_s'<2.1. These checks
% concern a real coordinate chart; no complex analyticity radius or
% later-time probability bound is asserted from them.

% This chart is fixed by each angular channel and the background. Changing
% the carrier labels in the MP8 preparation does not change that channel's
% operator. The physical spacing source is conjugated by the SAME fixed
% phase/frame/half-density map, so all its endpoint-ordering and native
% preparation terms remain. Between distinct angular channels, Cartesian
% coordinate/source blocks keep both channel phases and spin frames.
% A smooth global extension may be used to define the unitary, but(CP1--3)
% are only asserted in the annulus. Any finite approximate propagation
% must compute its actual equation residual and the chart-cutoff flux;
% the initial annular tail alone does not control those later errors.
% This is a concrete conditioned representation for that computation,
% not a new execution specification or a completed return calculation.
% SOURCES: direct factorization of the exact Schwarzschild radial energy
% surface, explicit Pauli rotations and symmetric first-order operator
% conjugation, and exact characteristic linearization. The construction
% was saved before reading any counterpart critical-phase chart argument.

% ACTUAL MP73--76 COMPARISON AND LOW-ENERGY LAUNCH COLUMNS [R, WIP].
% The complete actual original RED MP73--76/SOURCES were read. MP73 and
% MP75 clear RK/GR's actual kernels, domains and full native readout.
% MP74's exact two-root identity follows from H=M+alpha.z+S and
% Q(M)=0; its935/8 graph sum and outer-interval monotonicity check.
% Its higher moments are still unevaluated, so it supplies no high-energy
% truncation. MP76's squeezed quadratic row sum3200 gives the actual
%64th-moment factor<.069. The static interval is wholly exterior:
% x(.87)<-.45<x(1.2), and U>=25/81 there. Its heat/Schur estimate has
% the required square root from the row/column test and the correct
% exp(b E_*^2) spectral factor; both scalar partners are nonnegative on
% their inherited regular domains. The actual full-band low-energy input
% norm<2e-24, E3<4e-24 and raw spacing-response omission<5e-8 follow
% with the original initial rows and unrestricted receiving energies.
% No physical sign was deleted and no new normalized packet was introduced.

% The following extension bounds the missing BASELINE LAUNCH columns for
% that same omission. The fixed projector P_l=1_(|H|<=1/2) is independent
% of launch, so partial_j(P_l u)=P_l partial_j u. Use the original
% coordinates(R,k_cx,k_cy) and the already established actual uncut
% tangent form P0 f_j, where f_j is degree<=1 on the fixed translated
% Gaussian and ||f_j||<K=2e8. Its exact normalization, spinor derivative
% and actual centering derivative are included in f_j; the differentiated
% cut comparison remains<1e16 exp(-128) in L2.

% Extend MP61's degree79 A_39 approximation to that same tangent. In
% normalized Gaussian oscillator coordinates, write
% f_j=(c_0+sum_i c_i a_i^dag)g with |c_0|^2+sum_i|c_i|^2<=K^2,
% allowing spinor coefficients. Pointwise Cauchy--Schwarz gives
% |f_j(xi)|^2<=K^2(1+2|xi|^2)|g(xi)|^2. Consequently, for0<=l<=79,
% |||xi|^l f_j||<=K sqrt(2l+4)|||xi|^l g||
%                                  <=K sqrt162|||xi|^l g||.
% The extra Hermite degree costs another factor2 at rho=1/4. The SAME
% termwise estimate used for MP61's2.88 therefore gives
% ||rho^(-Nosc/2)A_39 f_j||
%             <(1.02)(2)(sqrt162)(2.88)K<75K=1.5e10.      (LC1)
% The omitted projector binomial remainder has the same moment estimate
% with degree<=81, costing less than13K in L2. Its previously proved
% bound and the actual cut/normalizer comparison give a conservative
% total tangent-approximation error<1e-25. No derivative of a newly
% selected polynomial approximation is being substituted for f_j.

% Apply the three-mode Mehler bound with(LC1), the SAME measured centering
% translation and thresholds(.04399,.09999,.02499). The largest tail is
% less than (1.5e10)(4/3)^(3/2) exp[-(3/10)(13.85)^2]<3e-15.
% The other two tails and the actual-input error leave
% ||1_(r outside[.94,1.06])partial_j u||<1e-14.
% The fixed exact inverse-Cauchy map carries this to the SAME static
% interval |x|<=.39 by finite propagation, with no extra multiplier.

% The MP76 Jz<1.5 estimate also extends directly: its64th power on the
% degree-one input reaches degree129, and the last application starts
% at degree127, where its norm is still<.069. Thus the actual tangent
% angular tail is at most K(.069/.2)^64 plus the cut comparison,
% which is<1e-21. P0 commutes with TOTAL Jz. Applying the same static
% barrier/heat estimate to the retained tangent, whose norm is<2e8,
% proves the uniform actual physical bound
%                  ||P_l partial_j u||<2e-14.            (LC2)
% The original low-energy component has norm<2e-24. These norms persist
% under the common fixed H; no derivative of the spectral boundary occurs.

% Keep the exact original initial centroid and its launch derivative,
% and use the physical-velocity integrals for the retained approximation
% as in MP76/GR. For each baseline launch column their vector error is
% <=2(2.38)(12)[2e-14+(2e8)(2e-24)]<1.2e-12.             (LC3)
% The time column has error<1e-23. Initial mean-frequency rows are kept
% exactly and acquire no omission error in this prescription. The fixed
% conversion to MP8's stated dimensionless labels has norm<4.
% If the ACTUAL centroid radius is>=1/2 throughout, the already bounded
% base centroid error<1.2e-22 and ||partial_j<q>||<1.2e10 give the
% radius-column bound
% 4[1.2e-12+2.01(1.2e-22)(1.2e10)]<2e-11.               (LC4)
% The radius premise is still to be certified in the retained propagation.
% This completes the low-energy INPUT omission's baseline launch control,
% not its high-energy complement, a restriction on internal source
% energies, a retained dimension or the inverse return Jacobian. No
% mixed spacing/launch derivative or final coefficient is claimed.
% SOURCES: complete actual MP73--76/SOURCES; independent elementary
% polynomial/spectral and heat/Schur checks; direct degree-one Gaussian
% extension of MP61 and the actual already differentiated preparation.
% LC1--4 were derived and saved before any counterpart launch extension
% was read. No new execution or state definition is used.

% ACTUAL MP77--78 COMPARISON [R, WIP]. Read their complete actual text and
% SOURCES. The signed CP phase, positive connection, inherited negative-
% kappa convention and half-density Jacobian agree. MP78 independently
% checks every LC degree, normalization, angular threshold and velocity
% column; no repair remains. These are representations/input errors,
% not a propagated annular bound or an actual return-root enclosure.

% A FINITE SPATIAL DOMAIN FOR THE FULL NATIVE RESPONSE [R, independent WIP].
% This is an alternative literal spatial realization of the ORIGINAL
% continuum matrix problem. It does not restrict the receiving energies
% or infer finite propagation for the nonlocal free-band compression.
% The same selected input, spacing derivative and exact initial readouts
% are approximated with the following explicitly bounded compact pair.

% Put t=2r-3 and P(t)=35t^4-84t^5+70t^6-20t^7. Define chi(r)=1 for
% r<=3/2, chi(r)=1-P(2r-3) for3/2<r<2, and chi(r)=0 for r>=2.
% This C^3 radial function obeys0<=chi<=1. Its first three Cartesian
% component derivatives are bounded by(5,110,2000), respectively:
% P'=140t^3(1-t)^3, P''=420t^2(1-t)^2(1-2t), and
% P'''=840t(1-t)(1-5t+5t^2), with r>=3/2 in the transition.
% It is only a numerical representation cutoff, not a change of the
% preparation or the smooth core. Set u_c(0)=chi u(0),
% w_c(0)=chi w(0), and u_c,j(0)=chi u_j(0), without renormalizing.
% In particular w_c(0) is not replaced by P2 u_c+i m.p u_c: that
% replacement would introduce a different, nonlocal initial map.

% Here are explicit actual-input graph margins. The already established
% degree79 A39 Gaussian approximant has L2 error<3e-36 uniformly in
% the launch box, and its cut momentum support is p<4.25. Multiplication
% by p^s for s<=3 therefore gives a cut error<4.25^s(3e-36).
% Removing the cut from this polynomial adds a smaller Gaussian-tail
% term; using the same degree<=82 moments gives, conservatively,
%             T_s(u-u_pol)<1e-31, 0<=s<=3.                 (FD1)
% The actual normalizer and centering translation are retained. For
% the polynomial itself, conjugating one momentum by rho^(-Nosc/2),
% rho=1/4, costs at most
% |k_c|+2sqrt(epsilon)||S_k||sqrt(2(D+3))<5
% on degrees D<=82, since |k_c|<=3.01 and ||S_k||<128.
% Thus the weighted norm of every ordered p-tensor through degree3
% is<1000, including its Euclidean tensor sum and the actual normalizer.
% If r>=3/2, |q-q_c|>.48999, so some |q_i-q_ci|>.28.
% The initial position rows(13,19,1) make every corresponding unscaled
% Mehler threshold>60. Hence each polynomial graph tail is bounded by
% 1000(4/3)^(3/2)exp[-(3/10)60^2], which is<1e-460.
% Combining(FD1), the three cutoff derivatives, ordered-tensor product
% rules and the saved H,H^2,H^3 bounds proves the loose uniform margins
% ||(1-chi)u||<1e-35, E3((1-chi)u)<1e-27.                 (FD2)
% The degree83 P2P0 polynomial and the degree-one tangent construction
% give, with their ORIGINAL coefficients m and normalization,
% ||(1-chi)w(0)||<1e-30,
% ||(1-chi)u_j(0)||<1e-24.                               (FD3)
% For the first bound, the existing L2 error<1e-33 is sufficient. Its
% rho-weighted polynomial norm can be bounded by1e6: the three inverse-
% power series have coefficient bound80 through order39, the previous
% geometric ratio<.551, and at most five momentum multiplications,
% each costing<5. The centering term is lower degree. For the tangents,
% LC1's1.5e10 weighted norm and L2 error<1e-25 suffice. Their .28
% threshold tails are again far below the quoted margins. No initial
% spatial tail is being promoted to a propagated graph estimate.

% The exact local continuity equation is
% partial_t |u|^2+div[u^dag(alpha+v)u]=0.
% Its normal flux has norm<=c_*=1+sqrt(5/6)<1.913. An exterior-cone
% energy estimate therefore gives finite propagation speed c_* for H.
% B is a LOCAL third-order differential operator, so the source Bu_c
% has the same support as u_c. Duhamel gives the same cone for w_c.
% Consequently all compact base, response and launch solutions obey
% support subset {r<=2+c_*t} subset {r<25}, 0<=t<=12.       (FD4)
% This is a cone for the continuum generator and its EXACT spacing
% derivative system; finite-a lattice hopping is not assigned compact
% support inside this cone.
% The regular center is included; no artificial inner boundary is used.
% One may solve the full radial Dirac channels on0<r<26. At the outer
% endpoint take any fixed maximal zero-flux boundary line for the matrix
% sigma_y+v(26)I. For example, in its sigma_y eigenbasis impose
% c_- = i sqrt[(1+v(26))/(1-v(26))] c_+.
% The Green boundary form vanishes on this maximal isotropic line, giving
% the self-adjoint regular-endpoint extension of each radial channel.
% All compact solutions vanish in an open neighborhood of that boundary
% throughout the prescribed time interval, so its choice has NO effect
% on(FD4) or on their full Cartesian source/detector matrix elements.
% Negative-kappa component conventions and every angular coherence stay.

% Compare on the original infinite space before any radial discretization.
% E3 of the base error is conserved by the SAME H, hence
% ||w(t)-w_c(t)|| <=1e-30+(12/epsilon)(1180)(1e-27)<2.4e-16.
% Keep the exact original initial centroid, spacing derivative and mean-
% frequency rows, and evaluate their physical-velocity integrals with
% (u_c,w_c). With ||w(t)||<4e14 and c_v<2.38, the raw vector response
% error is less than
% 2c_v(12)[(1e-35)(4e14)+2.4e-16]
%       +12(307)(1e-27)(2*50)<2e-14.                      (FD5)
% The last deliberately loose term bounds the direct V_B quadratic form
% by the saved actual base E2<50. The base-centroid error is<6e-34,
% and each baseline-launch column error is<6e-23 using ||u_j||<2e8.
% These are proxies for the ORIGINAL observable, not centroids or energies
% of a newly normalized compact packet. A radius/Jacobian use still needs
% the actual radius margin and return-root certificate.

% FD supplies a literal finite spatial interval for full two-component
% radial propagation and the complete native insertion. It supplies NO
% mesh, spectral dimension, angular cutoff or scalar-kernel quadrature.
% In particular sharp energy/free-sign projections are nonlocal: FD4
% must NOT be asserted for a subsequently projected compact state, nor
% may infinite-space ER/GR block kernels simply be replaced by finite-box
% kernels without a new comparison. A numerical realization may instead
% use the original local finite-box matrix operator with(FD5), checking
% its retained propagation and all source/readout errors directly.
% SOURCES: independent actual polynomial/Mehler graph estimates and
% ordered cutoff product rules; local Dirac continuity, cone energy
% estimate, inhomogeneous Duhamel and radial Green boundary form. This
% block was derived/saved before reading the original RED's announced
% MP79--80 broad energy cuts. No scientific execution or new state is used.

% ACTUAL MP79--80 FINITE-ENERGY COMPARISON [R, WIP].
% Read both complete actual sources/SOURCES after FD was saved. The
% endpoint complex disk gives |R|<.628 and the central disk keeps
% |Im h|<64/49<pi/2, proving the stated eta Cauchy bounds. The radial
% Taylor majorant8t/(1-10t), its15-fold argument scaling and product
% with q_i correctly yield M_j<=13000^j(j!)^2 for the actual Cartesian
% flow. The finite product-rule constants retain all div(v) terms.
% The positive majorant characteristic y-.075y^2=2t gives exponent
% (7+5/2)y; its t=1 value is<24. Thus ||H^64u||<96^64 and the
% actual input tail at |H|>384 has norm<3e-39 and E3<1.7e-31.
% Its original native initial map, cubic source and velocity proxy give
% the checked raw spacing error<3e-18 and baseline columns<7e-29.

% MP80's operator elliptic induction uses sqrt3/d<20 and the genuine
% finite-energy range of H, so its384 term is an operator norm, not a
% statement about a single moment. For m<=63 the displayed next-step
% ratio is<8061<8192. The ordered source coefficients3.174,10.1,32
% follow from epsilon*13000*64^2 and its powers;32 includes the three
% zero-order Cartesian terms. The finite generating-function step at
% t=n/(2K) is valid through n=61, without an analytic-vector premise.
% With K=8192 and Lambda_o=24K, the exact powers are
% 2K^3 4^-61=2^-82 and K^2 4^-61=2^-96. Consequently the checked
% receiving-energy response omission is<2.5e-15 through time12.
% Both signs and zero receiving energies remain; no small internal
% energy was deleted from the original native source.

% These controls make the infinite-space INPUT and RESPONSE energy
% ranges finite with explicit component errors. They do not select a
% radial/angular quadrature, efficient degree, actual propagated state,
% radius margin or inverse return Jacobian. FD instead supplies a local
% finite-space realization. Sharp spectral projections and spatial
% cutoffs are not interchangeable; any combined realization must use
% the correct operator/domain and explicitly add its comparisons.
% SOURCES: complete actual original RED MP79--80/SOURCES; direct complex
% estimates, finite ordered-product/elliptic recurrences and powers of2
% checked above. Same-pair[R] comparison only; no new run or coefficient.

% A LITERAL RADIAL MATRIX PENCIL AND ITS EQUATION RESIDUAL [R, WIP].
% This fixes a concrete retained representation, not an asserted degree
% or a new packet. Use either signed CP phase in the exterior. Because
% P_s=L p_s(r), e_s=L ebar_s(r), its spin frame U_s, velocity matrix
% A_s and connection b_s=v phi_s'/2 are INDEPENDENT of L. Thus exactly
% U_s^dag exp(-i L Sbar_s/epsilon)(H_kappa-L/sqrt3)
%          exp(i L Sbar_s/epsilon)U_s
%       ={A_s,p_r}/2+epsilon b_s sigma_x-2L ebar_s Pi_-,   (RP1)
% where Sbar_s'=p_s and Pi_-=diag(0,1). Both components remain.
% The shifted radial Galerkin matrix is therefore the AFFINE pencil
% T_N+L V_N; the physical matrix adds (L/sqrt3)I. There is no numerical
% differentiation of a carrier-size
% phase, and the same matrices serve every positive-kappa channel.
% The fixed negative-kappa component map is applied in its own channels.

% One explicit patch can use y=log(r)/sigma, sigma=2^-8, and normalized
% Hermite functions h_n(y) on the real line. Set
% b_n(r)=(sigma r)^(-1/2)h_n(log(r)/sigma).
% The half-density makes {b_n}_{n>=0} an EXACT orthonormal radial basis
% on(0,infinity). It vanishes to every algebraic order at the regular
% center. Its width is a basis parameter only, independent of the
% selected preparation. For a compact exterior patch define
% zeta(r)=eta((r-.87)/.03) eta((1.19-r)/.04),
% with the SAME explicit eta as the core. Then zeta=1 on[.9,1.15],
% zero outside(.87,1.19), and the Schwarzschild CP expressions are
% valid everywhere on its support. Let f_n=zeta b_n,0<=n<=N.
% No approximation to the original core occurs in this basis definition.

% At N<=128 the Gram matrix G_N=(<f_m,f_n>) is uniformly conditioned.
% Indeed outside[.9,1.15] one has |y|>26.9. Mehler with rho=1/2 gives
% ||I-G_N|| <=2^(N+1)erfc(26.9/sqrt3)<1e-65.             (RP2)
% Here erfc(x)<=exp(-x^2)/(sqrt(pi)x); the two radial tails are already
% covered by the two-sided Gaussian integral. Thus G_N^-1/2 defines
% an explicit exact isometry on this patch. This is a finite-span Gram
% bound, not a claim that the propagated packet lies in that span.

% All matrix elements are one-dimensional integrals. Before applying
% G_N^-1/2 on both sides, write f_m,f_n as real functions and use
% T_mn=-(i epsilon/2)int [f_m A_s f_n'-f_m' A_s f_n]dr
%                            +epsilon int f_m b_s sigma_x f_n dr,
% V_mn=-2 int f_m ebar_s Pi_- f_n dr.                    (RP3)
% The derivative convention in the antisymmetric first integral includes
% the A_s' term exactly; no midpoint ordering or spin connection is lost.
% T and V are Hermitian finite matrices. With G normalization the full
% radial dimension is2(N+1), before any angular labels or extra patches.
% In particular N=128 would mean258 radial spinor columns, but NO input
% projection or propagation tolerance is assigned to that number here.

% Let W_L be this exact finite isometry INCLUDING its phase and spin
% frame. Define K_L=W_L^dag H_kappa W_L and the actual differential
% defect R_L=H_kappa W_L-W_L K_L. It is computable on the same compact
% patch, including derivatives of zeta. For c(t)=exp(-itK_L/epsilon)c0,
% the full-space equation residual is exactly -R_L c(t).
% Its squared time integral needs no sampled residual curve. If
% K_L=Z diag(lambda_a)Z^dag and D=Z^dag R_L^dag R_L Z, then
% int_0^T||R_L c(t)||^2dt
%  =sum_ab (Z^dag c0)_a^* D_ab(Z^dag c0)_b
%       T exp[i(lambda_a-lambda_b)T/(2epsilon)]
%          sinc[(lambda_a-lambda_b)T/(2epsilon)].         (RP4)
% The coincident-eigenvalue limit is T. Replacing R_L by H^j R_L
% gives the same computable graph-residual formula whenever that graph
% is used in the previously derived full native source certificate.
% Compact smooth f_n lie in every such domain. Interval integration,
% finite-eigensystem error and roundoff must be enclosed; an approximate
% eigenpair residual is added, not presumed zero.

% The original cubic insertion changes angular labels, so its finite
% matrix is W_(L')^dag B_(kappa',kappa) W_L, with all spinor/Gaunt
% factors. The phase in this integral is
% exp[i(L-L')Sbar_s/epsilon]=exp[i(|kappa|-|kappa'|)Sbar_s],
% and only the bounded tensor-allowed differences occur. For positive
% kappa channels this is exp[i(kappa-kappa')Sbar_s]. Cartesian
% detector blocks keep the corresponding phase too. If different signed
% phase charts or radial patches are used on the two sides, their actual
% two phases remain; the last cancellation is not asserted across them.
% Source defects B W-W B_N and the initial native response are tested
% separately. Launch differentiation changes only the coefficients of
% the original projected/centered input, since this basis is fixed.
% Its exact derivative, initial projection error, angular sums/coherences,
% radial source/coordinate defects and four-return Jacobian are retained.

% RP is a concrete discretization and residual formula for retained
% radial blocks. The exterior patch alone is NOT assumed to contain the
% inward opposite-band packet, or any later source-created component.
% Extra global/core patches or a proved omission would be needed. A union
% of patches uses its actual Gram matrix, not presumed orthogonality.
% The full input projection, necessary degree/patch/angular counts and
% residual acceptance have not been evaluated, so RP authorizes no run
% and claims no finite-frequency value. It can be combined with FD only
% by computing the SAME local operator residual; it does not transfer
% ER/GR's infinite-space kernels to a finite box without comparison.
% SOURCES: direct L-homogeneity of CP, exact logarithmic half-density,
% Hermite Mehler bound, radial Green form, full differential defect and
% exact finite exponential integral. This representation and its error
% diagnostic were derived/saved before counterpart radial discretization
% intake. No numerical dimension was inferred from a linearized width.

% FINITE-ELL MATTER COUPLING IN THE RETAINED REPRESENTATION [R, WIP].
% Let W_ell be the specified normalized radial source kernel and use its
% actual l=1 adjoint on the radial geometric flow: v_ell=W_ell^*v. At
% N=1, the scalar back-average is exactly1. The binding continuum matrix
% and its native spacing derivative are therefore H(v_ell),B(v_ell),
% using the SAME full ordered expressions above. The geometric solution
% at fixed E is still v=-sqrt(r_s/r)n outside the enlarged source support.
% Neither the Weyl endpoint average nor W_ell^* is to replace the other:
% the finite-a symbol uses [v_ell(x+ae_i)+v_ell(x-ae_i)]/2.

% For an even isotropic convolution with second moment sigma_ell^2 delta,
% Taylor's formula in any region separated from the core/source gives
% v_ell=v+(sigma_ell^2/2)Delta v+O(mu_4 sup|D^4v|)
%       =v[1-9sigma_ell^2/(8r^2)]+O(mu_4 r^(-9/2)).
% The last power includes the fixed sqrt(r_s) coefficient. This is a
% finite-ell matter correction, not a geometric spacing correction.
% A different allocation/kernel uses its actual adjoint and moments;
% no exact convolution identity is asserted for an unspecified cell rule.
% The fixed point-coupled MP8 value, if evaluated, is not automatically
% the binding finite-ell value or an ell-independent result.

% RP remains a literal basis for the binding operator without changing
% the prepared packet or its detector. Add the exact first-order matrix
% delta F={v_ell-v,p}/2 to RP's physical K_L and differential residual,
% and delta B=-(1/3)sum_i{(v_ell-v)_i,p_i^3} to every full native source
% block/defect. Their coordinate commutators, physical frequency row and
% launch/root derivatives are included too. This preserves both bands.
% The CP frame is then a representation, not an exact new critical orbit.
% If W_ell is nonnegative/normalized, convolution does not increase the
% global component derivative bounds or sup|v|, so the saved global graph
% estimates and finite propagation speed hold with v_ell. Otherwise its
% actual L1 norm multiplies those bounds and must be retained. This states
% how the finite-ell operator enters the actual residual; no evaluated
% point-control tolerance or full propagation/root is inferred.
% SOURCES: fully reread current final4/GOAL_alt.md; direct adjoint
% substitution, ordered native expansion and isotropic Taylor formula.
% Existing fixed-E geometry and independent Laplacian calculation above.

% ACTUAL MP82 SMOOTH COMPACT-DOMAIN COMPARISON [R, WIP].
% Read the complete actual MP82/SOURCES after RP was saved. The smooth
% representation cutoff is applied to the ORIGINAL base, native response
% and launch vectors separately. Its higher derivative composition gives
% 1600+16/3<1610; including the three Cartesian choices gives
% 3 epsilon*1610*64<.0185. The finite product-rule sum is bounded by
% sum_(j=0)^m[m(.0185)/7]^j, so its factor is<1.21 through m=64.
% In particular1.21*.753<1 proves the displayed response moment without
% changing the nonlocal preparation map. The altered first-three cutoff
% derivatives retain FD's graph error: their largest third tensor cost
% is sqrt(27)*34000*epsilon^3<4e-17.
% Every compact smooth input and all its differential powers vanish near
% r=26; the regular-center angular domains are inherited from the smooth
% Cartesian spinor. Thus the H_b/H moment equality is a common-support
% domain identity, not a comparison of sharp spectral projectors. The
% tail constants follow from4^-64=2^-128: multiplying by1.21 or
% 1.21(2e8), respectively, gives the stated3.6e-39 and7.2e-31 bounds;
% the cubic energy graph remains<2.1e-31. No repair is needed.
% This closes the high-moment INPUT compatibility for FD's finite box.
% It does not establish a bound for B on boundary-touching spectral
% vectors, a finite receiving-space approximation, propagated residual,
% angular count or return root. RP's compact differential residual can
% be tested in this same finite domain; no infinite-space compressed
% kernel is thereby imported. No run or coefficient is inferred.
% SOURCES: complete actual original RED MP82/SOURCES, direct finite
% derivative/product-rule arithmetic, compact-domain powers and spectral
% tail inequalities above. Same-pair[R] only; original state unchanged.

% ACTUAL MP83 FULL DRIVEN-RESIDUAL COMPARISON [R, WIP].
% Read the complete actual MP83/SOURCES. The retained response equation
% has BOTH defects R d and S c, S=B W-W B_N; omitting the latter loses
% the original native source even when the finite ODE is solved exactly.
% Direct integration gives the displayed divided difference with denominator
% lambda_b-lambda_a and limit -it exp(-i lambda_a t/epsilon)/epsilon.
% The block Gram generator i(A^dag D-D A)/epsilon has the correct sign.
% The original full-space Duhamel convolution is int_0^T(T-s)r_0(s)ds,
% so its Cauchy--Schwarz factor is T^(3/2)/sqrt3. At epsilon=2^-24
% this independently reproduces the three displayed multipliers and the
% sufficient8.04e-8 response-norm tolerance. Including the base error
% times the full4e14 response norm gives the raw vector proxy budget
% below5e-6, with the direct307 E2 velocity-response term retained.
% These are sufficient tests on actual finite residuals, not achieved
% numbers. They neither select N=128 nor allow omission of inward bands,
% source patches, the native preparation, launch derivatives or the root.
% In particular compact smooth columns permit measuring H^jR and S in
% the original differential domain without a spectral-box B assumption.
% SOURCES: full actual original RED MP83/SOURCES; independent triangular
% variation-of-constants, entire frequency limits and original-velocity
% error arithmetic above. No scientific execution or coefficient inferred.

% ACTUAL FINITE-ELL RAY COMPARISON [R, WIP].
% Read complete actual original RED MP84/SOURCES. For its stated even
% isotropic convolution and Schwarzschild kernel collar, direct expansion
% gives F_ell=1-r_s/r+9r_s sigma_ell^2/(4r^3)+O(ell^4).
% The circular principal-ray equation rF_ell'-2F_ell=0 has unperturbed
% derivative -3r_s/r0^2, r0=3r_s/2. Thus delta r=-15sigma_ell^2/(4r0),
% and r_ph,ell/r0=1-(5/3)sigma_ell^2/r_s^2+O(ell^4/r_s^4).
% Since d(sqrt(F)/r)/dr=0 at r0, its frequency ratio changes by
% delta(omega/L)/(omega/L)=sigma_ell^2/r_s^2+O(ell^4/r_s^4).
% This checks the signs and factors without a new orbit run. It is an
% a=0 principal-ray averaging control, not the finite-epsilon MP8 detector
% or its native spacing coefficient. It gives no permission to replace
% the selected packet by a ray. Actual cell/allocation kernels require
% their own adjoint moments and collar, and the binding finite-ell full
% operator remains the H(v_ell),B(v_ell) assembly above.
% SOURCES: full actual MP84/SOURCES; independent radial force and simple-
% root variation, after the author's saved vector Laplacian calculation.

% ROTATING TWO-CARRIER INPUT MOMENTS FOR RETAINED RESOLUTION [R, WIP].
% The next radial resolution question concerns the ACTUAL packet's
% residual frequency, not just the large carrier label. Keep point coupling
% here, as in the fixed selected member, and define
% K=H-Omega J_z, Omega=1/sqrt3,
% b(q)=v(q)-Omega(e_z cross q), d(q,k)=b(q).k,
% k_mat(q,k)=alpha.k+d(q,k)I. The spin part of J_z must be retained:
% in native Cartesian spinors J_z=(q cross k)_z+(epsilon/2)Sigma_z.
% Consequently the EXACT Weyl symbol of K is
% k_mat-(epsilon Omega/2)Sigma_z, not k_mat alone. Set
% C_s=-(epsilon Omega/2)Sigma_z. Since C_s is constant, the exact
% symbol of the polynomial Q=K(K+6I) is
% q_W=|k|^2+d^2+6d+(epsilon^2/4)tr[(D b)^2]
%       +(2d+6)alpha.k+{alpha.k+d,C_s}+C_s^2+6C_s.     (RM1)
% Matrix products in the last three terms are literal. The derivative
% square is tr[(Db)^2], not the Frobenius square. Formula(RM1) follows
% because d is linear in k: its self-star product stops after order2,
% and d star d=d^2+(epsilon^2/4)sum_ij(partial_i b_j)(partial_j b_i).
% The first star corrections in {alpha.k,d} cancel. No spin or rotational
% connection has been discarded. This gives an explicit second-order
% differential operator for finite actual-input moment calculations.
% Equivalently, with Gamma=alpha_z Sigma_z,
% q_W=Q_principal-epsilon Omega[(d+3)Sigma_z+k_z Gamma]
%                  +(epsilon^2/4)[tr((Db)^2)+Omega^2]I.
% At the stated carrier Db has flow diagonal(c/2,-c,-c), c=sqrt(2/3),
% plus the antisymmetric rotation part. Thus tr((Db)^2)=5/6 and the
% exact local symbol is q_W(q_*,k_*)=(7epsilon^2/24)I. This is a
% local algebra check, not the expectation or spectral width of the packet.

% At q=e1,k=(sqrt6,sqrt3,0), d=-3 and the PRINCIPAL constant matrix
% k_mat has eigenvalues0,-6. Its Q polynomial vanishes on both bands.
% The upper principal first differential also vanishes because this is
% the critical rotating orbit; the lower one does not. Thus a high power
% of K alone would spuriously charge the small lower-band component with
% powers of6. A bound on Q^n retains it without assuming equal upper and
% lower spectral widths. The exact epsilon spin terms in(RM1) stay in
% every power and are not absorbed into a changed carrier or preparation.

% There is an exact block recurrence that exposes the required cancellation.
% The upper free projector P=P0(k) commutes with the FULL J_z, so write
% K=[[A,V^dag],[V,D]] in P plus (1-P), with V=(1-P)F P and
% F={v,k}/2. Then
% Q_++=A^2+6A+V^dag V,
% Q_-+=V A+(D+6)V,
% Q_+-=A V^dag+V^dag(D+6),
% Q_--=D(D+6)+V V^dag.                                  (RM2)
% For the original u=P u, set a0=u,b0=0 and iterate
% a_(n+1)=Q_++ a_n+Q_+- b_n,
% b_(n+1)=Q_-+ a_n+Q_-- b_n.
% Hence ||Q^n u||^2=||a_n||^2+||b_n||^2. The term(D+6)V,
% and not6V by itself, is the lower input source. Its actual size still
% has to be bounded/evaluated; a formal cancellation is not a tail bound.
% This full-matrix recurrence is equivalent to applying(RM1) in the
% native representation and does not require an approximate band deletion.

% Since H and J_z commute for the prescribed radial background, Q is a
% self-adjoint polynomial of K on its natural domain. Any actual bound
% ||Q^n u||<=mu_n gives the exact input estimate
% ||1_{|K(K+6)|>Lambda}u||<=mu_n/Lambda^n.                (RM3)
% For0<Lambda<9 the retained intervals are
% [-3-sqrt(9+Lambda),-3-sqrt(9-Lambda)] and
% [-3+sqrt(9-Lambda),-3+sqrt(9+Lambda)]. No internal receiving energy
% of B is thereby removed.
% Initial native response and launch columns need their separate moments.
% Finite angular labels must retain the spin in J_z. The actual continued
% spectral reference is irrelevant to this prescribed rank-one propagation
% calculation; no transport-reference identification is used.

% RM fixes the literal operator/moment to assess before assigning efficient
% radial resolution. No mu_n, useful Lambda, retained dimension, matrix
% quadrature or runtime is asserted. In particular MP72's reference RMS
% estimate alone does not justify a narrow finite spectral projection.
% The actual projection/centering, cutoff and weighted derivative errors
% must be included in any polynomial/Gaussian evaluation of(RM1). This
% calculation is for the unchanged full MP8 input, not a scalar substitute.
% SOURCES: independent exact Weyl products of a linear-momentum flow,
% full orbital-plus-spin rotation, free-band block multiplication and
% spectral Markov inequality. Derived/saved before any new counterpart
% rotating two-root input-moment calculation was read. Earlier MP74's
% announced unrotated two-root idea is acknowledged, not imported as a
% rotating bound or as a numerical result.

% STATIONARY OPTICAL ACTION FOR THE DIRECTIONAL EXTERIOR TESTS [scoped X].
% These are frozen principal-symbol observables, with point coupling and
% N=1, v=-sqrt(r_s/r)n in a Schwarzschild collar. They are not the full
% matrix MP8 centroid-return observable. Put lambda_a=(a omega)^2,
% H_a=|p|+v.p+a^2 h_2+O(a^4|p|^5),
% h_2=-sum_i p_i^4/(6|p|)-(2/3)sum_i v_i p_i^3.
% For a unit tangent t, the fixed-energy support function is the
% stationary value p.t on H_a(p)=omega, with its group velocity parallel
% to t. Write w=v.t, D=sqrt(1-|v|^2+w^2),
% L=(D-w)/(1-|v|^2)=1/(D+w). The continuum saddle is
% p_0=omega(t-Lv)/D, |p_0|=omega L/D; the Lagrange multiplier is L.
% Consequently its spacing derivative is -L h_2(p_0), and the exact
% order-a^2 optical action density is
% F_a(t)/omega=L+lambda_a D^-3[
%  c(t)/6-L^2 sum_i t_i^2 v_i^2
%       +(4/3)L^3 sum_i t_i v_i^3-(1/2)L^4 sum_i v_i^4]+O(lambda_a^2),
% c(t)=sum_i t_i^4.                                      (OT1)
% Extend homogeneously to nonunit tangents. In particular every term
% linear in v cancels in the spacing derivative. This cancellation uses
% the declared two-step current's coefficient 2/3; a one-step flow would
% give a different optical action. The multiplier, hence conserved omega
% and group-velocity aiming, is held fixed by this Legendre calculation.

% In the weak field the continuum odd part -v.t/(1-v^2) is an exact
% radial one-form and has no Euler--Lagrange force. The even continuum
% optical potential starts with (v^2+w^2)/2. The spacing terms through
% cubic order in v are, after the free c(t)/6 term,
% U_2=c(t)(v^2-w^2)/4-sum_i t_i^2 v_i^2,
% U_3=2w sum_i t_i^2v_i^2+(4/3)sum_i t_i v_i^3.          (OT2)
% Take the unperturbed incoming line q=z d+b e, d.e=0, |d|=|e|=1,
% b>0, and define c=sum d_i^4, t31=sum d_i^3e_i,
% t22=sum d_i^2e_i^2, t13=sum d_i e_i^3. Then
% U_2=r_s[-c z^2-2t31 bz+(c/4-t22)b^2]/(z^2+b^2)^(3/2).
% Its integrated b derivative is +2c r_s/b, opposite to the continuum
% derivative -2r_s/b. The free optical transverse Hessian in direction e
% is 1+lambda_a(2t22-c/2), because the homogeneous free correction is
% sum_i t_i^4/(6|t|^3). Dividing the force by this Hessian therefore gives
% (alpha_a/alpha_0-1)/lambda_a=-c/2-2t22+o(1),           (OT3)
% at matched conserved omega, incoming group direction and impact datum.
% This is the in-plane bending component; an order-lambda_a transverse
% component does not change the magnitude at first order about the
% nonzero continuum in-plane deflection. It is not obtained by identifying
% group direction with canonical momentum direction on the free lattice.

% There is also a definite half-integer weak-field term. The even part of
% U_3 along the incoming line is
% -r_s^(3/2)[8t31 b z^2+(4/3)t13 b^3]/(z^2+b^2)^(9/4).
% Set I=int_R s^2(1+s^2)^(-9/4)ds
%       =(8sqrt(pi)/5)Gamma(3/4)/Gamma(1/4).
% The companion integral without s^2 is (3/2)I. Thus int U_3 dz
% =-2I(4t31+t13)r_s^(3/2)b^(-1/2), and
% (alpha_a/alpha_0-1)/lambda_a
% =-c/2-2t22-(I/2)(4t31+t13)sqrt(r_s/b)
%       +O((r_s/b)log(R_far/b))+O(lambda_a)+O(b/R_far).  (OT4)
% The last remainder denotes a specified finite-boundary impact convention;
% it is not a claim of an R_far-independent coefficient of r_s/b.
% Continuum ray displacement begins at order r_s, after removing the exact
% radial one-form, so evaluating U_3 on the straight line is sufficient
% at order r_s^(3/2). Continuum/spacing trajectory variations first enter
% the next, integer, weak-field order. Boundary tails must be controlled
% before converting (OT4) to an asymptotic scattering convention.

% For normal[001], d=(cos theta,sin theta,0),
% e=(-sin theta,cos theta,0): c=1-(1/2)sin^2(2theta),
% t22=(1/2)sin^2(2theta), t31=-(1/4)sin(4theta),
% t13=(1/4)sin(4theta). Hence C0=-1/2-(3/4)sin^2(2theta)
% and C_(1/2)=(3I/8)sin(4theta). In particular theta=pi/8
% gives C0=-7/8 and C_(1/2)=3I/8, while theta=0,pi/4
% gives zero half-integer coefficient.
% For normal[011], use d=(cos theta,sin theta/sqrt2,-sin theta/sqrt2),
% e=(-sin theta,cos theta/sqrt2,-cos theta/sqrt2). If C=cos theta,S=sin theta,
% c=C^4+S^4/2, t22=3C^2S^2/2,
% 4t31+t13=CS[S^2-(7/2)C^2]. Therefore
% C0=-C^4/2-S^4/4-3C^2S^2,
% C_(1/2)=(I/2)CS[(7/2)C^2-S^2]. At theta=pi/4 these
% equal -15/16 and5I/16. For normal[111], d,e lie in the plane
% q1+q2+q3=0, so c=1/2, t22=1/6 and t31=t13=0:
% C0=-7/12 and C_(1/2)=0 for every in-plane angle. Out-of-plane
% symbol motion is not asserted to vanish by this in-plane invariant.

% The integer weak-field coefficient needs an optical impact convention.
% At a far incoming point q=z d with z<0, the spacing potential gives
% partial_e U_2=-2t31 r_s sign(z)/z^2,
% d/dz(partial_{t_e}U_2)=+2t31 r_s sign(z)/z^2.
% Thus y''=-4lambda_a t31 r_s sign(z)/z^2+... . Starting with
% the same physical group direction at z=-R, the inner transverse
% displacement changes by +4lambda_a t31 r_s log(R2/R1) when R is
% increased from R1 to R2. Since alpha_0 is proportional to1/b,
% delta C=-4t31(r_s/b)log(R2/R1).                        (OT5)
% This explains a boundary-dependent coefficient without treating the
% numerical exit-sphere displacement as an invariant Shapiro delay.
% A renormalized impact prescription can subtract this logarithm, but
% its finite subtraction must be stated before quoting the integer term.

% EXACT SPACING DERIVATIVE OF A RADIAL ROUND-TRIP DELAY [scoped X].
% Use a symmetry direction [100], [110] or [111], with equal nonzero
% Cartesian components so that the group velocity remains radial. Let
% c=sum d_i^4 (respectively1,1/2,1/3), u=sqrt(r_s/r), r_R>r_E>r_s.
% For outward/inward positive momentum magnitudes P_+/P_-,
% omega=(1-mp u)P_pm-(a^2 c/6)(1-mp4u)P_pm^3+O(a^4).
% At fixed omega, the positive speed reciprocals are
% 1/|dot r_pm|=1/(1-mp u)
%       +(lambda_a c/2)(1-mp4u)/(1-mp u)^4+O(lambda_a^2).
% The observable here subtracts the EXACT flat-lattice round-trip time
% at the same omega and direction; it is not the unsubtracted flight time.
% D_0=2r_s log[(r_R-r_s)/(r_E-r_s)]. With s=r_s/r,
% D_a-D_0=lambda_a c int_(r_E)^(r_R)
%       [-6s-21s^2+4s^3-s^4]/(1-s)^4 dr+O(lambda_a^2).  (OT6)
% The free reciprocal sum is2+lambda_a c+O(lambda_a^2), which is
% essential for this formula. Set y=r/r_s-1 and
% g(y)=-6log y+39/y+28/y^2+8/y^3. Then the integral is
% r_s[g(y_R)-g(y_E)]. Equivalently
% (D_a/D_0-1)/lambda_a=-3c F,
% F=1+[39(1/y_E-1/y_R)+28(1/y_E^2-1/y_R^2)
%                         +8(1/y_E^3-1/y_R^3)]/[6log(y_R/y_E)].
% This is the exact a^2 coefficient at fixed exterior endpoints. Its
% weak-field expansion is
% F=1+(13/2)r_s(1/r_E-1/r_R)/log(r_R/r_E)+O((r_s/r_E)^2). (OT7)
% The endpoint expansion must stay uniformly away from u=1 for the
% displayed perturbative remainder. A numerical row inside a chosen
% stellar collar tests only an explicitly continued prescribed PG field.

% BINDING/PROVENANCE SCOPE FOR OT. The finite-ell binding coupling replaces
% v by the actual l=1 adjoint back-average before forming this action.
% For an even isotropic convolution outside a complete Schwarzschild
% kernel collar, delta v/v=-9sigma_ell^2/(8r^2)+O(ell^4/r^4).
% Thus the point-control bending/time formulas acquire relative
% O(sigma_ell^2/b^2) or O(sigma_ell^2/r_E^2) terms, respectively;
% these are separate from lambda_a and no exact ell independence of
% matter observables is inferred. Arbitrary allocation kernels require
% their own adjoint moments. Endpoint ordering, band transport, finite
% packet width and the native MP8 return conditions are also separate.
% SOURCES: independent fixed-energy Legendre stationary action, elementary
% beta integrals, free optical Hessian and radial inverse group speed.
% The complete task7_v3 handoff at
% <internal>/TO_ICLOUD_task7_v3.md
% was read before this derivation and disclosed the coefficient targets;
% no underlying task7 derivation or producer was read before this save.
% Its tables alone are not an independent proof or our numerical acceptance.
% The completed original RED and actual other-lane argument comparisons
% are recorded below; their X scope is the stated point-symbol observable.

% MASSIVE PERIASTRON FROM RADIAL ACTION AND KEPLER AVERAGING [scoped X].
% Point-symbol control, with mu=GE=r_s/2, lambda_m=(am)^2, fixed
% conserved energy m Ecal and angular momentum m ell. In the111 plane
% the restricted quartic invariant is c|p|^4 with c=1/2, and its
% polarization gives sum_i v_i p_i^3=c|p|^2(v.p). An order-a^2
% out-of-plane displacement does not feed the in-plane scalar angle at
% this order about the rotationally invariant continuum orbit.
% Write u=sqrt(2mu/r), f=1-u^2, j=ell/r and
% Rcal=sqrt(Ecal^2-f(1+j^2)). The continuum radial momenta and local
% normal energies in mass units are
% z_pm=(Ecal u pm Rcal)/f, e_pm=(Ecal pm u Rcal)/f.
% At fixed Ecal,ell the first radial-action variation is exactly
% delta I_r/m=(lambda_m c)/(2pi) int_(r_p)^(r_a)
%                         [A(e_+)+A(e_-)]/Rcal dr,
% A(e)=(e^2-1)^2/6-(2/3)(e-Ecal)e(e^2-1).              (PA1)
% The moving turning-point contribution vanishes because the two
% unperturbed radial roots coincide there. The integrable1/Rcal
% endpoint factor is retained; no fixed-r expansion across a root is used.

% Parameterize the CONTINUUM turning points by p=A_orb(1-ecc^2),
% q=sqrt(1-ecc^2), x=mu/p, h=1+ecc cos chi, 0<=chi<=pi.
% Ecal^2=[(1-2x)^2-4x^2 ecc^2]/[1-x(3+ecc^2)],
% ell^2=mu p/[1-x(3+ecc^2)], r=p/h, and
% dr/Rcal=p/(h^2 sqrt x)
%          sqrt[1-x(3+ecc^2)]/sqrt[1-2x(h+2)] dchi.
% The explicit action integrand through the required order is
% A(e_+)+A(e_-)=x^2 S2+x^3 S3+O(x^4),
% S2=[8h^3-32h^2+8h q^2+q^4]/3,
% S3=[112h^4-352h^3-20h^3q^2+192h^2q^2
%                                    -20h q^4-2q^6]/3.
% The additional measure coefficient is h+q^2/2. Elementary integrals
% <h>=1, <h^2>=(3-q^2)/2, <h^-1>=1/q, <h^-2>=1/q^3
% therefore give
% delta I_r/m=lambda_m c mu^(3/2)/sqrt(p)
%   [-4+3q/2+x(-34+18q^2-11q^3/4)+O(x^2)].            (PA2)
% Here brackets mean the uniform chi average over[0,pi]. The coefficient
% -34+18q^2-11q^3/4 includes the perturbed radial measure, not just
% the lattice term evaluated on a Newtonian ellipse.

% Put y=mu/A_orb. Since Ecal^2=1-y+y^2+O(y^3), holding Ecal holds
% A_orb through the order needed here. Also
% ell=sqrt(mu A_orb)q[1+y(2/q^2-1/2)+O(y^2)]. Thus
% delta(varpi)=-2pi partial_ell(delta I_r/m)|Ecal
% =-8pi lambda_m c y/q^2
%       [1+y(55/(2q^2)-4)+O(y^2)].
% The geodesic advance is
% varpi_0=6pi y/q^2[1+y(19/(4q^2)-1/4)+O(y^2)].
% Consequently
% (varpi_a/varpi_0-1)/lambda_m
% =-(4c/3)[1+(mu/A_orb)(91/[4(1-ecc^2)]-15/4)
%                  +O((mu/A_orb)^2)]+O(lambda_m).      (PA3)
% In particular the leading111 coefficient is-2/3. This is a per-radial-
% period in-plane angle on the bound weak-field branch, uniformly away
% from ecc=1 and the separatrix; the circular value is its smooth limit.
% If the two symbols instead start at the SAME geodesic phase-space
% point, delta ell=0 and delta Ecal=O(lambda_m y^2). At fixed ell,
% partial_Ecal varpi_0=O(y), as follows from
% varpi_0=6pi k+(105pi/2)k^2+(15pi/2)k(Ecal^2-1)+...,
% k=mu^2/ell^2. The launch-energy correction is therefore O(lambda_m y^3)
% in the advance, and changes only the O(y^2) remainder of(PA3).
% It cannot be omitted from an exact finite-radius comparison or higher PN.

% The leading non-isotropic secular result can be derived without a
% direction-average substitution. Use the011 plane basis
% e1=(1,0,0), e2=(0,1,-1)/sqrt2 and positive circulation. The plane's
% restricted quartic form on a unit vector of angle theta is
% c(theta)=c0+c2 cos(2theta)+c4 cos(4theta),
% (c0,c2,c4)=(9/16,1/4,3/16).
% For a Kepler ellipse with periastron angle theta0, true anomaly f,
% h=1+ecc cos f, radial velocity W_r=ecc sin f and tangential W_t=h
% in units sqrt(mu/p), the canonical PG momentum is
% p_vec/m=sqrt(mu/p)[W+sqrt(2h)n]. Expanding the ACTUAL two-step
% symbol, including its flow term, gives
% delta H/m=lambda_m[-sum w_i^4/6+u^2 sum n_i^2w_i^2
%                         +(4/3)u^3 sum n_i^3w_i+(u^4/2)sum n_i^4].
% Let I(q,theta0) be its dimensionless Kepler average defined by
% <delta H>/m=lambda_m mu^2/(A_orb^2 q) I(q,theta0).
% The odd flow contribution uses
% sum n_i^3W_i=c(theta)W_r+c'(theta)W_t/4,
% and integration by parts with (sqrt h)'=-ecc sin f/(2sqrt h) gives
% I_odd=6sqrt2 <c'(theta0+f)sqrt h>_f.                  (PA4)
% It is not generally a radial exact derivative. It reverses sign with
% the orbital sense, unlike the even quartic part.

% Define z=(1-q)/(1+q), and A_n(ecc)=<cos(nf)sqrt(1+ecc cos f)>_f.
% Elementary rational Fourier integration of the even part gives
% I_even=c0(4-3q/2)-(3q/2)[c2 z cos(2theta0)+c4 z^2 cos(4theta0)],
% I_odd=-12sqrt2 c2 A_2 sin(2theta0)-24sqrt2 c4 A_4 sin(4theta0).
% A check without any unevaluated orbital dynamics is
% A_n=sqrt((1+q)/2) binom(1/2,n) t^n
%       * hypergeom2F1(n-1/2,-1/2;n+1;t^2), t=ecc/(1+q).
% This converges for every fixed0<=ecc<1 and fixes all signs. Put
% Q_n=A_n-q partial_q A_n=A_n+(q^2/ecc)partial_ecc A_n,
% with the continuous circular limit. Delaunay differentiation gives
% C011=(q partial_q I-I)/3
% =-3/4+q^2/[4(1+q)^2]cos(2theta0)
%    +3q^2(1-q)/[8(1+q)^3]cos(4theta0)
%    +sqrt2 Q_2 sin(2theta0)+(3sqrt2/2)Q_4 sin(4theta0). (PA5)
% This supplies the full leading angle-dependent coefficient, with
% closed special-function definitions of the handoff's Q2,Q4. It is the
% first-order secular per-orbit result at fixed eccentricity/orientation
% in the weak-field limit. Reversing circulation reverses the two sine
% terms when advance is measured in the direction of that circulation.
% An estimator averaged over several finite-radius radial periods must
% retain the ordinary GR rotation of theta0, the induced action response
% and its exact initialization. A seven-period number at finite mu/A_orb
% is not automatically this fixed-orientation coefficient or a unique PN
% correction. No numerical agreement beyond the actual tested scope is
% inferred, and no v10 result has been used.
% SOURCES: author's direct radial-root/action expansion(PA1--3), followed
% by canonical PG-to-Kepler shift, quartic polarization, elementary Fourier
% integrals and Delaunay differentiation(PA4--5). The target formula in
% the actual task7_v3 handoff was known before this derivation. Original
% RED MP86's independently saved leading111 check was read before PA
% was saved; its first PN and full011 were expressly unevaluated there.
% No underlying other-lane periastron derivation or producer was read
% before this save. The completed matching/source and same-pair checks
% are recorded in the actual duty3 comparison below.

% ACTUAL MP85--87 ORIGINAL-PAIR COMPARISON [R, WIP].
% Read the complete actual original RED MP85/86/SOURCES after OT was
% saved, and MP87 after its announcement. The local-flow Sylvester proof
% uses the RIGHT input gap1.99 only; its left Poisson semigroup remains
% gapless. Direct radial Fourier integration gives
% int_(|x|>d)P_t=(2/pi)[atan(t/d)+td/(t^2+d^2)]<=4t/(pi d),
% so no locality of the free projector is being assumed. I checked the
% exterior Hessian op-to-op norm on its stated box: at radial n the
% nontrivial eigenvalue is c r^(-5/2)[3a/8+(3/2)sqrt(1-7a^2/16)],
% |a|<=1, whose maximum remains below2 on that box. The local weighted
% term, global106 bound, Poisson-tail split at12 and anticommutator
% derivative give .312+.0294+.369+.000635<.72, all times epsilon.
% The two physical-resolvent estimates then give5.75e-8 and1.16e-7,
% with the actual low-angular input error added. No mathematical repair
% is identified. These are the actual input amplitude/first graph only;
% neither the old cubic graph nor the full inward response is replaced.
% MP86's independent constrained-action notation g=D+w=1/L maps
% exactly to(OT1). Its beta integral equals our Gamma expression. The
% full impact gradient, which also differentiates t22 with respect to
% the impact direction, gives
% grad_b int U2=(r_s/b)[(2c+4t22)e-4P_perp diag(d_i^2)e].
% Applying the free optical Hessian cancels the leading off-plane piece.
% Thus the leading vector correction is indeed along e; our scalar
% coefficient and the subsequent half-integer/logarithmic controls agree.
% The exact radial-delay integrand agrees before weak-field expansion.
% MP86's independent leading111 Kepler average also agrees with the
% leading term of(PA2); its first PN/full011 were not imported.
% SOURCES: full actual MP85--87/SOURCES and the explicit independent
% Hessian, kernel integral, resolvent arithmetic and optical vector check
% above. Same-pair[R], no new run, no quantum/packet or numerical promotion.

% ACTUAL OTHER-LANE DUTY3 ARGUMENT AND PRODUCER COMPARISON [scoped X].
% This is the STATUS-mapped critical exterior/PPN/dispersion sentence,
% requirements_new F1/B2 and paper 04_results.tex:312--355. The critical-
% first addendum is read in full. OT1--7 and PA1--5 were substantively
% saved before reading the underlying arguments or producers below;
% the target formulas in TO_ICLOUD_task7_v3.md had been disclosed first.
% The entire actual sources read are, beneath
% <internal>/:
% reviewb/review_output/redteam_T7v67/
%   reviewb_azure_redteam_T7v67_20261015_041859/result.md;
% reviewb/review_output/redteam_T7e/
%   reviewb_azure_redteam_T7e_20261015_075728/result.md;
% reviewb/review_output/
%   reviewb_azure_redteam_T7g_20261015_090039/result.md;
% code/engine/engine_ppn_v8.py and code/engine/engine_ppn_v9.py.
% The full task7_v3 handoff and SPEC_final4b_engine.md section17.2 were
% also read. No v10, executed-producer identity or complete raw numerical
% result is claimed read or accepted by this comparison.
%
% The constrained optical Legendre action and its full free transverse
% Hessian in OT1--5 agree with the actual T7v67 sections3--4 and T7e
% sections1--2 arguments. In particular the coefficient is -c/2-2t22,
% not the force correction alone. The beta integral, half-integer term
% and finite-aiming logarithm agree at matched physical group direction.
% OT6--7's direct inverse-speed integration agrees with T7e's fixed-
% energy radial-action derivative, including same-frequency flat
% subtraction. Its closed rational/logarithmic primitive gives the full
% first-spacing radial correction; the 13/2 expression is its weak-
% field expansion, not the exact finite-radius value.
%
% PA1--3's turning-point substitution independently evaluates the111
% radial action. The same original RED's MP88/90 saved action-contour/
% elementary Kepler-integral route gives precisely the action and first
% PN ratio, before its PA intake. This agrees with the actual T7v67
% radial-action statement and T7e section3 at the shared launch. The
% full T7g canonical-shift/orbit-average argument agrees with PA4--5:
% its F=I/q, S_n=A_n and b^2=(1-q)/(1+q), so its q^2 F_q/3 is our
% (q I_q-I)/3, including both sine terms. Original RED MP91 separately
% verifies the quartic polarization, integration by parts and Fourier
% coefficients, and MP92 records the actual complete T7g intake.
% I have read complete actual MP88--92/SOURCES and checked these maps;
% no mandatory mathematical repair remains. This is named analytic
% cross-lane agreement, not a second independent numerical execution.
%
% The actual v8/v9 RHS has the required sine energy, two-step Q and
% cos(2ap) flow velocity. launch matches conserved energy and group
% direction by iteration. orbit_start instead supplies the SAME
% geodesic-apoapsis phase-space point to both systems: delta l=0 and
% delta(E/m)=O((am)^2(mu/A)^2). PA3/MP90 show this energy difference
% first affects the relative periastron ratio at order(mu/A)^2, beyond
% the stated first PN term. The measured precession is a finite-time
% projected periapsis-azimuth average, with ordinary GR apsidal rotation;
% it is not automatically PA5 at a fixed orientation throughout.
% The producers' deflection observable uses a finite boundary and adds
% a continuum tail. In v9 the column called 'min radius / collar'
% actually prints bo*rg/r_collar, the INPUT IMPACT divided by the collar,
% not a computed trajectory minimum. No corrected trajectory-minimum
% evidence or numerical error bar is inferred from that column.
% Fields.all uses exact PG only outside its supplied collar; inside it
% uses the input interpolated N,v. An inside-collar delay therefore needs
% its actual profile before being identified with OT6's exact PG control.
%
% The X scope is point-coupled principal-symbol analytic coefficients:
% directional bending and its half-integer/tail conventions,111 first PN,
% full011 leading apsidal dependence, and radial first-spacing delay.
% It excludes generic integer bending coefficients without an impact
% prescription, unsettled finite-window011 tables, unaccepted v10,
% endpoint/band/width effects, finite-ell cancellation, and the absolute
% finite-frequency MP8 centroid-return coefficient. The earlier accepted
% radial-probe and first-amplitude numerical records remain in this same
% source; their removal from the compact below changes no disposition.
% SOURCES: exact complete argument/producer paths above, full actual
% original RED MP88--92/SOURCES, and the independently saved OT/PA
% calculations. No new derivation branch or scientific execution is
% needed for this critical source comparison. Optional MP8 refinements
% are held under the critical-first addendum; its paper mark stays open.

% NONMODEL DERIVATIVE CONTROL, retained only as algebraic provenance. If one
% chose D_i=(T_i-T_i^-1)/(2a) also for geometry, the particular constraint branch
% v=D Phi, Phi=-2C sqrt r+a^2 C r^(-3/2)(91 A4-45)/32 would have
% v2_i=C r^(-5/2)[672n_i^3+(183-1001A4)n_i]/64 and
% v^2=r_s/r[1+a^2(329A4-183)/(32r^2)]+... . Its zero-normal-speed coefficient
% is (329A4-183)/32; its metric marginal-expansion coefficient is
% (2009A4-903)/672. Its frozen-Bloch orbit has the additional mean
% +(425/288)(a/r_s)^2 and harmonic
% R_geo=425/128+(4333/6528)cos4phi-(1057sqrt2/3264)sin4phi.
% These coefficients solve that chosen stencil, not the binding continuum
% Gauss law. They are excluded from every active exterior result above/below.

\paragraph{Exterior [X: metric, axial, rays; R: averaging, orbit].}
With the asymptotic clock fixed, declared vacuum H2 gives $N=1$,
$v=-\sqrt{r_s/r}$, $r_s=2GE$; $G$ is defined by the Cavendish measurement.
At fixed $E$ the metric is independent of $\ell$, with $\beta=\gamma=1$,
trapping radius $r_s$ and null circle $r_*=3r_s/2$ when these lie in vacuum.

Point-symbol controls match conserved $\omega$, incoming group direction
$d$ and impact $b e$, with unit $d\perp e$; set $c=\sum d_i^4$,
$t_{22}=\sum d_i^2e_i^2$. To first order in $a^2$,
\[
\begin{aligned}
 {\alpha_a/\alpha_0-1\over(a\omega)^2}
 &=-{c\over2}-2t_{22}+O(\sqrt{r_s/b}),\\
 {\varpi_a/\varpi_0-1\over(am)^2}
 &=-{2\over3}\left[1+{GE\over A}
 \left({91\over4(1-e^2)}-{15\over4}\right)+\cdots\right],\\
 {\overline r\over r_*}&=1-{5\over12}(a/r_s)^2
                         -{57\over4}(a\omega)^2+\cdots .
\end{aligned}
\]
Here bending is massless and weak-field; periastron uses a shared
phase-space launch on a weakly bound [111] ellipse with semimajor axis $A$
and eccentricity $e$; $\overline r$ is the azimuthal mean of the
positive-circulation, massless fixed-$\omega$ upper Weyl orbit in the $xy$
plane. The source supplies the bending half-integer/logarithmic terms and
full [011] law. Same-frequency flat-subtracted radial [100]/[110]/[111]
trips have leading relative delay $-3c(a\omega)^2$.

Binding matter uses the adjoint back-average. A normalized even isotropic
convolution of covariance $\sigma_\ell^2 I$, supported in the Schwarzschild
collar, gives $v_\ell/v=1-9\sigma_\ell^2/(8r^2)+O(\ell^4/r^4)$.
The absolute finite-frequency native centroid-return coefficient remains
open. The local axial outgoing branch exists for every finite inflow and
$0<a\omega<1$; at unit inflow, $\dot r\sim\tfrac32(2a\omega)^{2/3}$.
Neither the symbol orbit nor this branch defines a universal packet
trapping shift.
